% Lutte contre les maladies nutritionnelles — SVT 1ère A — Cours signé OnBuch+
% Programme officiel MINESEC. Compiler avec : tectonic NA05-lutte-maladies-nutritionnelles.tex
\newcommand{\DOCMATIERE}{SVT}
\newcommand{\DOCNIVEAU}{1ère A}
\newcommand{\DOCMODULE}{Module 2 : L'Éducation à la Santé}
\newcommand{\DOCLECON}{NA05}
\newcommand{\DOCTITRE}{Lutte contre les maladies nutritionnelles}
% OnBuch+ — préambule « cours signés » ENRICHI (compilé avec tectonic / XeLaTeX)
% Système visuel « Pop » d'OnBuch+ : fond crème (jamais blanc), bordures encre
% épaisses, ombres « dures » décalées, titres Archivo Black en capitales.
% Version MATHÉMATIQUES : graphes (pgfplots/tikz), boîtes pédagogiques
% (définition, propriété, méthode, attention, le savais-tu, exercice résolu,
% exercice, TP, bilan), page de garde et signature OnBuch+ ancrée en bas.
%
% Macros attendues dans main.tex AVANT \input{preamble} :
%   \DOCMATIERE  \DOCNIVEAU  \DOCMODULE  \DOCLECON (numéro)  \DOCTITRE
\documentclass[11pt]{article}

\usepackage{fontspec}
\usepackage[french]{babel}
\usepackage[a4paper,margin=2cm,top=3.4cm,bottom=2.8cm,headheight=1.4cm,footskip=1.3cm]{geometry}
\usepackage{amsmath,amssymb,amsfonts}
\usepackage{mathtools} % \xmapsto, \coloneqq... souvent utilisés par les modèles
\usepackage{cancel} % simplifications barrées (bilans de forces)
% \abs{z} (module, valeur absolue) et \norm{u} (norme), utilisés par les modèles
\providecommand{\abs}{}\let\abs\relax\DeclarePairedDelimiter{\abs}{\lvert}{\rvert}
\providecommand{\norm}{}\let\norm\relax\DeclarePairedDelimiter{\norm}{\lVert}{\rVert}
\usepackage[table]{xcolor} % [table] : fournit aussi \rowcolors (alternance de lignes)
\usepackage{pagecolor}
\usepackage{tikz}
\usetikzlibrary{arrows.meta,positioning,calc,shapes.geometric,shapes.misc,shapes.arrows,decorations.pathmorphing,decorations.pathreplacing,decorations.markings,patterns,shadows,mindmap,trees,fit,backgrounds,matrix,intersections,angles,quotes}
\usepackage{pgfplots}
\usepackage[version=4]{mhchem} % équations nucléaires \ce{^{235}_{92}U}
\usepackage[european,siunitx]{circuitikz} % schémas électriques
\pgfplotsset{compat=1.18}
\usepgfplotslibrary{fillbetween,groupplots} % aires entre courbes (calcul intégral)
\usepackage{enumitem}
\usepackage{fancyhdr}
% [most] charge toutes les bibliothèques tcolorbox (listings, minted, raster,
% documentation...) et gonfle inutilement la mémoire TeX sur un document long
% (~300 boîtes) : on ne charge que ce qui est réellement utilisé.
\usepackage[skins,breakable]{tcolorbox}
\usepackage{graphicx}
\usepackage{colortbl}
\usepackage{booktabs}
\usepackage{tabularx}
\usepackage{diagbox} % case d'en-tête diagonale des tableaux à double entrée
% Colonnes extensibles centrée / gauche / droite (C, L, R), souvent utilisées par les modèles
\newcolumntype{C}{>{\centering\arraybackslash}X}
\newcolumntype{L}{>{\raggedright\arraybackslash}X}
\newcolumntype{R}{>{\raggedleft\arraybackslash}X}
\usepackage{array}
\usepackage{multicol}
\usepackage{siunitx}
% parse-numbers=false : accepte aussi \SI{2,5 \times 10^{-3}}{...}
\sisetup{locale=FR,output-decimal-marker={,},group-minimum-digits=4,parse-numbers=false}
\DeclareSIUnit\ohm{\ensuremath{\symup{\Omega}}} % U+2126 absent de la police
\DeclareSIUnit\division{div} % réglages d'oscilloscope (ms/div, V/div)
\DeclareSIUnit\tr{tr} % tours (vitesse de rotation, ex. \SI{1500}{\tr\per\minute})
\DeclareSIUnit\ch{ch} % cheval-vapeur (puissance moteur, unité usuelle hors SI)
\DeclareSIUnit\franccfa{FCFA} % franc CFA (coûts, contexte camerounais)
\DeclareSIUnit\dioptre{\ensuremath{\delta}} % vergence d'une lentille (dioptrie)
\usepackage{qrcode}
% \degree : commande fréquemment utilisée par les modèles pour un angle ou une
% température en dehors de \SI{}{} (non fournie par les paquets chargés).
\providecommand{\degree}{^{\circ}}
% \qed (fin de démonstration), utilisé par les modèles sans amsthm
\providecommand{\qed}{\hfill$\square$}
% \barre : double barre des valeurs interdites dans un tableau de variations (tkz-tab)
\providecommand{\barre}{\ensuremath{\|}}
% environnement proof (amsthm non chargé) utilisé par les modèles
\ifdefined\proof\else\newenvironment{proof}[1][Démonstration]{\par\noindent\textit{#1.}\ }{\qed\par}\fi
\usepackage{lastpage}
\usepackage{hyperref}
\hypersetup{hidelinks}

% ============================================================
% Palette « Pop » OnBuch+ — mêmes hex que la classe P (plus_ui.dart)
% ============================================================
\definecolor{poporange}{HTML}{F59321}
\definecolor{popdark}{HTML}{E07A0C}
\definecolor{popcream}{HTML}{FFF6E8}
\definecolor{popink}{HTML}{1C1714}
\definecolor{poppurple}{HTML}{7A5AE0}
\definecolor{popgreen}{HTML}{2E9E6A}
\definecolor{poppink}{HTML}{E0457B}
\definecolor{popblue}{HTML}{2F7DE1}
\definecolor{popgold}{HTML}{C98A12}
\definecolor{popmuted}{HTML}{8A7F76}
\definecolor{popline}{HTML}{EADFD2}
\definecolor{popgray}{HTML}{8A7F76}
\colorlet{popgrey}{popgray}
% Alias tolérants (le modèle nomme parfois les couleurs différemment)
\colorlet{popwhite}{white}
\colorlet{popblack}{popink}
\colorlet{popred}{poppink}
\colorlet{popyellow}{popgold}
% Teintes claires pour fonds de tableaux / zones
\colorlet{popblueL}{popblue!10!white}
\colorlet{poporangeL}{poporange!14!white}
\colorlet{popgreenL}{popgreen!12!white}
\colorlet{poppurpleL}{poppurple!12!white}
\colorlet{poppinkL}{poppink!12!white}
\colorlet{popgoldL}{popgold!14!white}

\pagecolor{popcream}

% ============================================================
% Typographie
% ============================================================
\defaultfontfeatures{Ligatures=TeX}
\setmainfont{PlusJakartaSans-Medium.ttf}[Path=../../../fonts/,BoldFont=PlusJakartaSans-Bold.ttf]
% Maths en sans-serif (Fira Math) avec chiffres et lettres droites de Plus
% Jakarta, pour que les formules se fondent dans le texte.
\usepackage[math-style=ISO,bold-style=ISO]{unicode-math}
\setmathfont{FiraMath-Regular.otf}[Path=../../../fonts/]
\setmathfont{PlusJakartaSans-Medium.ttf}[Path=../../../fonts/,range={up/num,up/{Latin,latin},\mathup}]
\usepackage{icomma} % virgule décimale sans espace en mode math
% Caractères mathématiques Unicode absents des polices de texte (Plus Jakarta,
% Archivo Black) : rendus par leur équivalent mathématique, en texte comme en
% formule — sinon ils s'affichent en carré vide (ex. « Arithmétique dans ℤ »).
\usepackage{newunicodechar}
\newunicodechar{ℝ}{\ensuremath{\mathbb{R}}}
\newunicodechar{ℤ}{\ensuremath{\mathbb{Z}}}
\newunicodechar{ℂ}{\ensuremath{\mathbb{C}}}
\newunicodechar{ℕ}{\ensuremath{\mathbb{N}}}
\newunicodechar{ℚ}{\ensuremath{\mathbb{Q}}}
\newunicodechar{𝔻}{\ensuremath{\mathbb{D}}}
\newunicodechar{α}{\ensuremath{\alpha}}
\newunicodechar{β}{\ensuremath{\beta}}
\newunicodechar{γ}{\ensuremath{\gamma}}
\newunicodechar{δ}{\ensuremath{\delta}}
\newunicodechar{ε}{\ensuremath{\varepsilon}}
\newunicodechar{θ}{\ensuremath{\theta}}
\newunicodechar{λ}{\ensuremath{\lambda}}
\newunicodechar{μ}{\ensuremath{\mu}}
\newunicodechar{σ}{\ensuremath{\sigma}}
\newunicodechar{τ}{\ensuremath{\tau}}
\newunicodechar{φ}{\ensuremath{\varphi}}
\newunicodechar{ψ}{\ensuremath{\psi}}
\newunicodechar{ω}{\ensuremath{\omega}}
\newunicodechar{Σ}{\ensuremath{\Sigma}}
% majuscules produites par \MakeUppercase dans les titres (α-aminés -> Α)
\newunicodechar{Α}{\ensuremath{\alpha}}
\newunicodechar{Β}{\ensuremath{\beta}}
\newunicodechar{Γ}{\ensuremath{\Gamma}}
\newunicodechar{Δ}{\ensuremath{\Delta}}
\newunicodechar{Ε}{\ensuremath{\varepsilon}}
\newunicodechar{Θ}{\ensuremath{\Theta}}
\newunicodechar{Λ}{\ensuremath{\Lambda}}
\newunicodechar{Μ}{\ensuremath{\mu}}
\newunicodechar{Τ}{\ensuremath{\tau}}
\newunicodechar{Φ}{\ensuremath{\Phi}}
\newunicodechar{Ψ}{\ensuremath{\Psi}}
\newunicodechar{Ω}{\ensuremath{\Omega}}
\newunicodechar{↦}{\ensuremath{\mapsto}}
\newunicodechar{∈}{\ensuremath{\in}}
\newunicodechar{∉}{\ensuremath{\notin}}
\newunicodechar{∧}{\ensuremath{\wedge}}
\newunicodechar{⇔}{\ensuremath{\Leftrightarrow}}
\newunicodechar{⟺}{\ensuremath{\Longleftrightarrow}}
\newunicodechar{⇒}{\ensuremath{\Rightarrow}}
\newunicodechar{⟹}{\ensuremath{\Longrightarrow}}
\newunicodechar{∘}{\ensuremath{\circ}}
\newunicodechar{∪}{\ensuremath{\cup}}
\newunicodechar{∩}{\ensuremath{\cap}}
\newunicodechar{≡}{\ensuremath{\equiv}}
\newunicodechar{⊂}{\ensuremath{\subset}}
\newunicodechar{⊥}{\ensuremath{\perp}}
\newunicodechar{∃}{\ensuremath{\exists}}
\newunicodechar{∀}{\ensuremath{\forall}}
\newunicodechar{∛}{\ensuremath{\sqrt[3]{\,}}}
\newunicodechar{∜}{\ensuremath{\sqrt[4]{\,}}}
\newunicodechar{⃗}{\ensuremath{{}^{\to}}}
\newunicodechar{ⁿ}{\ifmmode^{n}\else\textsuperscript{\ensuremath{n}}\fi}
\newunicodechar{ˣ}{\ifmmode^{x}\else\textsuperscript{\ensuremath{x}}\fi}
\newunicodechar{ᵇ}{\ifmmode^{b}\else\textsuperscript{\ensuremath{b}}\fi}
\newunicodechar{ʳ}{\ifmmode^{r}\else\textsuperscript{\ensuremath{r}}\fi}
\newunicodechar{ᵘ}{\ifmmode^{u}\else\textsuperscript{\ensuremath{u}}\fi}
\newunicodechar{ʸ}{\ifmmode^{y}\else\textsuperscript{\ensuremath{y}}\fi}
\newunicodechar{ᵃ}{\ifmmode^{a}\else\textsuperscript{\ensuremath{a}}\fi}
\newunicodechar{ᵉ}{\ifmmode^{e}\else\textsuperscript{\ensuremath{e}}\fi}
\newunicodechar{ᵏ}{\ifmmode^{k}\else\textsuperscript{\ensuremath{k}}\fi}
\newunicodechar{ᵖ}{\ifmmode^{p}\else\textsuperscript{\ensuremath{p}}\fi}
\newunicodechar{ᴺ}{\ifmmode^{N}\else\textsuperscript{\ensuremath{N}}\fi}
\newunicodechar{ᶜ}{\ifmmode^{c}\else\textsuperscript{\ensuremath{c}}\fi}
\newunicodechar{ʰ}{\ifmmode^{h}\else\textsuperscript{\ensuremath{h}}\fi}
\newunicodechar{ᵠ}{\ifmmode^{q}\else\textsuperscript{\ensuremath{q}}\fi}
\newunicodechar{ₙ}{\ifmmode_{n}\else\textsubscript{\ensuremath{n}}\fi}
\newunicodechar{ₐ}{\ifmmode_{a}\else\textsubscript{\ensuremath{a}}\fi}
\newunicodechar{ᵢ}{\ifmmode_{i}\else\textsubscript{\ensuremath{i}}\fi}
\newunicodechar{ᵣ}{\ifmmode_{r}\else\textsubscript{\ensuremath{r}}\fi}
\newunicodechar{ₖ}{\ifmmode_{k}\else\textsubscript{\ensuremath{k}}\fi}
\newunicodechar{ᵦ}{\ifmmode_{\beta}\else\textsubscript{\ensuremath{\beta}}\fi}
\newunicodechar{ₓ}{\ifmmode_{x}\else\textsubscript{\ensuremath{x}}\fi}
\newfontfamily\popdisplay{ArchivoBlack-Regular.ttf}[Path=../../../fonts/]
\newfontfamily\poplogo{SpaceGrotesk-Bold.ttf}[Path=../../../fonts/]
\newfontfamily\popsemi{PlusJakartaSans-SemiBold.ttf}[Path=../../../fonts/]
\newfontfamily\popxbold{PlusJakartaSans-ExtraBold.ttf}[Path=../../../fonts/]
\newfontfamily\popxboldls{PlusJakartaSans-ExtraBold.ttf}[Path=../../../fonts/,LetterSpace=3.0]

\setlist[enumerate]{leftmargin=1.7em,itemsep=5pt,topsep=4pt}
\setlist[itemize]{leftmargin=1.7em,itemsep=4pt,topsep=4pt,label={\color{poporange}\textbullet}}
\setlength{\parindent}{0pt}
\setlength{\parskip}{5pt}
\renewcommand{\arraystretch}{1.25}

% pgfplots : style Pop par défaut
\pgfplotsset{
  every axis/.append style={axis line style={popink,line width=1pt},tick style={popink},
    grid style={popline,line width=0.5pt},label style={font=\small},tick label style={font=\footnotesize}},
  popcurve/.style={poporange,line width=1.8pt,no markers,smooth},
}
% Même style disponible dans un \draw TikZ simple (hors pgfplots)
\tikzset{popcurve/.style={poporange,line width=1.8pt}}
\tikzset{popgrid/.style={popline,very thin}}
\colorlet{popcurve}{poporange} % le nom du style est parfois utilisé comme couleur

% ============================================================
% Pill / squiggle
% ============================================================
\newcommand{\poppill}[3]{%
  \tikz[baseline=-0.55ex]\node[draw=popink,line width=1.2pt,rounded corners=2.6pt,%
    fill=#2,inner xsep=6.5pt,inner ysep=3pt]{%
    \popxboldls\fontsize{7}{7}\selectfont\color{#3}\MakeUppercase{#1}};%
}
\newcommand{\popsquiggle}[1]{%
  \tikz[baseline=-0.4ex]{
    \draw[#1,line width=2.6pt,line cap=round]
      plot [smooth] coordinates {(0,0.09) (0.34,0.01) (0.68,0.21) (1.02,0.01) (1.36,0.10)};
  }%
}
% Mot-clé surligné (marqueur orange sous le texte)
\newcommand{\cle}[1]{{\popxbold\color{popdark}#1}}

% ============================================================
% En-tête / pied
% ============================================================
\pagestyle{fancy}
\fancyhf{}
\renewcommand{\headrulewidth}{0pt}
\renewcommand{\footrulewidth}{0pt}
\lhead{\raisebox{-0.15ex}{\poplogo\fontsize{13}{13}\selectfont\color{popink}OnBuch\color{poporange}+}}
\rhead{\poppill{\DOCMATIERE\ \textbullet\ \DOCNIVEAU\ \textbullet\ Leçon \DOCLECON}{white}{popink}}
\lfoot{\popxbold\fontsize{7.6}{7.6}\selectfont\color{popmuted}\MakeUppercase{Cours signé OnBuch+}\ \textbullet\ {\popsemi\DOCTITRE}}
\rfoot{\tikz[baseline=-0.6ex]\node[rounded corners=8.5pt,fill=popink,minimum height=17pt,minimum width=17pt,inner xsep=5pt,inner sep=0]{\popxbold\fontsize{8}{8}\selectfont\color{popcream}\thepage\,/\,\pageref*{LastPage}};}

% ============================================================
% Titres
% ============================================================
\newcommand{\chaptitre}[2]{%
  \begin{tcolorbox}[enhanced,colback=poporange,colframe=popink,boxrule=2.6pt,arc=11pt,
      drop shadow=popink,left=15pt,right=15pt,top=12pt,bottom=11pt,before skip=3pt,after skip=18pt]
    \poppill{#2}{popink}{popcream}\par\vspace{9pt}
    {\raggedright\hyphenpenalty=10000\exhyphenpenalty=10000\popdisplay\fontsize{22}{24}\selectfont\color{popcream}\MakeUppercase{#1}\par}
  \end{tcolorbox}
}
% Section numérotée : pastille orange avec le numéro + titre Archivo
\newcounter{popsec}
\newcommand{\coursec}[1]{%
  \stepcounter{popsec}\setcounter{popsub}{0}%
  \par\vspace{16pt}\noindent
  \begin{minipage}{\linewidth}
  \tikz[baseline=-0.7ex]\node[draw=popink,line width=1.6pt,fill=poporange,rounded corners=4pt,minimum size=22pt,inner sep=0]{\popdisplay\fontsize{12}{12}\selectfont\color{popcream}\thepopsec};%
  \hspace{8pt}{\popdisplay\fontsize{15}{17}\selectfont\color{popink}\MakeUppercase{#1}}\par
  \vspace{4pt}\noindent{\color{poporange}\rule{36pt}{3.6pt}}
  \end{minipage}\par\nopagebreak\vspace{8pt}
}
\newcounter{popsub}
\newcommand{\courssub}[1]{%
  \stepcounter{popsub}%
  \par\vspace{9pt}\noindent{\popxbold\fontsize{11.5}{13}\selectfont\color{popink}{\color{poporange}\thepopsec.\thepopsub}\hspace{6pt}#1}\par\nopagebreak\vspace{4pt}
}

% ============================================================
% Boîtes pédagogiques — même recette : fond blanc, bordure épaisse colorée,
% ombre dure encre, badge plein. Titre optionnel [#1].
% ============================================================
\tcbset{popbase/.style={breakable,enhanced,colback=white,boxrule=2.3pt,arc=9pt,
  shadow={3pt}{-3pt}{0pt}{popink},left=14pt,right=14pt,top=11pt,bottom=11pt,before skip=11pt,after skip=15pt,
  fontupper=\popsemi\fontsize{10.6}{14.5}\selectfont\color{popink},
  fontlower=\popsemi\fontsize{10.6}{14.5}\selectfont\color{popink}}}
% Coche dessinée (le glyphe ✓ n'existe pas dans les polices du document)
\newcommand{\popcheck}[1]{\tikz[baseline=-0.5ex]\draw[#1,line width=1.6pt,line cap=round,line join=round](-0.13,0.0)--(-0.04,-0.09)--(0.13,0.1);}
\providecommand{\coche}{\popcheck{popgreen}}
\providecommand{\resultat}[1]{\par\smallskip\noindent\fcolorbox{popgreen}{popgreenL}{\parbox{\dimexpr\linewidth-2\fboxsep-2\fboxrule\relax}{\textbf{Résultat.} #1}}\par\smallskip}
\newcommand{\popboxhead}[3]{\poppill{#1}{#2}{white}\ifx\relax#3\relax\else\hspace{6pt}{\popxbold\fontsize{10.6}{12}\selectfont\color{popink}#3}\fi\par\vspace{7pt}}

\newtcolorbox{aretenir}[1][]{popbase,colframe=popgreen,before upper={\popboxhead{À retenir}{popgreen}{#1}}}
\newtcolorbox{exemplebox}[1][]{popbase,colframe=poporange,before upper={\popboxhead{Exemple}{poporange}{#1}}}
\newtcolorbox{definition}[1][]{popbase,colframe=popblue,before upper={\popboxhead{Définition}{popblue}{#1}}}
\newtcolorbox{propriete}[1][]{popbase,colframe=poppurple,before upper={\popboxhead{Propriété}{poppurple}{#1}}}
\newtcolorbox{methode}[1][]{popbase,colframe=popgold,before upper={\popboxhead{Méthode}{popgold}{#1}}}
\newtcolorbox{attention}[1][]{popbase,colframe=poppink,before upper={\popboxhead{Attention}{poppink}{#1}}}
\newtcolorbox{experience}[1][]{popbase,colframe=popblue,colback=popblueL,before upper={\popboxhead{Expérience}{popblue}{#1}}}
% « Le savais-tu ? » : carte violette pleine (contexte, culture, Cameroun)
\newtcolorbox{savaistu}[1][]{popbase,colback=poppurple,colframe=popink,
  fontupper=\popsemi\fontsize{10.4}{14.2}\selectfont\color{white},
  before upper={\poppill{Le savais-tu\,?}{white}{poppurple}\ifx\relax#1\relax\else\hspace{6pt}{\popxbold\color{white}#1}\fi\par\vspace{7pt}}}
% Exercice résolu : énoncé en haut, \tcblower puis la solution sur fond crème
\newcounter{popexo}\newcounter{popexr}
\newtcolorbox{exoresolu}[1][]{popbase,colframe=poporange,colbacklower=popcream,
  segmentation style={popink,line width=1.2pt,dashed},
  before upper={\stepcounter{popexr}\popboxhead{Exercice résolu \thepopexr}{poporange}{#1}},
  before lower={\poppill{Solution}{popink}{popcream}\par\vspace{6pt}}}
% Exercice (non résolu en place) — la correction est dans la partie Corrigés
\newtcolorbox{exercice}[1][]{popbase,colframe=popink,
  before upper={\stepcounter{popexo}\popboxhead{Exercice \thepopexo}{popink}{#1}}}
% Corrigé
\newtcolorbox{corrige}[1][]{popbase,colframe=popgreen,colback=popgreenL,
  before upper={\popboxhead{Corrigé}{popgreen}{#1}}}
% Objectifs / prérequis / bilan
\newtcolorbox{objectifs}{popbase,colframe=popink,colback=poporangeL,
  before upper={\popboxhead{Objectifs de la leçon}{popink}{}}}
\newtcolorbox{prerequis}{popbase,colframe=popink,colback=popblueL,
  before upper={\popboxhead{Prérequis}{popblue}{}}}
\newtcolorbox{bilan}{popbase,colframe=popink,colback=white,boxrule=2.8pt,
  before upper={\popboxhead{Fiche bilan}{poporange}{L'essentiel à connaître pour l'examen}}}
% Figure encadrée avec légende
\newtcolorbox{popfigure}[1][]{enhanced,breakable=false,colback=white,colframe=popline,boxrule=1.4pt,arc=8pt,
  left=10pt,right=10pt,top=10pt,bottom=8pt,before skip=10pt,after skip=12pt,halign=center,
  lower separated=false,halign lower=center,
  fontlower=\popsemi\fontsize{9}{11.5}\selectfont\color{popmuted},
  #1}
\newcommand{\legende}[1]{\par\vspace{6pt}{\popsemi\fontsize{9}{11.5}\selectfont\color{popmuted}#1}}

% ============================================================
% Signature OnBuch+
% ============================================================
\newcommand{\poplogoinline}[1]{{\poplogo\fontsize{#1}{#1}\selectfont\color{popink}OnBuch\color{poporange}+}}
\newcommand{\signatureonbuch}{%
  \begin{tcolorbox}[enhanced,colback=popink,colframe=popink,boxrule=2.2pt,arc=10pt,
      drop shadow=poporange,left=14pt,right=14pt,top=10pt,bottom=10pt,before skip=0pt,after skip=0pt]
    \tikz[baseline=-0.7ex]\node[circle,fill=popgreen,draw=popcream,line width=1.3pt,minimum size=22pt,inner sep=0]{\popcheck{white}};%
    \hspace{9pt}\begin{minipage}[c]{0.62\linewidth}
      {\popxbold\fontsize{12}{13}\selectfont\color{popcream}Signé\ \textbullet\ {\poplogo OnBuch\color{poporange}+}}\par\vspace{2pt}
      {\raggedright\popsemi\fontsize{8}{10.5}\selectfont\color{popline}\MakeUppercase{Équipe pédagogique OnBuch+}\\\MakeUppercase{Conforme au programme officiel MINESEC}\par}
    \end{minipage}\hfill
    {\poplogo\fontsize{22}{22}\selectfont\color{popcream}OnBuch\color{poporange}+}
  \end{tcolorbox}%
}

% ============================================================
% Page de garde
% #1 titre · #2 matière · #3 niveau · #4 module · #5 n° leçon · #6 durée ·
% #7 description / objectifs (texte ou itemize)
% ============================================================
\newcommand{\pagegarde}[7]{%
  \thispagestyle{empty}
  \newgeometry{a4paper,margin=1.7cm,top=1.6cm,bottom=1.4cm}%
  \begin{tikzpicture}[remember picture,overlay]
    \fill[poporange!18] ([xshift=-3.2cm,yshift=-3.6cm]current page.north east) circle (4.6cm);
    \fill[poppurple!14] ([xshift=2.4cm,yshift=7.6cm]current page.south west) circle (3.8cm);
  \end{tikzpicture}%
  {\poplogo\fontsize{30}{30}\selectfont\color{popink}OnBuch\color{poporange}+}\hfill
  \raisebox{0.6ex}{\poppill{Cours signé \textbullet\ Premium}{popink}{popcream}}\par
  \vspace{1pt}\popsquiggle{poppurple}\par
  \vspace{8pt}
  \begin{tcolorbox}[enhanced,colback=poporange,colframe=popink,boxrule=2.8pt,arc=13pt,
      drop shadow=popink,left=18pt,right=18pt,top=12pt,bottom=13pt,after skip=10pt]
    \poppill{#2 \textbullet\ #3}{popink}{popcream}\hspace{5pt}\poppill{#4}{white}{popink}\par\vspace{8pt}
    {\popdisplay\fontsize{14}{14}\selectfont\color{popink}LEÇON #5}\par\vspace{4pt}
    {\raggedright\hyphenpenalty=10000\exhyphenpenalty=10000\popdisplay\fontsize{25}{27}\selectfont\color{popcream}\MakeUppercase{#1}\par}
    \vspace{8pt}
    {\popxbold\fontsize{9.5}{9.5}\selectfont\color{popink}\MakeUppercase{Durée conseillée : #6}}
  \end{tcolorbox}
  \begin{tcolorbox}[enhanced,colback=white,colframe=popink,boxrule=2.2pt,arc=10pt,
      drop shadow=popink,left=16pt,right=16pt,top=10pt,bottom=10pt,after skip=8pt]
    \poppill{Dans cette leçon}{poppurple}{white}\par\vspace{5pt}
    \popsemi\fontsize{9.3}{12.6}\selectfont\color{popink}#7
  \end{tcolorbox}
  \noindent\begin{minipage}[t]{0.487\linewidth}
    \begin{tcolorbox}[enhanced,colback=popgreen,colframe=popink,boxrule=2.2pt,arc=10pt,
        drop shadow=popink,left=11pt,right=11pt,top=10pt,bottom=10pt,height=3.1cm,valign=center]
      \tcbox[colback=white,colframe=popink,boxrule=1.3pt,arc=5pt,boxsep=0pt,nobeforeafter,
        left=3pt,right=3pt,top=3pt,bottom=3pt,valign=center]{\qrcode[height=1.9cm]{https://play.google.com/store/apps/details?id=com.luvvix.onbuch&hl=fr}}%
      \hspace{8pt}\begin{minipage}[c]{0.5\linewidth}
        {\popdisplay\fontsize{11}{12.5}\selectfont\color{white}\MakeUppercase{Télécharge OnBuch}}\par\vspace{4pt}
        {\raggedright\popsemi\fontsize{8.4}{11}\selectfont\color{white}Gratuit \textbullet\ Google Play\\Tuteur IA, annales, concours, résultats\par}
      \end{minipage}
    \end{tcolorbox}
  \end{minipage}\hfill
  \begin{minipage}[t]{0.487\linewidth}
    \begin{tcolorbox}[enhanced,colback=poppurple,colframe=popink,boxrule=2.2pt,arc=10pt,
        drop shadow=popink,left=11pt,right=11pt,top=10pt,bottom=10pt,height=3.1cm,valign=center]
      \tikz[baseline=-0.7ex]\node[circle,fill=white,draw=popink,line width=1.3pt,minimum size=1.6cm,inner sep=0]{\popdisplay\fontsize{24}{24}\selectfont\color{poppurple}+};%
      \hspace{8pt}\begin{minipage}[c]{0.6\linewidth}
        {\popdisplay\fontsize{11}{12.5}\selectfont\color{white}\MakeUppercase{Passe à OnBuch+}}\par\vspace{4pt}
        {\raggedright\popsemi\fontsize{8.4}{11}\selectfont\color{white}Tous les cours signés, Tuteur IA illimité, fiches PDF — onglet OnBuch+ de l'app\par}
      \end{minipage}
    \end{tcolorbox}
  \end{minipage}
  \vspace{8pt}
  \popcontenu
  \vfill
  \signatureonbuch
  \restoregeometry
  \newpage
}

% Rangée « Ce cours contient » (page de garde)
\newcommand{\poptile}[3]{% #1 couleur #2 gros texte #3 légende
  \begin{tcolorbox}[enhanced,colback=#1,colframe=popink,boxrule=2pt,arc=9pt,shadow={2.5pt}{-2.5pt}{0pt}{popink},
      left=6pt,right=6pt,top=9pt,bottom=9pt,halign=center,valign=center,height=1.75cm,width=\linewidth,nobeforeafter]
    {\popdisplay\fontsize{12.5}{13}\selectfont\color{white}\MakeUppercase{#2}}\par\vspace{3pt}
    {\centering\popsemi\fontsize{7.8}{9.5}\selectfont\color{white}#3\par}
  \end{tcolorbox}}
\newcommand{\popcontenu}{%
  \par\noindent{\popxboldls\fontsize{7.5}{7.5}\selectfont\color{popmuted}CE COURS SIGNÉ CONTIENT}\par\vspace{7pt}
  \noindent\begin{minipage}[t]{0.235\linewidth}\poptile{popblue}{Cours}{complet, illustré, pas à pas}\end{minipage}\hfill
  \begin{minipage}[t]{0.235\linewidth}\poptile{poporange}{Méthodes}{et exercices résolus}\end{minipage}\hfill
  \begin{minipage}[t]{0.235\linewidth}\poptile{poppink}{Exercices}{type examen + corrigés}\end{minipage}\hfill
  \begin{minipage}[t]{0.235\linewidth}\poptile{popgreen}{Fiche bilan}{l'essentiel à retenir}\end{minipage}\par}

% Sommaire
\newtcolorbox{sommaire}{enhanced,colback=white,colframe=popink,boxrule=2.2pt,arc=10pt,
  drop shadow=popink,left=18pt,right=18pt,top=13pt,bottom=13pt,before skip=6pt,after skip=20pt,
  before upper={\poppill{Sommaire}{poppurple}{white}\par\vspace{9pt}\popsemi\fontsize{10.6}{14.5}\selectfont\color{popink}}}

% Signature de fin de cours — poussée tout en bas de la dernière page
\newcommand{\findecours}{%
  \par\vfill\nopagebreak
  \begin{center}\popsquiggle{poporange}\end{center}\vspace{4pt}
  \signatureonbuch
}
\AtBeginDocument{\def\llbracket{\mathopen{[\mkern-3.5mu[}}\def\rrbracket{\mathclose{]\mkern-3.5mu]}}} % intervalles d'entiers (glyphe ⟦ absent de la police)
% Glyphes absents de Fira Math : redéfinis à partir de glyphes disponibles
\AtBeginDocument{%
  \def\setminus{\mathbin{\backslash}}\let\smallsetminus\setminus
  \def\vdots{\mathord{\vbox{\baselineskip3.2pt\lineskiplimit0pt\kern4pt\hbox{.}\hbox{.}\hbox{.}}}}%
  \def\ddots{\mathinner{\mkern1mu\raise6.5pt\hbox{.}\mkern2mu\raise3.6pt\hbox{.}\mkern2mu\raise0.7pt\hbox{.}\mkern1mu}}%
  \def\triangle{\mathord{\scalebox{0.9}{$\symup{\Delta}$}}}%
  \def\nabla{\mathord{\raisebox{\depth}{\scalebox{1}[-1]{$\symup{\Delta}$}}}}%
  \def\checkmark{\popcheck{popdark}}%
  \def\star{\mathbin{\ast}}\def\bigstar{\mathord{\ast}}%
  \def\diamond{\mathbin{\scalebox{0.6}{$\Diamond$}}}%
  \def\bigcup{\mathop{\mathchoice{\raisebox{-0.3em}{\scalebox{1.7}{$\cup$}}}{\scalebox{1.25}{$\cup$}}{\cup}{\cup}}\displaylimits}%
  \def\bigcap{\mathop{\mathchoice{\raisebox{-0.3em}{\scalebox{1.7}{$\cap$}}}{\scalebox{1.25}{$\cap$}}{\cap}{\cap}}\displaylimits}%
  \def\lesseqqgtr{\mathrel{\lessgtr}}\def\bigoplus{\mathop{\oplus}\displaylimits}%
}
\providecommand{\ssi}{si et seulement si} % abréviation fréquente
\AtBeginDocument{\providecommand{\wideparen}{\overparen}} % arc de cercle

%% FIGFIX-BEGIN v5 — correctifs globaux des figures (docs/onbuch-generation/tools/figfix)
\usepackage{adjustbox}\usepackage{etoolbox}\tcbuselibrary{raster}
% toute figure TikZ/pgfplots plus large que la ligne est réduite à la largeur disponible
\BeforeBeginEnvironment{tikzpicture}{\begin{adjustbox}{max width=\linewidth}}
\AfterEndEnvironment{tikzpicture}{\end{adjustbox}}
% légendes pgfplots lisibles par-dessus les courbes
\pgfplotsset{every axis/.append style={legend style={fill=white,fill opacity=0.92,text opacity=1,draw=black!15,inner sep=3pt}}}
% étiquettes d'arêtes (syntaxe edge["..."]) avec fond blanc
\tikzset{every edge quotes/.append style={fill=white,inner sep=1.2pt}}
%% FIGFIX-END
\begin{document}

\pagegarde{Lutte contre les maladies nutritionnelles}{SVT}{1ère A}{Module 2 : L'Éducation à la Santé}{NA05}{3 h}{Cette leçon présente les principales maladies nutritionnelles : les maladies de carence (malnutrition protéino-énergétique, avitaminoses, anémie) et les maladies de surcharge (obésité). Elle donne les outils pour identifier leurs signes, comprendre leurs causes et conséquences, et proposer des mesures de prévention fondées sur une alimentation équilibrée.
\vspace{4pt}
\begin{multicols}{2}\small
\begin{itemize}
\item À la fin de cette leçon, je sais définir une maladie nutritionnelle et distinguer maladies de carence et maladies de surcharge.
\item À la fin de cette leçon, je sais décrire les signes cliniques du kwashiorkor et du marasme à partir d'un cas concret.
\item À la fin de cette leçon, je sais citer les principales avitaminoses (A, C, D) et l'anémie ferriprive, avec leurs causes et leurs signes.
\item À la fin de cette leçon, je sais expliquer les causes et les conséquences de l'obésité sur la santé.
\item À la fin de cette leçon, je sais calculer et interpréter l'indice de masse corporelle (IMC) d'un individu.
\item À la fin de cette leçon, je sais lire et interpréter une courbe de croissance (poids/âge) d'un enfant.
\end{itemize}\end{multicols}}

\begin{sommaire}
\begin{multicols}{2}
\begin{enumerate}
\item Généralités sur les maladies nutritionnelles
\item La malnutrition protéino-énergétique (MPE)
\item Les avitaminoses et l'anémie
\item Les maladies de surcharge : l'obésité
\item Interprétation des données : courbe de croissance et IMC
\item Principes d'une alimentation préventive
\item Activité d'intégration
\item Méthodes et exercices résolus
\item Exercices
\item Corrigés des exercices
\item Fiche bilan
\end{enumerate}
\end{multicols}
\end{sommaire}

\begin{objectifs}
\begin{itemize}
\item À la fin de cette leçon, je sais définir une maladie nutritionnelle et distinguer maladies de carence et maladies de surcharge.
\item À la fin de cette leçon, je sais décrire les signes cliniques du kwashiorkor et du marasme à partir d'un cas concret.
\item À la fin de cette leçon, je sais citer les principales avitaminoses (A, C, D) et l'anémie ferriprive, avec leurs causes et leurs signes.
\item À la fin de cette leçon, je sais expliquer les causes et les conséquences de l'obésité sur la santé.
\item À la fin de cette leçon, je sais calculer et interpréter l'indice de masse corporelle (IMC) d'un individu.
\item À la fin de cette leçon, je sais lire et interpréter une courbe de croissance (poids/âge) d'un enfant.
\item À la fin de cette leçon, je sais proposer des mesures de prévention et de correction adaptées à chaque maladie nutritionnelle.
\item À la fin de cette leçon, je sais énoncer les principes d'une alimentation préventive équilibrée adaptée au contexte camerounais.
\end{itemize}
\end{objectifs}

\begin{prerequis}
\begin{itemize}
\item Les groupes d'aliments et leurs rôles (aliments énergétiques, de construction, protecteurs) — notions de début d'année
\item Les besoins nutritionnels selon l'âge, le sexe et l'activité
\item La notion de ration alimentaire équilibrée
\item Les notions de vitamines et de sels minéraux (fer, calcium, iode)
\end{itemize}
\end{prerequis}

\newpage

\coursec{Généralités sur les maladies nutritionnelles}

Dans un quartier de Yaoundé, deux enfants inquiètent leurs parents. Le premier, âgé de 3 ans, a les cheveux décolorés tirant sur le roux, le ventre ballonné et les jambes très maigres ; il pleure sans cesse et refuse de jouer. Le second, âgé de 15 ans, pèse 95~kg pour 1,60~m et s'essouffle au moindre effort ; monter un escalier le met en sueur. Au centre de santé intégré, l'infirmier examine les deux enfants et explique aux familles une chose surprenante : ces deux enfants souffrent de \cle{maladies nutritionnelles}, mais de nature \emph{opposée}. Le premier manque de nourriture adaptée, le second en consomme trop, et mal.

Comment deux problèmes si différents peuvent-ils avoir la même origine : une alimentation déséquilibrée ? Comment reconnaître ces maladies et surtout comment les prévenir dans nos familles ? C'est à ces questions que cette leçon va répondre, en commençant par poser des bases solides : qu'est-ce qu'une maladie nutritionnelle, comment les classe-t-on, quels facteurs les favorisent au Cameroun, et comment le personnel de santé les identifie-t-il ?

\courssub{Qu'est-ce qu'une maladie nutritionnelle ?}

\textbf{Intuition : la balance des apports et des besoins.}

Imagine ton organisme comme une maison en construction. Chaque jour, il faut des matériaux (protéines pour fabriquer les muscles et les défenses, lipides et glucides pour l'énergie, vitamines et sels minéraux pour faire fonctionner les « machines » cellulaires) et de l'énergie pour les ouvriers. Si les camions de matériaux n'arrivent pas en quantité suffisante, la construction ralentit puis s'arrête : c'est la \cle{carence}. Si au contraire les camions livrent sans cesse plus que nécessaire, les matériaux s'entassent, encombrent le chantier et finissent par l'étouffer : c'est la \cle{surcharge}. Dans les deux cas, la maison — ton corps — tombe malade.

\begin{definition}[Maladie nutritionnelle]
Une \cle{maladie nutritionnelle} est une maladie qui résulte d'un \textbf{déséquilibre durable entre les apports alimentaires (énergétiques et nutritionnels) et les besoins de l'organisme}. Ce déséquilibre peut porter sur la \emph{quantité} des aliments (trop ou pas assez d'énergie) ou sur leur \emph{qualité} (manque ou excès d'un nutriment précis : protéines, vitamines, fer, etc.).
\end{definition}

\begin{attention}[Ce qu'une maladie nutritionnelle n'est pas]
Une maladie nutritionnelle n'est pas une maladie \emph{contagieuse} : on ne « attrape » pas le kwashiorkor comme on attrape le paludisme. Elle n'est pas non plus une simple question de « manger beaucoup » : on peut manger en grande quantité et être quand même carencé, si le repas est pauvre en protéines, en vitamines ou en fer (par exemple un régime constitué presque uniquement de manioc ou de maïs bouillis).
\end{attention}

\begin{propriete}[Principe fondamental (admis)]
Tout déséquilibre \textbf{durable} entre les apports alimentaires et les besoins de l'organisme entraîne une pathologie : apports insuffisants $\Rightarrow$ maladie de \cle{carence} ; apports excessifs $\Rightarrow$ maladie de \cle{surcharge}. Cette propriété est admise : elle se justifie par le fait que l'organisme ne peut ni fonctionner sans matériaux ni stocker indéfiniment les excédents sans dommage.
\end{propriete}

\courssub{Classification des maladies nutritionnelles}

Reprenons nos deux enfants de Yaoundé. Le premier présente une \textbf{maladie de carence} : ses apports sont insuffisants (ici, surtout en protéines et en énergie). Le second présente une \textbf{maladie de surcharge} : ses apports énergétiques dépassent durablement ses dépenses. Toutes les maladies nutritionnelles se rangent ainsi en deux grandes catégories.

\begin{popfigure}
\begin{center}
\begin{tikzpicture}[
  grow=right,
  level 1/.style={sibling distance=3.8cm, level distance=4.5cm},
  level 2/.style={sibling distance=1.2cm, level distance=4.8cm},
  every node/.style={font=\small},
  racine/.style={rectangle, rounded corners, draw=popdark, fill=popblueL, thick, align=center, inner sep=5pt},
  cat/.style={rectangle, rounded corners, draw=popdark, thick, align=center, inner sep=4pt},
  ex/.style={rectangle, draw=popmuted, fill=popcream, rounded corners, align=center, font=\footnotesize, inner sep=3pt},
  edge from parent/.style={draw=popmuted, thick}
]
\node[racine, align=center]{Maladies\\ nutritionnelles}
  child { node[cat, fill=poporangeL, align=center]{Maladies de \textbf{surcharge}\\ (apports excessifs)}
    child { node[ex, align=center]{\textbf{Obésité}\\ (excès de graisses\\ et de sucres)} }
    child { node[ex, align=center]{Complications fréquentes :\\ diabète de type 2,\\ hypertension artérielle,\\ maladies cardiovasculaires} }
  }
  child { node[cat, fill=popgreenL, align=center]{Maladies de \textbf{carence}\\ (apports insuffisants)}
    child { node[ex, align=center]{\textbf{MPE}\\ (Malnutrition\\ Protéino-Énergétique :\\ marasme, kwashiorkor)} }
    child { node[ex, align=center]{\textbf{Avitaminoses}\\ (scorbut : vit. C ;\\ rachitisme : vit. D ;\\ xérophtalmie : vit. A)} }
    child { node[ex, align=center]{\textbf{Anémie ferriprive}\\ (manque de fer,\\ folates, vit. B12)} }
  };
\end{tikzpicture}
\end{center}
\legende{Arbre de classification des maladies nutritionnelles. Deux grandes catégories : les maladies de carence (apports insuffisants, en vert) et les maladies de surcharge (apports excessifs, en orange), avec des exemples qui seront détaillés dans les sections suivantes.}
\end{popfigure}

\begin{exemplebox}[Deux visages opposés du même problème]
\begin{itemize}
  \item \textbf{Carence} : un nourrisson sevré brutalement puis nourri presque exclusivement de bouillie de maïs reçoit peu de protéines et peu d'énergie. Il peut développer une \emph{malnutrition protéino-énergétique} (marasme ou kwashiorkor).
  \item \textbf{Surcharge} : un adolescent qui consomme chaque jour boissons sucrées, beignets, frites et plats très gras, tout en passant ses journées assis, accumule les graisses et devient \emph{obèse}.
\end{itemize}
Dans les deux cas, la cause profonde est la même : l'alimentation n'est pas adaptée aux besoins.
\end{exemplebox}

Le tableau suivant synthétise les différences entre les deux catégories.

\begin{center}
\begin{tabularx}{\linewidth}{|l|X|X|}
\hline
\rowcolor{popblueL} \textbf{Critère} & \textbf{Maladies de carence} & \textbf{Maladies de surcharge} \\
\hline
Cause principale & Apports alimentaires insuffisants (quantité ou qualité) & Apports excédentaires par rapport aux dépenses \\
\hline
Exemples principaux & MPE (marasme, kwashiorkor), avitaminoses, anémie ferriprive & Obésité et ses complications (diabète de type 2, hypertension) \\
\hline
Signes cliniques typiques & Amaigrissement sévère ou œdèmes, retard de croissance, fatigue, pâleur, cheveux décolorés & Excès de masse graisseuse, essoufflement à l'effort, IMC élevé \\
\hline
Populations les plus touchées & Enfants de moins de 5 ans, femmes enceintes et allaitantes, personnes en situation de précarité & Adultes urbains, adolescents sédentaires, mais de plus en plus les enfants \\
\hline
Prévention & Alimentation variée et suffisante, allaitement maternel exclusif, lutte contre les infections & Alimentation équilibrée, limitation des sucres ajoutés et graisses saturées, activité physique régulière \\
\hline
\end{tabularx}
\end{center}

\courssub{Facteurs favorisants au Cameroun}

Pourquoi ces maladies sont-elles fréquentes dans notre pays ? L'investigation épidémiologique (enquêtes de santé publique menées dans les régions, comme l'enquête EDS-MICS) met en évidence plusieurs facteurs qui agissent souvent \emph{ensemble}.

\begin{itemize}
  \item \textbf{La pauvreté et l'insécurité alimentaire} : quand le budget familial est très limité, l'accès aux aliments riches (viande, poisson, œufs, fruits, légumes variés) devient difficile ; le repas se réduit souvent à un féculent unique (manioc, maïs, riz) sans accompagnement protéiné ni légumier.
  \item \textbf{L'ignorance des règles d'équilibre alimentaire} : certaines familles disposent de ressources suffisantes mais composent mal les repas (bouillie trop diluée pour le nourrisson, absence de légumineuses, excès de produits industriels sucrés et salés).
  \item \textbf{Les infections répétées} : le paludisme, les parasitoses intestinales (vers) et les diarrhées aiguës répétées diminuent l'absorption des nutriments, augmentent les pertes et majorent les besoins métaboliques. Un enfant souvent malade se carence même en mangeant « normalement » : c'est le cercle vicieux \emph{malnutrition–infection}.
  \item \textbf{Le sevrage précoce ou mal conduit} : l'arrêt trop tôt de l'allaitement maternel exclusif (avant 6 mois), ou son remplacement par des bouillies pauvres en énergie et en protéines, expose le nourrisson aux carences.
  \item \textbf{La sédentarité et la transition nutritionnelle} : en ville (Yaoundé, Douala, Bafoussam), la diminution de l'activité physique (transports motorisés, écrans) et la consommation croissante de sodas, snacks, plats de rue très gras et sucrés favorisent la surcharge pondérale, dès l'adolescence.
\end{itemize}

\begin{savaistu}[Le double fardeau de la malnutrition au Cameroun]
Au Cameroun, selon les dernières enquêtes (EDS-MICS), environ \textbf{un enfant de moins de 5 ans sur trois} souffre de retard de croissance (retard statural, ou « \emph{stunting} ») : sa taille est nettement inférieure à celle attendue pour son âge, signe d'une sous-nutrition chronique. Ce retard n'est pas seulement une question de centimètres : il s'accompagne souvent d'un développement cognitif et scolaire diminué, avec des conséquences économiques à long terme. À l'inverse, le surpoids et l'obésité progressent rapidement chez les adultes des villes (près de 30\% des femmes en zone urbaine) et commencent à toucher les adolescents. Le pays fait donc face à un \emph{double fardeau} : carences et surcharges coexistent, parfois au sein d'une même famille, voire chez un même individu (obésité sarcopénique ou carences en micronutriments chez l'obèse).
\end{savaistu}

\courssub{Conséquences générales des déséquilibres nutritionnels}

Qu'elle soit de carence ou de surcharge, une maladie nutritionnelle ne reste jamais sans effet. Les conséquences générales sont :

\begin{itemize}
  \item le \textbf{retard de croissance} chez l'enfant (taille et poids insuffisants pour l'âge, retard statural et/ou pondéral) ;
  \item l'\textbf{affaiblissement du système immunitaire} : l'organisme carencé synthétise moins d'anticorps et de cellules immunitaires, d'où des infections plus fréquentes, plus longues et plus graves (rougeole, diarrhées, infections respiratoires aiguës) ;
  \item la \textbf{fatigue et la baisse des capacités physiques et intellectuelles} : difficultés de concentration à l'école, moindre résistance à l'effort physique, productivité au travail réduite ;
  \item les \textbf{complications chroniques} à l'âge adulte : l'obésité prépare le terrain au diabète de type 2, à l'hypertension artérielle et aux maladies cardiovasculaires (AVC, infarctus) ; les carences anciennes laissent parfois des séquelles irréversibles (retard mental par carence iodée ou ferrique sévère, cécité par xérophtalmie vitamine A, rachitisme séquellaire).
\end{itemize}

\begin{attention}[Piège fréquent au Probatoire]
Ne pas croire qu'un enfant « gros » est forcément « bien nourri ». Un enfant peut présenter un excès de graisses (surcharge pondérale) et souffrir en même temps de carences en fer, en zinc ou en vitamines si son alimentation est riche en calories vides (sucre, graisse) mais pauvre en nutriments essentiels : c'est la \emph{faim cachée}. De même, un ventre ballonné n'est pas un signe de bonne santé : chez l'enfant carencé, il peut traduire des œdèmes (kwashiorkor), une hépatomégalie ou une distension par parasites intestinaux (ascaridiose).
\end{attention}

\courssub{Comment identifier une maladie nutritionnelle ?}

Au centre de santé, l'infirmier ne devine pas : il suit une démarche rigoureuse, en trois étapes complémentaires, que tu dois savoir décrire et appliquer sur un cas concret.

\begin{methode}[Identifier une maladie nutritionnelle en 3 étapes]
\begin{enumerate}
  \item \textbf{Observer} les signes cliniques : aspect général (amaigrissement extrême ou obésité), état de la peau (sécheresse, dermatose), des cheveux (décoloration, finesse, signe du « drapeau » pour le kwashiorkor), des yeux (xérophtalmie, conjonctives pâles), présence d'œdèmes (gonflement bilatéral des pieds, puis visage, signe du godet), pâleur des conjonctives et des paumes (signe d'anémie), fatigue à l'effort.
  \item \textbf{Mesurer} (mesures anthropométriques) : le \emph{poids} (balance électronique ou mécanique), la \emph{taille} (tose pour l'enfant < 2 ans, stadiomètre pour l'enfant > 2 ans et l'adulte), le \emph{périmètre brachial} (PB) pour le dépistage rapide de la MPE chez l'enfant de 6 à 59 mois. Calculer l'\emph{Indice de Masse Corporelle} (IMC = Poids / Taille$^2$) pour l'adulte et l'adolescent, et \textbf{comparer les valeurs aux courbes de croissance de référence (OMS)} ou aux seuils standards (IMC pour l'âge chez l'adolescent).
  \item \textbf{Interroger} (anamnèse alimentaire et médicale) : questionner le patient ou ses parents sur la composition et la fréquence des repas (diversité alimentaire), l'allaitement (exclusif ? durée) et le sevrage (âge, qualité des aliments de complément), les infections récentes (paludisme, diarrhées, VIH/TB), l'activité physique et le mode de vie.
\end{enumerate}
La confrontation des trois étapes permet de conclure : type de maladie (carence ou surcharge), gravité (aiguë, modérée, sévère), et mesures de prise en charge ou de prévention à proposer.
\end{methode}

\begin{exoresolu}[À toi de jouer l'infirmier]
\textbf{Énoncé.} Deux enfants consultent au centre de santé intégré.
\begin{itemize}
  \item \textbf{Enfant A}, 3 ans : poids 9~kg (poids attendu vers 14~kg pour 3 ans), cheveux décolorés et fins, œdèmes bilatéraux des pieds, ventre ballonné, apathique. La mère explique qu'il a été sevré à 8 mois et mange surtout de la bouillie de maïs peu épaisse ; il a eu trois épisodes de diarrhée ce trimestre.
  \item \textbf{Enfant B}, 15 ans : 95~kg pour 1,60~m, essoufflé à l'effort. Il déclare boire un soda par jour, manger souvent beignets et frites, et ne pratiquer aucun sport.
\end{itemize}
Pour chaque enfant : (1) indique la catégorie de maladie nutritionnelle ; (2) cite deux arguments tirés de la démarche d'identification ; (3) propose une mesure adaptée.
\tcblower
\textbf{Solution détaillée.}

\textbf{Enfant A.} (1) Il s'agit d'une \textbf{maladie de carence}, plus précisément d'une \emph{malnutrition protéino-énergétique (MPE) de type kwashiorkor} (présence d'œdèmes + signes cutanéo-phanétaires). (2) \emph{Observation} : œdèmes bilatéraux des pieds, cheveux décolorés (signe du drapeau), ventre ballonné (hépatomégalie ou ascite), apathie. \emph{Mesure} : poids de 9~kg très inférieur au poids de référence pour l'âge (environ 14~kg), soit un rapport Poids/Âge < -3 ET (score Z), confirmant une malnutrition aiguë sévère. \emph{Interrogatoire} : sevrage précoce (8 mois), alimentation quasi exclusive de bouillie de maïs pauvre en protéines et en énergie (faible densité énergétique), diarrhées répétées qui aggravent les pertes et les besoins. (3) Mesures : prise en charge selon le protocole national/OMS (phase de stabilisation avec lait F-75, phase de réhabilitation avec lait F-100 ou alimentation locale enrichie type « pap enrichi » à base de maïs, soja/arachide, huile, sucre), traitement systématique des infections (antibiotiques, vermifuge, vitamine A, zinc), éducation nutritionnelle de la mère (fréquence, densité, diversité), suivi hebdomadaire du poids sur la courbe de croissance.

\textbf{Enfant B.} (1) Il s'agit d'une \textbf{maladie de surcharge} : l'obésité. (2) \emph{Observation} : excès de masse graisseuse visible, essoufflement à l'effort (intolérance à l'exercice). \emph{Mesure} : $\mathrm{IMC} = \dfrac{95}{1{,}60^2} = \dfrac{95}{2{,}56} \approx 37,1$~kg/m$^2$. Chez l'adolescent de 15 ans, cette valeur se situe bien au-dessus du percentile 97 (seuil d'obésité) sur les courbes de croissance IMC pour l'âge (OMS), confirmant l'obésité sévère. \emph{Interrogatoire} : alimentation hypercalorique déséquilibrée (excès de sucres simples et graisses saturées via sodas, fritures, beignets), sédentarité totale. (3) Mesures : suppression des boissons sucrées et réduction drastique des fritures/beignets, rééquilibrage alimentaire (augmentation légumes, fruits, protéines maigres, céréales complètes), prescription d'une activité physique quotidienne adaptée (marche rapide 30--45 min/jour), bilan biologique à la recherche de complications (glycémie à jeun, profil lipidique, transaminases, urée/créatinine), suivi médical régulier (poids, IMC, TA) et soutien psychologique si nécessaire.
\end{exoresolu}

\begin{aretenir}[L'essentiel de cette section]
\begin{itemize}
  \item Une \cle{maladie nutritionnelle} résulte d'un déséquilibre durable entre les apports alimentaires (quantité et qualité) et les besoins de l'organisme.
  \item On distingue les \textbf{maladies de carence} (MPE : marasme/kwashiorkor, avitaminoses, anémie ferriprive) et les \textbf{maladies de surcharge} (obésité et ses complications métaboliques et cardiovasculaires).
  \item Au Cameroun, elles sont favorisées par la pauvreté/insécurité alimentaire, l'ignorance de l'équilibre alimentaire, le couple malnutrition-infection (paludisme, parasitoses, diarrhées), le sevrage précoce/mal conduit, la sédentarité et la transition nutritionnelle (produits ultra-transformés).
  \item Leurs conséquences : retard de croissance (statural/pondéral), affaiblissement immunitaire, fatigue/baisse des performances, complications chroniques irréversibles.
  \item L'identification clinique repose sur la triade : \textbf{Observer} (signes physiques), \textbf{Mesurer} (anthropométrie, IMC, courbes de référence OMS), \textbf{Interroger} (anamnèse alimentaire et infectieuse).
\end{itemize}
\end{aretenir}

Maintenant que les bases sont posées, la section suivante étudiera en détail la première grande catégorie : la malnutrition protéino-énergétique (MPE), qui frappe encore trop d'enfants camerounais.

\coursec{La malnutrition protéino-énergétique (MPE)}

Dans la section précédente, nous avons vu que les maladies nutritionnelles regroupent deux grandes familles : les maladies de \cle{carence} (l'organisme manque de quelque chose) et les maladies de \cle{surcharge} (l'organisme reçoit trop). Parmi les maladies de carence, la plus redoutable chez le jeune enfant est sans conteste la \cle{malnutrition protéino-énergétique}. C'est elle que nous étudions maintenant en détail, car elle reste, dans de nombreuses régions du Cameroun et d'Afrique, une cause majeure de mortalité infantile.

\courssub{Une observation qui pose problème}

Dans un centre de santé intégré de l'Extrême-Nord, deux enfants de 2 ans sont amenés le même jour. Le premier, Alima, a le visage et les pieds gonflés, un ventre très rond, des cheveux roux et cassants ; il est apathique, il ne joue plus. Le second, Bouba, est squelettique : ses côtes se dessinent sous la peau, ses bras sont aussi fins que des bâtons, son visage ressemble à celui d'un vieillard ; mais il n'a aucun gonflement. Les deux enfants pèsent moins de \SI{7}{\kg}, alors qu'un enfant de 2 ans devrait peser environ \SI{12}{\kg}.

Deux questions se posent immédiatement :
\begin{itemize}
\item Ces deux enfants souffrent-ils de la même maladie, alors que leurs aspects sont si différents ?
\item Pourquoi des enfants qui mangent pourtant « quelque chose » chaque jour (une bouillie de maïs ou de mil) tombent-ils ainsi malades ?
\end{itemize}

C'est en répondant à ces questions que nous allons construire la notion de malnutrition protéino-énergétique.

\courssub{Définition de la malnutrition protéino-énergétique}

\begin{definition}[Malnutrition protéino-énergétique (MPE)]
La \cle{malnutrition protéino-énergétique} (MPE) est un ensemble de maladies de carence résultant d'un apport insuffisant en \cle{protéines} et/ou en \cle{énergie} (glucides et lipides). Elle touche surtout les \textbf{enfants de 0 à 5 ans}, période où les besoins de croissance sont élevés. Elle se manifeste par un arrêt ou un ralentissement de la croissance pondérale (le poids n'augmente plus) puis staturo-pondérale (la taille aussi est affectée).
\end{definition}

L'intuition est simple : le corps d'un enfant est un « chantier » en construction permanente. Les protéines sont les \textbf{matériaux} (briques) de ce chantier — elles servent à fabriquer les muscles, le sang, la peau, les anticorps — tandis que les glucides et les lipides fournissent l'\textbf{énergie} (le carburant) nécessaire au chantier et à toutes les fonctions vitales. Si les briques manquent, le chantier s'arrête ; si le carburant manque aussi, l'organisme détruit ses propres réserves (graisses, puis muscles) pour survivre.

\begin{aretenir}[Les deux formes cliniques majeures]
Selon que la carence porte surtout sur les protéines ou sur l'ensemble protéines + énergie, la MPE prend deux visages :
\begin{itemize}
\item le \cle{kwashiorkor} : carence \textbf{protéique dominante} (l'enfant reçoit à peu près assez de calories, mais presque pas de protéines) ;
\item le \cle{marasme} : carence \textbf{globale}, à la fois énergétique et protéique (l'enfant manque de tout).
\end{itemize}
Il existe aussi des formes intermédiaires, le \cle{marasme-kwashiorkor}.
\end{aretenir}

\courssub{Le kwashiorkor : la carence protéique dominante}

\textbf{Situation typique.} Un enfant est allaité correctement pendant un an. Puis naît un petit frère : l'aîné est brutalement sevré et mis à la bouillie de maïs ou de manioc, aliment énergétique mais très pauvre en protéines. Quelques mois plus tard, il gonfle.

\begin{definition}[Kwashiorkor]
Le \cle{kwashiorkor} est la forme de MPE due à une carence dominante en \textbf{protéines}, avec un apport énergétique relativement conservé (alimentation à base de bouillies et de tubercules). Il apparaît classiquement entre \textbf{1 et 3 ans}, souvent après un sevrage brutal.
\end{definition}

\textbf{Les signes caractéristiques} sont :
\begin{itemize}
\item des \cle{œdèmes} : gonflements d'abord aux pieds et aux jambes, puis au visage (visage « de lune »). Mécanisme : les protéines du sang (albumine) retiennent l'eau dans les vaisseaux ; quand elles manquent, l'eau s'échappe vers les tissus ;
\item un \cle{ventre ballonné} : dû aux œdèmes de la paroi intestinale, à l'accumulation de liquide dans l'abdomen (ascite) et parfois à une hépatomégalie (gros foie stéatosique) ;
\item des \textbf{cheveux décolorés} (roux ou jaunâtres), fins, fragiles, qui s'arrachent facilement (signe du « drapeau ») ;
\item une \textbf{peau craquelée}, qui se desquame par plaques (« peau de serpent » ou « peau écaillée »), surtout aux jambes et aux fesses ;
\item l'\cle{apathie} : l'enfant est triste, indifférent, ne joue plus, ne réclame plus ;
\item paradoxalement, l'enfant ne paraît pas « maigre » au premier regard, car les œdèmes masquent la fonte musculaire.
\end{itemize}

\begin{attention}[Erreur fréquente : le « gros bébé » qui est en danger]
Ne jamais confondre le \textbf{ventre ballonné} et les joues gonflées du kwashiorkor avec un bon état nutritionnel ! Dans certaines familles, on se réjouit que l'enfant soit « bien portant » alors qu'il est en réalité gravement malade. Le ventre rond du kwashiorkor est un \textbf{signe de gravité} : c'est de l'eau et non de la graisse. Un enfant au ventre ballonné, aux pieds œdématiés et aux cheveux décolorés doit être conduit d'urgence au centre de santé.
\end{attention}

\begin{savaistu}[D'où vient le mot « kwashiorkor » ?]
Le mot \emph{kwashiorkor} vient d'une langue du Ghana (le ga) et signifie littéralement « la maladie de l'enfant sevré quand arrive le suivant ». Ce nom décrit à lui seul la cause la plus fréquente : l'enfant est chassé du sein maternel à la naissance d'un nouveau bébé et nourri de bouillies pauvres en protéines. C'est la pédiatre jamaïcaine Cicely Williams qui, en 1935, a introduit ce terme dans la littérature médicale après l'avoir observée en Afrique de l'Ouest.
\end{savaistu}

\courssub{Le marasme : la carence globale}

\begin{definition}[Marasme]
Le \cle{marasme} est la forme de MPE due à une carence \textbf{globale}, à la fois en énergie et en protéines : l'enfant manque de tout, souvent parce qu'il n'est tout simplement pas assez nourri (allaitement insuffisant ou absent, bouillies trop diluées, infections répétées qui coupent l'appétit). Il survient plutôt \textbf{avant 1 an}.
\end{definition}

\textbf{Les signes caractéristiques} sont :
\begin{itemize}
\item une \cle{maigreur extrême} : les graisses puis les muscles ont été consommés par l'organisme pour produire de l'énergie (autophagie) ;
\item un \textbf{visage de vieillard} : joues creuses, pommettes saillantes, yeux enfoncés ;
\item les \textbf{côtes apparentes}, les bras et les jambes réduits à « la peau et les os », les fesses atrophiées (« fesses de vieillard ») ;
\item l'\textbf{absence d'œdèmes} : c'est le critère qui distingue le marasme du kwashiorkor ;
\item l'enfant reste souvent \textbf{alerte} et réclame à manger, contrairement à l'enfant kwashiorkor qui est apathique.
\end{itemize}

\textbf{Formes mixtes.} Dans la réalité, de nombreux enfants présentent un tableau intermédiaire : le \cle{marasme-kwashiorkor}, qui associe une maigreur sévère et des œdèmes. C'est la forme la plus grave, au pronostic le plus sombre.

\begin{popfigure}
\begin{center}
\begin{tikzpicture}[scale=0.9]
% ===== KWASHIORKOR =====
\begin{scope}
% tete gonflee
\fill[poporangeL] (0,3.4) circle (0.62);
% cheveux decolores
\draw[popgold,thick] (-0.45,3.85) -- (-0.55,4.15);
\draw[popgold,thick] (-0.15,4.0) -- (-0.2,4.3);
\draw[popgold,thick] (0.2,4.0) -- (0.25,4.3);
\draw[popgold,thick] (0.5,3.85) -- (0.6,4.15);
% yeux tristes
\fill[popdark] (-0.22,3.5) circle (0.05);
\fill[popdark] (0.22,3.5) circle (0.05);
\draw[popdark] (-0.12,3.15) .. controls (0,3.05) .. (0.12,3.15);
% corps avec ventre ballonne
\fill[poporangeL] (0,1.9) ellipse (0.95 and 1.15);
% bras
\draw[poporange,very thick] (-0.75,2.6) -- (-1.15,1.6);
\draw[poporange,very thick] (0.75,2.6) -- (1.15,1.6);
% jambes oedematees (epaisses)
\draw[poporange,line width=6pt] (-0.4,0.85) -- (-0.45,-0.4);
\draw[poporange,line width=6pt] (0.4,0.85) -- (0.45,-0.4);
% pieds oedematees
\fill[poporange] (-0.45,-0.5) ellipse (0.28 and 0.14);
\fill[poporange] (0.45,-0.5) ellipse (0.28 and 0.14);
% peau craquelee
\draw[popdark!60] (-0.3,1.3) -- (-0.1,1.15) -- (0.15,1.3);
\draw[popdark!60] (0.1,2.2) -- (0.3,2.05) -- (0.5,2.2);
% legendes pointees
\draw[popmuted,thin] (0.62,3.6) -- (2.2,4.3) node[right,align=left,font=\small]{visage gonflé (\cle{œdème})};
\draw[popmuted,thin] (0.95,2.0) -- (2.2,2.6) node[right,align=left,font=\small]{ventre ballonné};
\draw[popmuted,thin] (-0.2,4.1) -- (-1.6,4.6) node[left,align=right,font=\small]{cheveux décolorés, fragiles};
\draw[popmuted,thin] (-0.45,-0.3) -- (-1.9,-0.9) node[left,align=right,font=\small]{pieds œdématiés};
\draw[popmuted,thin] (0.35,1.2) -- (2.2,0.6) node[right,align=left,font=\small]{peau craquelée};
\node[font=\bfseries] at (0,-1.3) {Kwashiorkor};
\end{scope}
% ===== MARASME =====
\begin{scope}[xshift=8.5cm]
% tete de vieillard
\fill[popblueL] (0,3.4) circle (0.5);
\draw[popdark] (-0.18,3.45) circle (0.06);
\draw[popdark] (0.18,3.45) circle (0.06);
\draw[popdark] (-0.1,3.1) .. controls (0,3.0) .. (0.1,3.1);
% rides
\draw[popdark!50] (-0.3,3.7) -- (0.3,3.7);
\draw[popdark!50] (-0.25,3.6) -- (0.25,3.6);
% corps squelettique
\fill[popblueL] (0,1.9) ellipse (0.42 and 1.1);
% cotes apparentes
\draw[popdark!60] (-0.38,2.3) .. controls (0,2.15) .. (0.38,2.3);
\draw[popdark!60] (-0.4,2.05) .. controls (0,1.9) .. (0.4,2.05);
\draw[popdark!60] (-0.38,1.8) .. controls (0,1.65) .. (0.38,1.8);
% bras filiformes
\draw[popblue,thick] (-0.4,2.6) -- (-0.85,1.5);
\draw[popblue,thick] (0.4,2.6) -- (0.85,1.5);
% jambes filiformes
\draw[popblue,thick] (-0.2,0.85) -- (-0.3,-0.45);
\draw[popblue,thick] (0.2,0.85) -- (0.3,-0.45);
% legendes pointees
\draw[popmuted,thin] (0.5,3.5) -- (2.2,4.3) node[right,align=left,font=\small]{visage de vieillard};
\draw[popmuted,thin] (0.42,2.1) -- (2.2,2.6) node[right,align=left,font=\small]{côtes apparentes};
\draw[popmuted,thin] (-0.85,1.9) -- (-2.0,1.4) node[left,align=right,font=\small]{bras filiformes\\ (fonte musculaire)};
\draw[popmuted,thin] (0.3,-0.2) -- (2.2,-0.6) node[right,align=left,font=\small]{\textbf{pas d'œdème}};
\node[font=\bfseries] at (0,-1.3) {Marasme};
\end{scope}
\end{tikzpicture}
\end{center}
\legende{Comparaison schématique des deux formes majeures de MPE. À gauche, l'enfant atteint de kwashiorkor : œdèmes (pieds, visage), ventre ballonné, cheveux décolorés, peau craquelée. À droite, l'enfant marasmique : maigreur extrême, côtes visibles, visage de vieillard, sans œdème.}
\end{popfigure}

\begin{center}
\begin{tabularx}{\linewidth}{|l|X|X|}
\hline
\rowcolor{popblueL} \textbf{Critère} & \textbf{Kwashiorkor} & \textbf{Marasme} \\
\hline
Cause dominante & Carence en \textbf{protéines} (énergie à peu près suffisante : bouillies) & Carence \textbf{globale} : énergie + protéines (sous-alimentation totale) \\
\hline
Œdèmes & \textbf{Présents} (pieds, visage) & \textbf{Absents} \\
\hline
Aspect général & Gonflé, ventre ballonné, fonte musculaire masquée & Maigreur extrême, « peau et os », côtes apparentes \\
\hline
Cheveux et peau & Cheveux décolorés (roux), fragiles ; peau craquelée & Cheveux peu modifiés ; peau sèche, plissée \\
\hline
Comportement & Apathie, refus de manger & Enfant alerte, appétit souvent conservé \\
\hline
Âge habituel & 1 à 3 ans (après le sevrage) & Avant 1 an (mais possible plus tard) \\
\hline
\end{tabularx}
\end{center}

\courssub{Démarche d'investigation : comment le corps réagit-il au manque de nourriture ?}

\begin{experience}[Comprendre la différence entre kwashiorkor et marasme]
\textbf{Question.} Pourquoi deux enfants carencés présentent-ils des aspects opposés (l'un gonflé, l'autre squelettique) ?

\textbf{Hypothèse.} L'aspect dépend de la nature du nutriment qui manque le plus.

\textbf{Données expérimentales et cliniques.}
\begin{itemize}
\item Chez l'animal de laboratoire nourri avec un régime riche en glucides mais quasi sans protéines, on observe une baisse de l'\cle{albumine} sanguine (protéine fabriquée par le foie) et l'apparition d'œdèmes.
\item Chez l'animal soumis à une restriction calorique totale, on observe une fonte des graisses puis des muscles, sans œdème.
\item Les dosages sanguins chez les enfants hospitalisés confirment : hypoalbuminémie marquée dans le kwashiorkor, proche de la normale dans le marasme.
\end{itemize}

\textbf{Interprétation.} L'albumine retient l'eau dans les vaisseaux sanguins (pression oncotique). Quand les protéines alimentaires manquent, le foie ne peut plus fabriquer assez d'albumine : l'eau quitte le sang et s'accumule dans les tissus → œdèmes. En revanche, quand c'est l'énergie qui manque totalement, l'organisme « brûle » ses réserves : d'abord les graisses, puis les protéines musculaires → maigreur extrême.

\textbf{Conclusion.} Le kwashiorkor traduit une carence protéique dominante (avec œdèmes), le marasme une carence énergétique et protéique globale (sans œdèmes). L'hypothèse est validée.
\end{experience}

\courssub{Les causes de la MPE}

La MPE n'est jamais une fatalité : elle résulte de causes identifiées, souvent combinées.

\begin{enumerate}
\item \textbf{Le sevrage brutal.} Le passage soudain du lait maternel (riche et équilibré) à une bouillie pauvre est le déclencheur classique du kwashiorkor, surtout lorsqu'une nouvelle grossesse survient tôt.
\item \textbf{Une alimentation pauvre en protéines.} Les \textbf{bouillies exclusives} de maïs, de mil ou de manioc, surtout si elles sont trop diluées, apportent de l'énergie mais très peu de protéines de bonne qualité.
\item \textbf{Les infections répétées.} Paludisme, diarrhées, rougeole, infections respiratoires : chaque infection augmente les besoins, coupe l'appétit et fait perdre des protéines. C'est un cercle vicieux : la malnutrition favorise les infections, qui aggravent la malnutrition.
\item \textbf{La précocité des grossesses chez l'adolescente.} Une mère très jeune, elle-même en croissance et souvent mal nourrie, donne naissance à des enfants de faible poids, plus vulnérables à la MPE. Les naissances rapprochées épuisent aussi les réserves de la mère.
\item \textbf{Facteurs aggravants} : pauvreté, ignorance des règles d'alimentation du nourrisson, sécheresses et insécurité alimentaire dans certaines régions (Extrême-Nord, Nord, Adamaoua).
\end{enumerate}

\courssub{Les conséquences de la MPE}

\begin{itemize}
\item \textbf{Retard de croissance} : l'enfant ne prend plus de poids, puis sa taille se fige (retard staturo-pondéral parfois irréversible si la carence survient avant 2 ans).
\item \textbf{Retard du développement intellectuel} : le cerveau se construit massivement avant 5 ans ; une carence prolongée à cet âge compromet les apprentissages futurs (déficit cognitif).
\item \textbf{Sensibilité accrue aux infections} : la fabrication des anticorps (immunoglobulines) exige des protéines ; l'enfant carencé est sans défense face à la rougeole, aux diarrhées, aux pneumopathies.
\item \textbf{Mortalité infantile élevée} : la MPE sévère non traitée est souvent mortelle, directement ou par les infections qu'elle favorise. Elle reste l'une des premières causes de décès des moins de 5 ans dans le monde.
\end{itemize}

\begin{exemplebox}[Un chiffre pour comprendre : les besoins d'un enfant de 2 ans]
Un enfant de 2 ans doit recevoir environ \SI{1000}{kcal} par jour et des \textbf{protéines à chaque repas} (\SIrange{1}{2}{\gram\per\kilogram\per\day}). Or une bouillie de maïs seule, surtout diluée, apporte peu de calories et presque pas de protéines : il faudrait en boire un volume impossible pour l'estomac d'un petit enfant. En revanche, enrichir cette bouillie avec de la \textbf{pâte d'arachide}, du \textbf{lait}, de la \textbf{farine de soja} ou accompagner le repas de \textbf{niébé} ou de \textbf{poisson} couvre les besoins en un volume adapté. C'est toute la logique de la diversification alimentaire.
\end{exemplebox}

\begin{methode}[Identifier une MPE à partir d'un cas décrit]
Face à un cas clinique (exercice ou situation réelle), procéder par étapes :
\begin{enumerate}
\item \textbf{Relever l'âge} : avant 1 an → penser marasme ; 1 à 3 ans après sevrage → penser kwashiorkor.
\item \textbf{Chercher les œdèmes} (pieds, visage) et le \textbf{ventre ballonné} : présents → kwashiorkor ; absents avec maigreur extrême → marasme.
\item \textbf{Examiner cheveux et peau} : cheveux roux fragiles, peau craquelée → kwashiorkor.
\item \textbf{Comparer le poids à l'âge} (ou à la taille) sur la courbe de croissance : un poids inférieur à \SI{80}{\%} du poids médian de référence (classification de Gomez) signe la malnutrition ; moins de \SI{60}{\%} ou présence d'œdèmes → forme sévère (selon critères OMS : score Z < -3 ou œdèmes).
\item \textbf{Conclure} en nommant la forme (kwashiorkor, marasme ou forme mixte) et en justifiant par les signes relevés.
\end{enumerate}
\end{methode}

\begin{exoresolu}[Identifier la forme de MPE]
\textbf{Énoncé.} Salam, 18 mois, est amené au centre de santé. Sa mère raconte qu'il a été sevré à 10 mois à la naissance de sa petite sœur et nourri depuis de bouillie de maïs. À l'examen : poids \SI{6,5}{\kg} (au lieu d'environ \SI{10}{\kg}), pieds et visage gonflés, ventre proéminent, cheveux roux cassants, enfant prostré. 1) Identifier la maladie et sa forme. 2) Justifier le diagnostic. 3) Expliquer l'origine des œdèmes.
\tcblower
\textbf{Solution détaillée.}
1) Il s'agit d'une \textbf{malnutrition protéino-énergétique}, forme \textbf{kwashiorkor}.

2) Justification par les signes : présence d'\textbf{œdèmes} (pieds et visage gonflés), \textbf{ventre ballonné}, \textbf{cheveux décolorés et fragiles}, \textbf{apathie}, âge typique (18 mois, après un sevrage brutal), alimentation exclusive de bouillie pauvre en protéines. Le poids très inférieur à la normale confirme la malnutrition sévère.

3) Origine des œdèmes : la bouillie de maïs seule n'apporte pas assez de protéines. Le foie ne peut plus fabriquer suffisamment d'\textbf{albumine}, protéine sanguine qui retient l'eau dans les vaisseaux par la pression oncotique. L'eau s'échappe alors du sang et s'accumule dans les tissus : ce sont les œdèmes.
\end{exoresolu}

\courssub{La prévention de la MPE}

La MPE se prévient efficacement par des mesures simples et à la portée de toutes les familles :

\begin{enumerate}
\item \textbf{L'allaitement maternel exclusif pendant les 6 premiers mois} : le lait maternel couvre seul tous les besoins du nourrisson (énergie, protéines, anticorps, facteurs de croissance) ; il est gratuit, sûr et toujours disponible. L'Organisation mondiale de la santé recommande de le poursuivre, en complément, jusqu'à 2 ans ou plus.
\item \textbf{La diversification alimentaire à partir de 6 mois} : introduire progressivement, en plus du lait, des bouillies \textbf{enrichies} et des sources de protéines locales et bon marché : \cle{niébé}, \cle{arachide}, \cle{soja}, \cle{poisson} (sèche ou frais), œufs, lait, haricots. L'association céréale + légumineuse (maïs + niébé, par exemple) fournit des protéines de bonne qualité (profil d'acides aminés complémentaire).
\item \textbf{Le suivi de la croissance} : pesée régulière de l'enfant (mensuelle idéalement) et report sur la \textbf{courbe de croissance} du carnet de santé. Un poids qui stagne deux mois de suite, ou une courbe qui fléchit (casse), est un signal d'alarme qui doit conduire au centre de santé \emph{avant} l'apparition des signes graves.
\item \textbf{Mesures associées} : espacement des naissances (planification familiale), vaccination complète et traitement rapide des infections (paludisme, diarrhées, infections respiratoires), éducation nutritionnelle des mères et des adolescentes, amélioration de l'hygiène de l'eau et des aliments.
\end{enumerate}

\begin{attention}[Pièges à éviter dans la rédaction]
\begin{itemize}
\item Ne pas écrire que le kwashiorkor est dû à un manque de nourriture en général : c'est une carence \textbf{protéique dominante} ; l'enfant mange, mais mal (qualité).
\item Ne pas attribuer des œdèmes au marasme : leur \textbf{absence} est justement son critère distinctif majeur.
\item Ne pas limiter la prévention à « bien manger » : il faut citer des mesures précises (allaitement exclusif 6 mois, diversification avec protéines locales, suivi de la courbe de croissance, vaccination, espacement des naissances).
\item Ne pas confondre « protéines » et « viande » : les protéines végétales bien associées (céréales + légumineuses) sont une alternative efficace et accessible.
\end{itemize}
\end{attention}

\begin{exercice}[À toi de jouer]
Dans un village, on observe deux enfants de 20 mois. Enfant A : poids \SI{7}{\kg}, maigreur extrême, côtes visibles, pas d'œdème, réclame à manger. Enfant B : poids \SI{7,5}{\kg}, pieds gonflés, ventre ballonné, cheveux jaunâtres, apathique. 1) Identifier la forme de MPE de chaque enfant en justifiant. 2) Proposer trois mesures de prévention à expliquer aux mères du village.
\end{exercice}

\begin{corrige}[Exercice 1]
1) Enfant A : \textbf{marasme} — maigreur extrême, côtes apparentes, absence d'œdèmes, appétit conservé : carence globale (énergie + protéines). Enfant B : \textbf{kwashiorkor} — œdèmes, ventre ballonné, cheveux décolorés, apathie : carence protéique dominante.

2) Mesures de prévention : allaitement maternel exclusif pendant 6 mois puis poursuivi avec diversification ; enrichissement des bouillies avec des protéines locales (niébé, arachide, soja, poisson) ; pesée mensuelle et suivi de la courbe de croissance au centre de santé ; traitement rapide des infections et espacement des naissances.
\end{corrige}

En résumé, la MPE est une maladie de carence qui porte deux visages — le kwashiorkor (manque de protéines, avec œdèmes) et le marasme (manque de tout, sans œdèmes) — mais une seule logique : le jeune enfant ne reçoit pas ce dont son corps en construction a besoin. D'autres carences, plus ciblées, frappent aussi les enfants et les adultes : ce sont les avitaminoses et l'anémie, que nous étudions dans la section suivante.

\coursec{Les avitaminoses et l'anémie}

Dans la section précédente, nous avons étudié la malnutrition protéino-énergétique (MPE) : elle résulte d'un déficit global en énergie et en protéines, et frappe surtout les jeunes enfants sous la forme du kwashiorkor ou du marasme. Mais une personne peut manger à sa faim — en quantité — et pourtant tomber malade, car son alimentation manque d'un élément précis : une vitamine, un minéral. Ce sont les \cle{maladies de carence qualitative}, beaucoup plus insidieuses car la faim n'est pas toujours ressentie. C'est à elles que cette section est consacrée, à travers les \cle{avitaminoses} et l'\cle{anémie ferriprive}.

\courssub{Notion d'avitaminose}

\textbf{Une observation pour commencer.} Dans un village de l'Extrême-Nord, une maman remarque que son fils de 6 ans, pourtant bien portant le jour, se cogne partout dès que la nuit tombe et ne reconnaît plus les personnes à la tombée du jour. Au centre de santé, l'infirmier apprend que l'enfant mange presque exclusivement du mil bouilli, sans huile, sans légumes, sans fruits. Le diagnostic est posé : \emph{cécité nocturne}, premier signe d'une carence en vitamine A. L'enfant n'a pas faim, il mange tous les jours — mais il lui manque une molécule indispensable.

\begin{definition}[Avitaminose]
Une \cle{avitaminose} est une maladie due à une \textbf{carence prolongée en une vitamine donnée}. Les vitamines sont des substances organiques indispensables au fonctionnement de l'organisme, nécessaires en très petites quantités, et que le corps humain ne sait pas fabriquer (sauf exception, comme la vitamine D partiellement synthétisée sous l'effet des UVB solaires). Elles doivent donc être apportées par l'alimentation. On parle d'\emph{hypovitaminose} lorsque la carence est partielle.
\end{definition}

\begin{attention}[Ne pas confondre]
Une avitaminose n'est pas une maladie infectieuse : elle n'est causée ni par un microbe, ni par un parasite, et elle \textbf{ne se transmet pas}. Elle se prévient et se corrige par l'alimentation (ou par des suppléments prescrits). Autre piège : les vitamines ne fournissent \textbf{aucune énergie} — ce ne sont ni des sucres, ni des graisses — mais elles sont indispensables comme « clés » (coenzymes) de nombreuses réactions chimiques du corps.
\end{attention}

\courssub{La carence en vitamine A : de la cécité nocturne à la xérophtalmie}

\textbf{Situation concrète.} Reprenons le cas de l'enfant du village. Sa cécité nocturne (ou \emph{héméralopie}) s'explique ainsi : la vitamine A (rétinol) est indispensable à la fabrication de la \textbf{rhodopsine}, le pigment photosensible des bâtonnets de la rétine, ces cellules qui permettent la vision en lumière faible (vision scotopique). Sans vitamine A, la régénération de la rhodopsine est bloquée : la vision crépusculaire devient impossible. Si la carence se poursuit, l'épithélium conjonctival et cornéen se kératinise et se dessèche : c'est la \cle{xérophtalmie}, caractérisée par l'apparition de \emph{taches de Bitot} (dépôts mousseux blanchâtres sur la conjonctive), qui peut aboutir au ramollissement et à la nécrose de la cornée (\emph{kératomalacie}) puis à la \textbf{cécité définitive et irréversible}. La vitamine A joue aussi un rôle majeur dans l'immunité (intégrité des muqueuses, différenciation des lymphocytes) : les enfants carencés attrapent plus facilement rougeole et diarrhées sévères.

\textbf{Démarche d'investigation (modèle de raisonnement au Probatoire).}
\begin{enumerate}
\item \emph{Observation} : l'enfant voit mal la nuit (héméralopie) ; son alimentation est pauvre en corps gras et en produits animaux/colorés.
\item \emph{Hypothèse} : carence en vitamine A (rétinol / bêta-carotène).
\item \emph{Vérification} : examen ophtalmologique (conjonctive sèche, taches de Bitot, kératomalacie), enquête alimentaire (fréquence de consommation d'huile de palme rouge, foie, légumes orangés/verts).
\item \emph{Interprétation} : les signes oculaires cliniques et le régime alimentaire concordent avec l'hypothèse.
\item \emph{Conclusion et traitement} : supplémentation orale en vitamine A (haute dose selon protocole OMS) et correction durable du régime ; la cécité nocturne régresse en quelques jours/semaines si elle est prise au stade précoce (héméralopie, xérophtalmie conjonctivale), mais la kératomalacie laisse des séquelles irréversibles.
\end{enumerate}

\textbf{Sources alimentaires.} La vitamine A se trouve sous deux formes : le \emph{rétinol} (forme préformée, d'origine animale : foie, œufs, lait entier, beurre, poisson, huile de foie de morue) et les \emph{caroténoïdes} (provitamine A, d'origine végétale), dont le principal est le \textbf{bêta-carotène}, pigment orange-rouge transformé en rétinol dans la muqueuse intestinale : \textbf{huile de palme rouge non raffinée, carotte, mangue, patate douce à chair orangée, feuilles vertes foncées} (folon, feuilles de manioc, morelle noire, amarante).

\begin{savaistu}[L'huile de palme rouge, un trésor camerounais]
L'huile de palme rouge, très consommée au Cameroun (sauce jaune, mbongo, eru, ndolé...), est l'une des \textbf{meilleures sources naturelles de provitamine A} au monde : sa couleur rouge-orangée caractéristique vient précisément de sa richesse exceptionnelle en bêta-carotène (et autres caroténoïdes). Attention toutefois : une cuisson trop prolongée ou à feu trop vif détruit une partie significative des caroténoïdes thermosensibles ; il est préférable de l'ajouter en fin de cuisson. C'est aussi pour lutter contre la carence en vitamine A, problème de santé publique majeur, que le Ministère de la Santé Publique, avec l'appui de l'OMS et de l'UNICEF, organise des campagnes de supplémentation en vitamine A (haute dose) tous les 6 mois pour les enfants de 6 à 59 mois, souvent couplées aux journées de vaccination ou de déparasitage.
\end{savaistu}

\courssub{La carence en vitamine C : le scorbut}

\textbf{Situation concrète.} Historiquement, le scorbut décimait les équipages des navires au long cours (XVe–XVIIIe siècles), privés de fruits et légumes frais pendant des mois : gencives gonflées, violacées, qui saignent spontanément, dents qui se déchaussent et tombent, fatigue intense, hématomes spontanés, plaies qui ne cicatrisent plus, douleurs articulaires invalidantes. Aujourd'hui encore, on observe des formes frustes ou infantiles (maladie de Barlow) chez des personnes au régime très pauvre en produits frais (personnes âgées isolées, alcooliques, nourrissons non diversifiés, contextes de crise).

\begin{definition}[Scorbut]
Le \cle{scorbut} est la maladie due à la carence en \textbf{vitamine C} (acide L-ascorbique). Ses signes cliniques principaux sont : des \textbf{gingivites hémorragiques} (gencives enflées, saignantes), une \textbf{asthénie} marquée, des \textbf{ecchymoses} et hématomes spontanés (fragilité capillaire), un \textbf{retard de cicatrisation}, des \textbf{douleurs ostéo-articulaires} et un arrêt de la croissance osseuse chez l'enfant.
\end{definition}

\textbf{Explication biophysique.} La vitamine C est un cofacteur indispensable des \textbf{prolyl- et lysyl-hydroxylases}, enzymes qui assurent l'hydroxylation de la proline et de la lysine lors de la synthèse du \textbf{collagène}, la protéine structurelle majeure qui « cimente » les tissus conjonctifs (gencives, parodonte, paroi vasculaire, peau, os, cartilage). Sans elle, le collagène synthétisé est défectueux (instable thermiquement), les vaisseaux deviennent fragiles (purpura), les tissus de soutien se relâchent (chute des dents, ouverture des cicatrices). La vitamine C est aussi un puissant \textbf{antioxydant} hydrosoluble et \textbf{améliore l'absorption du fer non héminique} (fer végétal, Fe$^{3+}$ $\to$ Fe$^{2+}$) par réduction — un détail capital pour la prévention de l'anémie ferriprive (voir plus loin).

\textbf{Sources et fragilité.} Agrumes (oranges, citrons, mandarines), \textbf{goyave} (exceptionnellement riche, $\approx$ 200--300 mg/100 g), mangue, ananas, fruit du baobab (pain de singe), tomates, poivrons, légumes verts frais (persil, chou). La vitamine C est la plus fragile des vitamines : elle est détruite par la \textbf{chaleur} (cuisson), l'\textbf{oxydation} à l'air (contact avec le fer, découpe prolongée), la \textbf{lumière} et le \textbf{stockage prolongé} — d'où l'intérêt crucial des fruits et légumes \emph{frais}, crus ou cuits très brièvement (vapeur, sauté rapide), et de la consommation rapide après préparation.

\courssub{La carence en vitamine D : le rachitisme}

\textbf{Situation concrète.} Dans certains quartiers urbains denses ou milieux culturels couvrants, on voit des enfants de 6 mois à 3 ans dont les jambes s'arquent en « O » (genu varum) ou en « X » (genu valgum) lorsqu'ils marchent, avec un retard de dentition, un front bombé (boîte crânienne élargie), des costochondres saillants (\emph{rosaire rachitique}) et une hypotonie musculaire. Ces enfants, souvent gardés longtemps à l'intérieur, peu exposés au soleil ou couverts en permanence, souffrent de \cle{rachitisme}.

\begin{definition}[Rachitisme]
Le \cle{rachitisme} est la maladie de carence en \textbf{vitamine D} (calciférol) chez l'enfant en période de croissance osseuse active. La vitamine D (sous sa forme active calcitriol) est indispensable à l'\textbf{absorption intestinale active du calcium et du phosphore} et à leur fixation minérale sur la matrice osseuse (minéralisation du cartilage de croissance et de l'ostéoïde). En son absence, le cartilage de croissance s'épaissit mais ne se calcifie pas : les os, mal minéralisés, restent mous (ostéoïde non minéralisé) et se déforment sous le poids du corps et la traction musculaire : \textbf{jambes arquées (varum/valgum)}, déformation du thorax (bec de pigeon), retard de dentition et de la marche, retard staturo-pondéral, hypocalcémie possible (tétanie). Chez l'adulte, l'équivalent est l'\emph{ostéomalacie} (déminalisation osseuse diffuse, douleurs, fractures de fatigue).
\end{definition}

\textbf{Le rôle du soleil : une particularité métabolique.} Originalité de la vitamine D : elle est en grande partie \textbf{synthétisée \emph{in situ} par la peau} (couche basale de l'épiderme) sous l'action des \textbf{rayons ultraviolets B (UVB)} du soleil, à partir du \textbf{7-déhydrocholestérol} (dérivé du cholestérol). Sous nos latitudes tropicales, une exposition modérée et régulière (visage, bras, avant-bras, 10 à 15 minutes par jour, en évitant les heures de pic UV) suffit généralement à couvrir les besoins. Les sources alimentaires (poissons gras : sardine, maquereau, hareng ; jaune d'œuf ; foie ; produits laitiers enrichis ; champignons exposés aux UV) ne sont que complémentaires, sauf en l'absence totale d'exposition solaire.

\begin{popfigure}
\begin{center}
\begin{tikzpicture}[scale=0.95]
% Jambe normale
\begin{scope}[xshift=0cm]
% Contour jambe
\draw[fill=popcream,draw=popdark,thick] (0,0) -- (1,0) -- (0.9,3.5) -- (0.1,3.5) -- cycle;
% Os (tibia) droit
\draw[popblue,very thick] (0.5,-0.1) -- (0.5,3.6);
\draw[fill=popblueL,draw=popblue] (0.5,3.65) circle (0.15);
\node[popdark,font=\small] at (0.5,-0.6) {\textbf{Enfant sain}};
\draw[thin,popmuted,->] (0.5,1.75) -- (2.2,2.0) node[right,align=left,font=\small] {Os droit,\\ bien minéralisé};
\end{scope}
% Jambe rachitique
\begin{scope}[xshift=7.5cm]
% Contour jambe déformée
\draw[fill=popcream,draw=popdark,thick] (0.1,0) .. controls (-0.5,1.75) .. (0.1,3.5) -- (0.9,3.5) .. controls (1.5,1.75) .. (0.9,0) -- cycle;
% Os arqué
\draw[poporange,very thick] (0.5,3.6) .. controls (-0.3,1.75) .. (0.5,-0.1);
\draw[fill=poporangeL,draw=poporange] (0.5,3.65) circle (0.15);
\node[popdark,font=\small] at (0.5,-0.6) {\textbf{Enfant rachitique}};
\draw[thin,popmuted,->] (0.5,1.75) -- (-1.7,2.0) node[left,align=right,font=\small] {Os mou, arqué\\ sous le poids du corps};
\end{scope}
\end{tikzpicture}
\end{center}
\legende{Comparaison schématique d'une jambe d'enfant sain (os tibial droit, bien calcifié, résistant à la charge) et d'une jambe d'enfant rachitique (os tibial mal minéralisé, ramolli, qui s'incurve en « O » — genu varum — sous l'effet du poids du corps et de la traction musculaire), conséquence directe d'une carence en vitamine D (et/ou calcium).}
\end{popfigure}

\courssub{L'anémie ferriprive}

\textbf{Situation concrète.} Aïcha, 15 ans, élève en classe de Première, se plaint de fatigue permanente (asthénie), de maux de tête (céphalées) et d'essoufflement à la montée des escaliers (dyspnée d'effort). Sa maman remarque qu'elle est de plus en plus pâle (teint, conjonctives, paumes, lits unguéaux). Au centre de santé, l'hémogramme (numération-formule sanguine) montre un taux d'hémoglobine abaissé (< 12 g/dL chez l'adolescente), une \textbf{microcytose} (VGM bas : volume globulaire moyen diminué) et une \textbf{hypochromie} (TCHM bas : taux d'hémoglobine corpusculaire moyen diminué). L'enquête révèle une alimentation pauvre en viande et en légumes verts, des menstruations abondantes (ménorragies) et des épisodes de paludisme répétés non prévenus. Diagnostic : \textbf{anémie ferriprive microcytaire hypochrome}.

\begin{definition}[Anémie ferriprive]
L'\cle{anémie} est une diminution de la concentration en \textbf{hémoglobine} dans le sang (seuil selon âge/sexe : < 13 g/dL homme, < 12 g/dL femme/adolescente, < 11 g/dL femme enceinte/enfant 6--59 mois), pigment rouge des globules rouges (hématies/érythrocytes) qui assure le transport du dioxygène (O$_2$) des poumons vers les tissus. L'\cle{anémie ferriprive} (ou \emph{sidéroprive}), la plus fréquente dans le monde, est due à une \textbf{carence en fer (Fe)}, le fer étant l'atome central du groupe \emph{hème} de l'hémoglobine, indispensable à la fixation de l'O$_2$. Ses signes cardinaux : \textbf{pâleur} cutanéo-muqueuse, \textbf{asthénie} (fatigue anormale), \textbf{dyspnée d'effort}, vertiges, céphalées, \textbf{baisse des performances cognitives et scolaires} (troubles de l'attention, mémoire), tachycardie de compensation.
\end{definition}

\textbf{Mécanisme physiopathologique.} Sans fer disponible, la synthèse de l'hème est bloquée dans les précurseurs érythroïdes de la moelle osseuse. La moelle fabrique alors des globules rouges \textbf{plus petits (microcytose)} et \textbf{mieux pauvres en hémoglobine (hypochromie)}. Le sang transporte alors moins de dioxygène par unité de volume : pour compenser l'hypoxie tissulaire, le débit cardiaque et la ventilation augmentent — d'où la tachycardie, la dyspnée et la fatigue à l'effort. Le cerveau, grand consommateur d'O$_2$, souffre en premier (troubles cognitifs).

\textbf{Les causes multiples et souvent associées de la carence en fer :}
\begin{itemize}
\item une \textbf{alimentation pauvre en fer biodisponible} (peu de viande rouge, d'abats, de poisson, de légumineuses, de feuilles vertes) ;
\item les \textbf{pertes sanguines chroniques} : menstruations abondantes et prolongées chez l'adolescente et la femme (cause n°1 chez la femme en âge de procréer), hémorragies digestives (ulcère, hémorroïdes) ;
\item les \textbf{parasitoses hématophages} : les vers intestinaux, notamment les \textbf{ankylostomes} (\emph{Necator americanus}, \emph{Ancylostoma duodenale}) qui se fixent à la muqueuse duodénale et sucent le sang (perte de 0,1--0,3 mL/jour/ver) ;
\item le \textbf{paludisme} à répétition (plasmodiose), qui détruit massivement les globules rouges parasités (hémolyse intravasculaire) et déprime la moelle osseuse ;
\item les besoins physiologiques accrus de la \textbf{grossesse} (expansion du volume sanguin maternel, constitution des stocks fœtaux et placentaires) et de la \textbf{croissance} (adolescence, nourrisson).
\end{itemize}

\begin{attention}[Erreur fréquente au Probatoire et en pratique]
Beaucoup d'élèves (et de patients) confondent : « anémie = paludisme ». C'est \textbf{faux et dangereux}. Au Cameroun, la \textbf{carence en fer est une cause majeure, souvent prédominante, d'anémie}, fréquemment \emph{associée} au paludisme mais \textbf{indépendante} de lui sur le plan étiologique. Une personne peut être anémiée sans avoir jamais eu le paludisme, simplement parce que son alimentation manque de fer ou qu'elle perd du sang (ménorragies, ankylostomiase). Inversement, traiter correctement une crise de paludisme ne suffit \textbf{pas} à guérir une anémie ferriprive installée : il faut \textbf{impérativement} corriger l'apport en fer (alimentation + supplémentation médicamenteuse si Hb basse) et traiter la cause des pertes (déparasitage, traitement des ménorragies).
\end{attention}

\begin{exemplebox}[Des besoins en fer inégaux : la vulnérabilité de l'adolescente]
Les besoins quotidiens en fer absorbé d'une \textbf{adolescente} (environ \SI{15}{mg/jour} à \SI{18}{mg/jour} d'apport alimentaire pour couvrir les pertes et la croissance) sont presque \textbf{deux fois supérieurs} à ceux d'un garçon du même âge (environ \SI{10}{mg/jour} à \SI{12}{mg/jour}), en raison des pertes de fer liées aux \textbf{menstruations} (perte moyenne 30--40 mL sang/cycle $\approx$ 15--20 mg Fe/mois). C'est pourquoi les adolescentes et les femmes enceintes sont les premières victimes de l'anémie ferriprive, et pourquoi les supplémentations préventives en \textbf{fer et acide folique} (prévention aussi des anomalies de fermeture du tube neural) sont systématiquement proposées dès le début du 2e trimestre de la grossesse dans les centres de santé camerounais (Consultations Prénatales).
\end{exemplebox}

\textbf{Sources de fer et biodisponibilité.} Le fer \emph{héminique} (fer du groupe hème), d'origine animale, est le \textbf{mieux absorbé} (15--25 \% de l'apport) : viande rouge, abats (foie, rognons, cœur), \textbf{boudin} (très riche), poisson, sang. Le fer \emph{non héminique} (fer inorganique Fe$^{3+}$), d'origine végétale, est \textbf{mal absorbé} (1--5 \%) : haricots niébé, arachides, lentilles, soja, feuilles vertes (folon, feuilles de manioc, morelle noire, amarante), céréales complètes. \textbf{Astuces nutritionnelles majeures :} un apport de \textbf{vitamine C} au même repas (jus d'orange, de goyave, de citron, tomate crue, poivron) \textbf{multiplie par 2 à 4 l'absorption du fer végétal} (réduction Fe$^{3+}$ $\to$ Fe$^{2+}$ et chélation) ; à l'inverse, le \textbf{thé} (tanins), le \textbf{café} et le \textbf{lait} (calcium) pris \emph{au cours} du repas la diminuent fortement — les consommer \emph{entre} les repas.

\courssub{Complément utile au Probatoire : le goitre endémique (carence en iode)}

\begin{exemplebox}[Hors avitaminoses strictes (carence minérale), mais exigible au Probatoire]
Le \cle{goitre endémique} est une hypertrophie (gonflement) diffuse ou nodulaire de la \textbf{glande thyroïde} (logée à la base du cou, en avant de la trachée) due à la carence en \textbf{iode}, un oligo-élément indispensable à la synthèse des hormones thyroïdiennes (thyroxine T$_4$, triiodothyronine T$_3$). Privée d'iode, la thyroïde grossit (hyperplasie compensatrice) pour tenter de capter le peu d'iode circulant : c'est le goitre. La carence iodée sévère chez la femme enceinte (surtout au 1er trimestre) peut entraîner de graves retards irréversibles du développement neurologique et statural de l'enfant : le \textbf{crétinisme endémique} (déficience intellectuelle, surdité, paralysie spastique). Certaines régions montagneuses (ex. : contreforts de l'Adamaoua, zones volcaniques lessivées) ou éloignées de la mer sont plus touchées car leurs sols, leurs eaux et leurs aliments sont naturellement pauvres en iode. La prévention est simple, efficace, peu coûteuse et universelle : consommer exclusivement du \textbf{sel iodé} (vérifier la mention « sel iodé » ou logo sur l'emballage) et le conserver à l'abri de l'humidité et de la chaleur (cuisine), qui font sublimer l'iode (perte du fortifiant).
\end{exemplebox}

\courssub{Tableau récapitulatif des principales carences et prévention générale}

\begin{center}
\begin{tabularx}{\linewidth}{|l|X|X|}
\hline
\rowcolor{popblueL} \textbf{Élément déficitaire} & \textbf{Signes cliniques évocateurs (sémiologie)} & \textbf{Sources locales accessibles (Cameroun)} \\
\hline
Vitamine A (Rétinol / $\beta$-carotène) & Héméralopie (cécité nocturne), xérophtalmie (taches de Bitot, kératomalacie), infections respiratoires/digestives à répétition & Huile de palme rouge (crue/finale), carotte, mangue, patate douce orangée, foie, œufs, lait entier, feuilles vertes foncées (folon, manioc) \\
\hline
Vitamine C (Acide ascorbique) & Scorbut : gingivites hémorragiques, ecchymoses, mauvaise cicatrisation, asthénie, douleurs ostéo-articulaires & Goyave (champion), agrumes, ananas, mangue, fruit du baobab, tomate, poivron, légumes verts frais (crus ou cuisson vapeur brève) \\
\hline
Vitamine D (Calciférol) & Rachitisme (enfant) : jambes arquées (varum/valgum), rosaire rachitique, front bombé, retard dentition/marche, hypotonie. Ostéomalacie (adulte) & \textbf{Soleil (UVB sur peau)} : exposition modérée quotidienne. Alimentaires : poissons gras (sardine, maquereau), jaune d'œuf, foie, lait enrichi, champignons séchés au soleil \\
\hline
Fer (Fe) & Anémie microcytaire hypochrome : pâleur, asthénie, dyspnée d'effort, tachycardie, troubles cognitifs, ongles en cuillère (koïlonychie) & Viande rouge, abats (foie), boudin, poisson, sang ; niébé, lentilles, arachides, soja, feuilles vertes (folon, morelle) \textbf{+ Vitamine C au repas} \\
\hline
Iode (I) & Goitre (gonflement cou), hypothyroïdie, crétinisme endémique (si grossesse) & \textbf{Sel iodé} (obligatoire), poissons/crustacés de mer, algues \\
\hline
\end{tabularx}
\end{center}

\begin{aretenir}[Prévenir les carences nutritionnelles : trois réflexes quotidiens]
\begin{enumerate}
\item \textbf{Diversifier l'assiette chaque jour} en intégrant systématiquement des fruits et légumes locaux de saison (mangue, goyave, papaye, carotte, feuilles vertes : folon, manioc, morelle, amarante), en associant \textbf{une source de vitamine C} (fruit cru, jus frais) \textbf{aux repas riches en fer végétal} (niébé, feuilles) pour en multiplier l'absorption.
\item \textbf{Utiliser les ressources locales protectrices :} consommer de l'\textbf{huile de palme rouge} (source n°1 de provitamine A) en fin de cuisson ; utiliser exclusivement du \textbf{sel iodé} (bien refermer le sachet, loin du feu) ; exposer modérément les nourrissons et jeunes enfants au soleil (matin/tard, 10--15 min, visage/bras).
\item \textbf{Traiter et prévenir les causes infectieuses/parasitaires :} dormir sous \textbf{moustiquaire imprégnée} (MILD) pour prévenir le paludisme ; faire un \textbf{déparasitage intestinal} régulier (tous les 6 mois, albendazole/mébendazole) pour éliminer les ankylostomes ; respecter les \textbf{supplémentations officielles} (Vitamine A 6--59 mois, Fer/Folates femme enceinte) et consulter dès les premiers signes (pâleur, cécité nocturne, jambes arquées, gencives qui saignent).
\end{enumerate}
\end{aretenir}

\begin{exoresolu}[Identifier une carence et proposer une correction adaptée]
\textbf{Énoncé.} Dans un village, on observe deux enfants. L'enfant 1, âgé de 5 ans, se cogne aux objets dès la tombée de la nuit ; son alimentation est faite de mil et de riz blanc nature, sans sauce grasse ni légumes. L'enfant 2, âgé de 3 ans, a les jambes arquées en « O » (genu varum) et n'a encore que 4 dents ; il reste toute la journée à l'intérieur de la case, porté sur le dos de sa mère. Pour chaque enfant : a) identifier l'élément en carence ; b) nommer la maladie ou le signe clinique majeur ; c) proposer deux mesures de correction concrètes et adaptées au contexte.
\tcblower
\textbf{Solution détaillée.}
\emph{Enfant 1.} a) La cécité nocturne (héméralopie) évoque une carence en \textbf{vitamine A} (rétinol / bêta-carotène). b) Il s'agit de l'\textbf{héméralopie}, premier stade clinique de la carence en vitamine A, qui peut évoluer vers la \textbf{xérophtalmie} (taches de Bitot, kératomalacie, cécité). c) Correction : 1) Enrichir immédiatement l'alimentation familiale en \textbf{huile de palme rouge} (ajoutée en fin de cuisson des sauces), carottes râpées, mangues, feuilles de manioc/folon ; 2) Amener l'enfant au centre de santé pour une \textbf{supplémentation curative en vitamine A} (haute dose orale selon protocole OMS : 200 000 UI à J0, J1, J14).

\emph{Enfant 2.} a) Les jambes arquées (genu varum), le retard de dentition et l'absence d'exposition solaire évoquent une carence en \textbf{vitamine D} (et secondairement en calcium). b) Il s'agit du \textbf{rachitisme} nutritionnel. c) Correction : 1) \textbf{Exposer l'enfant au soleil} quotidiennement (visage, bras, jambes, 15--20 min vers 10h ou 16h) pour permettre la synthèse cutanée de vitamine D ; 2) Enrichir l'alimentation en sources de vitamine D et calcium (lait, poisson gras type sardine/maquereau, œuf) et, sur prescription médicale, administrer une \textbf{supplémentation en vitamine D (cholécalciférol) et calcium} (charge puis entretien).
\end{exoresolu}

\begin{exercice}[À toi de jouer — Cas clinique type Probatoire]
Une adolescente de 16 ans, scolarisée en Première, consulte pour une fatigue installée depuis 3 mois, l'empêchant de suivre en cours et de faire ses devoirs. Elle est très pâle (conjonctives, paumes), s'essouffle pour monter l'escalier de l'école (dyspnée d'effort). L'interrogatoire révèle : alimentation quasi végétarienne (riz, pâtes, niébé, peu de viande/poisson), règles abondantes (7--8 jours, nombreux caillots), 3 épisodes de paludisme traités « à la pharmacie » dans l'année, pas de déparasitage depuis 2 ans. L'hémogramme montre : Hb = 8,5 g/dL ; VGM = 68 fL (N : 80--100) ; TCHM = 28 pg (N : 27--32).
\begin{enumerate}
\item Quel type d'anémie ce bilan oriente-t-il ? Justifier par deux arguments biologiques.
\item Citer \textbf{trois} causes étiologiques qui concourent \textbf{simultanément} chez cette patiente.
\item Proposer \textbf{trois} mesures de prise en charge (préventives et curatives) prioritaires et argumentées.
\end{enumerate}
\end{exercice}

\begin{corrige}[Exercice : L'adolescente fatiguée — Anémie ferriprive multifactorielle]
\begin{enumerate}
\item Il s'agit d'une \textbf{anémie ferriprive (sidéroprive) microcytaire hypochrome}. Arguments biologiques : \textbf{microcytose} (VGM bas à 68 fL) et \textbf{hypochromie} (TCHM bas à 28 pg), signature biologique typique du défaut de synthèse de l'hème par manque de fer.
\item Trois causes concourantes : 1) \textbf{Apport insuffisant en fer héminique} (alimentation pauvre en viande/poisson/abats, fer végétal mal absorbé) ; 2) \textbf{Perte sanguine chronique gynécologique} (ménorragies abondantes = perte ferrique majeure mensuelle) ; 3) \textbf{Hémolyse et dépression médullaire} liées aux \textbf{épisodes répétés de paludisme} (destruction globules rouges) + \textbf{ankylostomiase probable} (absence déparasitage 2 ans, vers suceurs de sang).
\item Mesures prioritaires : 1) \textbf{Supplémentation médicamenteuse en fer} (sulfate/gluconate/fumarate ferreux + acide folique) per os pendant 3 à 6 mois pour reconstituer les stocks (ferritine) et normaliser l'Hb, sous contrôle biologique ; 2) \textbf{Correction diététique} : augmenter fréquence/quantité sources fer héminique (viande, foie, poisson, boudin) + associer systématiquement vitamine C (jus goyave/orange) aux repas à base de niébé/légumes ; supprimer thé/café au repas ; 3) \textbf{Traitement des causes de pertes} : déparasitage curatif (albendazole 400 mg dose unique) ; prise en charge des ménorragies (consultation gynécologique, contraception progestative si indiquée) ; prévention paludisme (MILD + chimioprévention si zone/ saison).
\end{enumerate}
\end{corrige}

En résumé, les avitaminoses (A, C, D) et l'anémie ferriprive illustrent parfaitement qu'une alimentation peut être \emph{quantitativement} suffisante (pas de faim, apport calorique correct) mais \emph{qualitativement} déficiente en micronutriments essentiels. Savoir reconnaître leurs signes cliniques d'alerte — héméralopie, gingivites hémorragiques, jambes arquées, pâleur muqueuse — et connaître les ressources alimentaires locales qui les préviennent (huile de palme rouge, goyave, soleil, niébé + citron, sel iodé) est une compétence de santé publique vitale, autant qu'une exigence incontournable du Probatoire. Dans la section suivante, nous aborderons le problème inverse, de plus en plus fréquent en milieu urbain camerounais : les maladies de surcharge nutritionnelle, et notamment l'obésité.

\coursec{Les maladies de surcharge : l'obésité}

Dans la section précédente, nous avons étudié les maladies de \cle{carence} : avitaminoses et anémie, qui résultent d'apports insuffisants en vitamines ou en fer. Mais une alimentation déséquilibrée peut aussi pécher par \emph{excès}. Dans les quartiers de Yaoundé ou de Douala, on observe de plus en plus de personnes — y compris des jeunes — qui consomment quotidiennement des boissons sucrées, des beignets, des frites et des plats très gras, tout en se déplaçant presque exclusivement en moto-taxi ou en voiture. Résultat : l'\cle{obésité}, autrefois rare dans nos pays, progresse rapidement. Cette « maladie de surcharge » est l'objet de cette section.

\courssub{Définition de l'obésité}

\textbf{Situation d'accroche.} Deux amis, tous deux âgés de 17 ans, mangent au même restaurant universitaire. Le premier marche 45 minutes par jour pour aller au lycée et joue au football le week-end ; le second prend le taxi-moto pour le moindre déplacement, grignote des chips et boit deux bouteilles de soda par jour. Au bout de quelques années, le second a pris 25 kg, essentiellement au niveau du ventre. Comment expliquer cette différence ?

\textbf{Observation.} Le corps humain fonctionne comme une \cle{balance énergétique} : d'un côté les \emph{apports} (l'énergie fournie par les aliments), de l'autre les \emph{dépenses} (l'énergie utilisée pour le fonctionnement des organes, la thermorégulation et l'activité physique). Lorsque les apports dépassent durablement les dépenses, l'énergie excédentaire n'est pas éliminée : elle est \emph{stockée} sous forme de graisses dans le tissu adipeux.

\begin{definition}[Obésité]
L'\cle{obésité} est une maladie nutritionnelle de surcharge caractérisée par un \textbf{excès de masse grasse} dans l'organisme, résultant d'apports énergétiques durablement supérieurs aux dépenses énergétiques. L'énergie non utilisée est stockée sous forme de \textbf{lipides dans le tissu adipeux}, ce qui entraîne une augmentation anormale de la masse corporelle et expose à de graves complications (diabète, hypertension, maladies cardiovasculaires).
\end{definition}

\begin{aretenir}[L'idée essentielle]
Obésité = \textbf{apports énergétiques} $>$ \textbf{dépenses énergétiques}, de façon prolongée. L'excédent est stocké dans le tissu adipeux. Ce n'est pas seulement un problème d'apparence : c'est une véritable maladie qui use le cœur, les vaisseaux, les articulations et le foie.
\end{aretenir}

\courssub{La balance énergétique : équilibre, carence, surcharge}

Pour bien comprendre, représentons les trois situations possibles de la balance énergétique.

\begin{popfigure}
\begin{center}
\begin{tikzpicture}[scale=0.9, every node/.style={font=\small}]
% --- Situation 1 : équilibre ---
\begin{scope}[xshift=0cm]
\draw[thick, popdark] (-1.6,0) -- (1.6,0);
\draw[thick, popdark] (0,0) -- (0,-0.35);
\draw[thick, popgreen] (-1.6,0) -- (-1.6,0.9);
\draw[thick, popgreen] (1.6,0) -- (1.6,0.9);
\draw[fill=popgreenL, draw=popgreen] (-2.3,0.9) rectangle (-0.9,1.5);
\node at (-1.6,1.2) {\textbf{Apports}};
\draw[fill=popblueL, draw=popblue] (0.9,0.9) rectangle (2.3,1.5);
\node at (1.6,1.2) {\textbf{Dépenses}};
\node[popgreen, font=\small\bfseries] at (0,-0.8) {Équilibre : poids stable};
\end{scope}
% --- Situation 2 : carence ---
\begin{scope}[xshift=6.2cm]
\draw[thick, popdark] (0,0.5) -- (0,0.15);
\draw[thick, popdark] (-1.6,0.5) -- (1.6,-0.3);
\draw[thick, poporange] (-1.6,0.5) -- (-1.6,1.1);
\draw[thick, popblue] (1.6,-0.3) -- (1.6,1.3);
\draw[fill=poporangeL, draw=poporange] (-2.3,1.1) rectangle (-0.9,1.7);
\node at (-1.6,1.4) {\textbf{Apports}};
\draw[fill=popblueL, draw=popblue] (0.9,1.3) rectangle (2.3,1.9);
\node at (1.6,1.6) {\textbf{Dépenses}};
\node[poporange, font=\small\bfseries] at (0,-0.8) {Carence : amaigrissement};
\end{scope}
% --- Situation 3 : surcharge ---
\begin{scope}[xshift=12.4cm]
\draw[thick, popdark] (0,0.5) -- (0,0.15);
\draw[thick, popdark] (-1.6,-0.3) -- (1.6,0.5);
\draw[thick, poppink] (-1.6,-0.3) -- (-1.6,1.3);
\draw[thick, popblue] (1.6,0.5) -- (1.6,1.1);
\draw[fill=poppinkL, draw=poppink] (-2.3,1.3) rectangle (-0.9,1.9);
\node at (-1.6,1.6) {\textbf{Apports}};
\draw[fill=popblueL, draw=popblue] (0.9,1.1) rectangle (2.3,1.7);
\node at (1.6,1.4) {\textbf{Dépenses}};
\node[poppink, font=\small\bfseries] at (0,-0.8) {Surcharge : prise de poids};
\end{scope}
\end{tikzpicture}
\end{center}
\legende{La balance énergétique. À gauche : apports égaux aux dépenses, le poids reste stable. Au centre : apports inférieurs aux dépenses, l'organisme puise dans ses réserves (carence, amaigrissement). À droite : apports supérieurs aux dépenses, l'excédent est stocké dans le tissu adipeux (surcharge, obésité).}
\end{popfigure}

\courssub{Causes de l'obésité}

D'où vient ce déséquilibre ? Menons la démarche d'investigation.

\textbf{Observation.} Des enquêtes de santé publique menées dans les grandes villes africaines montrent que la fréquence de l'obésité augmente parallèlement à l'urbanisation : consommation croissante de plats préparés très gras et très sucrés, de boissons sucrées, et diminution de la marche et des activités physiques.

\textbf{Hypothèse.} L'obésité résulte d'une alimentation trop riche en énergie combinée à une dépense physique trop faible.

\textbf{Données.} Une bouteille de soda de \SI{50}{cL} contient environ 50 g de sucre, soit environ 200 kcal ; il faut marcher près de 45 minutes pour dépenser cette énergie. Un élève qui boit deux sodas par jour sans activité physique accumule donc un excédent quotidien considérable.

\textbf{Interprétation et conclusion.} L'hypothèse est confirmée : le déséquilibre de la balance énergétique est la cause directe de l'obésité. On distingue plusieurs facteurs favorisants :

\begin{itemize}
\item une \textbf{alimentation trop riche en graisses et en sucres} : fritures, plats très huilés, pâtisseries, bonbons ;
\item la consommation régulière de \textbf{boissons sucrées} (sodas, jus industriels sucrés), qui apportent du sucre sans rassasier ;
\item la \textbf{sédentarité} : absence de marche, de sport, temps excessif devant les écrans ;
\item des \textbf{facteurs génétiques et hormonaux} (admis) : certaines personnes stockent plus facilement les graisses, et certains troubles hormonaux (par exemple de la glande thyroïde) favorisent la prise de poids. Ces facteurs n'expliquent qu'une minorité des cas : dans la grande majorité des situations, l'obésité est liée au mode de vie.
\end{itemize}

\begin{attention}[Ne pas tout mettre sur le compte de l'hérédité]
On entend souvent : « dans ma famille, on est gros de naissance ». Les facteurs génétiques existent, mais ils ne déterminent pas à eux seuls l'obésité. Une alimentation équilibrée et une activité physique régulière permettent à la plupart des personnes, même prédisposées, de maintenir un poids normal. À l'examen, cite toujours d'abord les causes liées au mode de vie (alimentation, sédentarité) avant les facteurs génétiques et hormonaux.
\end{attention}

\courssub{Conséquences de l'obésité sur la santé}

L'excès de masse grasse n'est pas inoffensif : il perturbe le fonctionnement de nombreux organes.

\begin{itemize}
\item \textbf{Diabète de type 2} : l'excès de graisse rend les cellules moins sensibles à l'insuline ; la glycémie (taux de sucre dans le sang) reste anormalement élevée, ce qui endommage à terme les reins, les yeux et les nerfs.
\item \textbf{Hypertension artérielle} : le cœur doit fournir un effort supplémentaire pour irriguer une masse corporelle plus importante ; la pression dans les artères augmente durablement.
\item \textbf{Maladies cardiovasculaires} : les graisses en excès se déposent sur les parois des artères (athérosclérose), ce qui favorise l'infarctus du myocarde et les accidents vasculaires cérébraux (AVC).
\item \textbf{Douleurs articulaires} : les genoux et les hanches supportent une charge excessive, ce qui use prématurément les cartilages (arthrose).
\item \textbf{Difficultés respiratoires} : la graisse abdominale gêne les mouvements du diaphragme ; on observe des essoufflements à l'effort et des arrêts de la respiration pendant le sommeil (apnées du sommeil).
\item \textbf{Retentissement psychologique} : moqueries, rejet, perte d'estime de soi, isolement, parfois dépression — en particulier chez les adolescents.
\end{itemize}

\begin{savaistu}[L'obésité, un problème croissant au Cameroun]
Selon l'Organisation mondiale de la santé (OMS), l'obésité progresse rapidement dans les pays en développement, où elle coexiste désormais avec la malnutrition par carence : on parle de « double fardeau » de la malnutrition. Au Cameroun, les campagnes de santé publique encouragent la consommation de fruits et légumes locaux (ndolé, gombo, fruits de saison) et la pratique régulière d'une activité physique, y compris la simple marche quotidienne.
\end{savaistu}

\courssub{L'indice de masse corporelle (IMC) : mesurer la surcharge}

Comment savoir objectivement si une personne est en surpoids ou obèse ? Le poids seul ne suffit pas : 80 kg est normal pour un homme de 1,90 m mais excessif pour une personne de 1,55 m. Il faut donc rapporter la masse à la taille. C'est le rôle de l'\cle{IMC}.

\begin{definition}[Indice de masse corporelle]
L'\cle{indice de masse corporelle} (IMC) est un indicateur qui rapporte la masse d'une personne au carré de sa taille :
\[ \mathrm{IMC} = \frac{\text{masse (en kg)}}{\text{taille}^2\ (\text{en m}^2)} \]
Il s'exprime en kg/m\textsuperscript{2} (souvent écrit sans unité).
\end{definition}

\begin{propriete}[Seuils d'interprétation de l'IMC chez l'adulte (normes OMS, admises)]
Chez l'adulte, l'OMS définit les catégories suivantes :
\begin{itemize}
\item IMC $< 18{,}5$ : \textbf{maigreur} ;
\item $18{,}5 \leq \mathrm{IMC} \leq 24{,}9$ : corpulence \textbf{normale} ;
\item $25 \leq \mathrm{IMC} \leq 29{,}9$ : \textbf{surpoids} ;
\item $\mathrm{IMC} \geq 30$ : \textbf{obésité}.
\end{itemize}
Ces normes sont admises : on ne demande pas de les démontrer, mais il faut savoir les utiliser pour interpréter un cas concret. Chez l'enfant et l'adolescent, on utilise plutôt les courbes de croissance (section suivante).
\end{propriete}

\begin{popfigure}
\begin{center}
\begin{tikzpicture}[scale=0.85, every node/.style={font=\small}]
% Axe principal
\draw[thick, ->] (0,0) -- (13,0) node[right] {IMC (kg/m\textsuperscript{2})};
% Graduations et valeurs seuils
\foreach \x/\label in {0/15, 1.75/18,5, 5.25/25, 7.5/30, 12.5/40} {
    \draw (\x, -0.1) -- (\x, 0.1) node[below=2pt] {\label};
}
% Zones colorées (rectangles sous l'axe)
\fill[poporangeL, draw=poporange] (0, -0.5) rectangle (1.75, -1.5) node[midway, text=poporange, font=\bfseries] {Maigreur};
\fill[popgreenL, draw=popgreen] (1.75, -0.5) rectangle (5.25, -1.5) node[midway, text=popgreen, font=\bfseries] {Normal};
\fill[popgoldL, draw=popgold] (5.25, -0.5) rectangle (7.5, -1.5) node[midway, text=popgold, font=\bfseries] {Surpoids};
\fill[poppinkL, draw=poppink] (7.5, -0.5) rectangle (12.5, -1.5) node[midway, text=poppink, font=\bfseries] {Obésité};
% Lignes de séparation pointillées montant vers l'axe
\draw[dashed, thick, popdark] (1.75, -0.5) -- (1.75, 0.5);
\draw[dashed, thick, popdark] (5.25, -0.5) -- (5.25, 0.5);
\draw[dashed, thick, popdark] (7.5, -0.5) -- (7.5, 0.5);
% Flèches indicatives au-dessus de l'axe
\draw[<->, thick, popgreen] (1.75, 0.7) -- (5.25, 0.7) node[midway, above=2pt, font=\bfseries] {Zone de corpulence normale};
\end{tikzpicture}
\end{center}
\legende{Échelle des valeurs de l'IMC chez l'adulte (normes OMS). Les seuils 18,5 ; 25 et 30 délimitent quatre zones : maigreur (< 18,5), normal (18,5–24,9), surpoids (25–29,9), obésité (≥ 30). La flèche verte souligne la zone cible pour la santé.}
\end{popfigure}

\begin{exemplebox}[Calcul de l'IMC d'un élève]
Un élève de Première pèse 70 kg et mesure 1,65 m. Calculons son IMC :
\begin{align*}
\mathrm{IMC} &= \frac{70}{1{,}65^2} = \frac{70}{2{,}7225} \approx 25{,}7
\end{align*}
Or $25 \leq 25{,}7 \leq 29{,}9$ : cet élève est en \textbf{surpoids}. Il n'est pas encore obèse, mais il doit réagir : réduire les boissons sucrées et les grignotages, et pratiquer une activité physique régulière pour revenir dans la zone de corpulence normale.
\end{exemplebox}

\begin{methode}[Calculer et interpréter un IMC en 3 étapes]
\begin{enumerate}
\item \textbf{Calculer le carré de la taille} : convertir d'abord la taille en \emph{mètres}, puis élever au carré. Exemple : 1,65 m $\rightarrow$ $1{,}65^2 = 2{,}7225$ m\textsuperscript{2}.
\item \textbf{Diviser la masse par ce carré} : $\mathrm{IMC} = \dfrac{\text{masse (kg)}}{\text{taille}^2\ (\text{m}^2)}$. Exemple : $70 / 2{,}7225 \approx 25{,}7$.
\item \textbf{Comparer aux seuils OMS} : $<18{,}5$ maigreur ; $18{,}5$–$24{,}9$ normal ; $25$–$29{,}9$ surpoids ; $\geq 30$ obésité. Conclure par une phrase complète.
\end{enumerate}
\end{methode}

\begin{attention}[Les deux erreurs classiques au Probatoire]
\begin{itemize}
\item \textbf{Oublier d'élever la taille au carré} : écrire $\mathrm{IMC} = 70 / 1{,}65 \approx 42{,}4$ est faux. La formule contient la taille \emph{au carré}.
\item \textbf{Utiliser la taille en centimètres sans conversion} : écrire $\mathrm{IMC} = 70 / 165^2$ donne un résultat absurde ($\approx 0{,}0026$). Il faut d'abord convertir : $165\ \text{cm} = 1{,}65\ \text{m}$. Vérifie toujours l'ordre de grandeur : un IMC se situe normalement entre 15 et 45.
\end{itemize}
\end{attention}

\courssub{Prévention et correction de l'obésité}

L'obésité étant essentiellement une maladie du mode de vie, elle se prévient et se corrige en rééquilibrant la balance énergétique, des deux côtés à la fois : diminuer les apports excessifs et augmenter les dépenses.

\begin{itemize}
\item \textbf{Réduire les sucres et les graisses} : limiter les sodas et boissons sucrées (les remplacer par de l'eau), les fritures, les plats trop huilés, les bonbons et pâtisseries.
\item \textbf{Augmenter la consommation de fruits et légumes} : peu caloriques, riches en fibres, vitamines et sels minéraux, ils rassasient sans surcharger. Les marchés camerounais offrent une grande variété de légumes (ndolé, gombo, feuilles de manioc) et de fruits (mangues, papayes, avocats avec modération).
\item \textbf{Pratiquer une activité physique régulière} : au moins \textbf{30 minutes par jour} — marche rapide, footing, football, danse, travaux champêtres. Aller à l'école à pied quand c'est possible est déjà un excellent exercice.
\item \textbf{Éviter les grignotages} entre les repas : chips, beignets et bonbons pris à tout moment de la journée apportent une énergie considérable sans valeur nutritive.
\end{itemize}

\begin{exoresolu}[Calculer et interpréter un IMC]
\textbf{Énoncé.} Mme Ngo Bell pèse 84 kg et mesure 1,60 m. 1) Calcule son IMC (arrondi au dixième). 2) Interprète le résultat à l'aide des normes OMS. 3) Propose deux mesures de correction adaptées.
\tcblower
\textbf{Solution détaillée.}

1) On convertit d'abord la taille en mètres (déjà en mètres ici) et on calcule son carré : $1{,}60^2 = 2{,}56$ m\textsuperscript{2}. Puis :
\[ \mathrm{IMC} = \frac{84}{2{,}56} \approx 32{,}8\ \text{kg/m}^2. \]

2) Comparaison aux seuils OMS : $32{,}8 \geq 30$, donc Mme Ngo Bell est en situation d'\textbf{obésité}. Elle s'expose au diabète de type 2, à l'hypertension artérielle et aux maladies cardiovasculaires.

3) Mesures de correction : réduire les graisses et les sucres (supprimer les sodas, limiter les fritures) et augmenter les fruits et légumes ; pratiquer au moins 30 minutes d'activité physique par jour (marche rapide quotidienne) et éviter les grignotages entre les repas. Un suivi médical au centre de santé est recommandé.
\end{exoresolu}

\begin{exercice}[À toi de jouer]
Un lycéen mesure 1,75 m et pèse 90 kg. 1) Calcule son IMC. 2) Dans quelle catégorie OMS se situe-t-il ? 3) Quelle masse maximale devrait-il atteindre pour revenir à la limite supérieure de la corpulence normale (IMC = 24,9) ?
\end{exercice}

\begin{corrige}[Exercice 1]
1) $1{,}75^2 = 3{,}0625$ ; $\mathrm{IMC} = 90 / 3{,}0625 \approx 29{,}4$.

2) $25 \leq 29{,}4 \leq 29{,}9$ : le lycéen est en \textbf{surpoids} (limite supérieure de la catégorie).

3) Masse maximale : $m = 24{,}9 \times 3{,}0625 \approx 76{,}3$ kg. Il devrait donc perdre environ 14 kg, progressivement, par une alimentation allégée en sucres et graisses et une activité physique régulière.
\end{corrige}

\begin{aretenir}[L'essentiel de la section]
\begin{itemize}
\item L'obésité est un \textbf{excès de masse grasse} dû à des apports énergétiques supérieurs aux dépenses.
\item Causes principales : alimentation riche en graisses et sucres, boissons sucrées, sédentarité ; facteurs génétiques et hormonaux (admis).
\item Conséquences : diabète de type 2, hypertension, maladies cardiovasculaires, douleurs articulaires, difficultés respiratoires, retentissement psychologique.
\item $\mathrm{IMC} = \dfrac{\text{masse (kg)}}{\text{taille}^2\ (\text{m}^2)}$ ; seuils OMS : $<18{,}5$ maigreur ; $18{,}5$–$24{,}9$ normal ; $25$–$29{,}9$ surpoids ; $\geq 30$ obésité.
\item Prévention : moins de sucres et de graisses, plus de fruits et légumes, 30 min d'activité physique par jour, pas de grignotage.
\end{itemize}
\end{aretenir}

\coursec{Interprétation des données : courbe de croissance et IMC}

Dans la section précédente, nous avons étudié l'obésité, maladie de surcharge pondérale dont le diagnostic repose sur un outil simple : l'indice de masse corporelle (IMC). Mais comment dépister \emph{précocement} une malnutrition chez un nourrisson, avant même l'apparition des signes visibles comme l'œdème du kwashiorkor ou l'amaigrissement du marasme ? Comment savoir si un adulte est réellement en surcharge pondérale ? Pour répondre à ces questions, les agents de santé disposent de deux outils simples, peu coûteux et utilisés partout au Cameroun, du centre de santé intégré de quartier à l'hôpital de district : la \cle{courbe de croissance} pour l'enfant et l'\cle{IMC} pour l'adulte. Cette section t'apprend à lire ces outils et à interpréter des données simples, une compétence exigible au Probatoire.

\courssub{La courbe de croissance : un outil de suivi du poids de l'enfant}

\textbf{Une situation concrète pour commencer.} Au centre de santé de Bafoussam, une mère présente le carnet de santé de sa fille de 14 mois. L'infirmier pèse l'enfant : 7,2 kg. Il reporte un point sur une feuille quadrillée du carnet, relie ce point aux précédents, puis fronce les sourcils : « Madame, le poids de votre fille n'a pas bougé depuis deux mois. Il faut revoir son alimentation. » Qu'a donc vu l'infirmier sur ce graphique ? C'est ce que nous allons comprendre.

\begin{definition}[Courbe de croissance]
La \cle{courbe de croissance} est un graphique, imprimé dans le \cle{carnet de santé} de chaque enfant, qui représente l'évolution du \textbf{poids de l'enfant en fonction de son âge}, de la naissance à 5 ans (voire plus). Elle comporte des \textbf{courbes de référence} (établies par l'Organisation mondiale de la Santé, OMS) qui délimitent une \textbf{zone normale}, une \textbf{zone de risque} et une \textbf{zone de malnutrition}. Elle permet de suivre la croissance de l'enfant et de dépister précocement une malnutrition.
\end{definition}

\textbf{Comment est construite cette courbe ?} Le graphique comporte deux axes :
\begin{itemize}
\item l'\textbf{axe des abscisses} (horizontal) porte l'\textbf{âge} de l'enfant, exprimé en \textbf{mois} ;
\item l'\textbf{axe des ordonnées} (vertical) porte le \textbf{poids} de l'enfant, exprimé en \textbf{kilogrammes}.
\end{itemize}
À chaque visite de pesée (idéalement mensuelle pendant la première année), l'agent de santé reporte un point : l'âge en abscisse, le poids en ordonnée. Les points successifs, reliés entre eux, forment la courbe de croissance \emph{individuelle} de l'enfant, que l'on compare aux courbes de référence.

\textbf{Les courbes de référence et les zones.} Les courbes de référence OMS sont obtenues à partir des mesures de milliers d'enfants en bonne santé, bien nourris, dans plusieurs pays du monde. Sur le carnet de santé, on distingue généralement :
\begin{itemize}
\item une \textbf{courbe supérieure} (souvent le z-score +2 ou +3) : au-dessus d'elle, l'enfant est en surcharge pondérale (risque d'obésité infantile) ;
\item une \textbf{courbe inférieure} (souvent le z-score -2) : en dessous d'elle, l'enfant est en \textbf{malnutrition} (poids insuffisant pour son âge) ;
\item l'espace compris \textbf{entre les deux courbes} constitue la \textbf{zone normale} ; la partie basse de cette zone, proche de la courbe inférieure, est une \textbf{zone de risque} qui impose une surveillance rapprochée.
\end{itemize}

\begin{exemplebox}[Un repère chiffré à retenir]
À la naissance, un nouveau-né camerounais en bonne santé pèse en moyenne environ 3 kg. Sa croissance pondérale est très rapide la première année :
\begin{itemize}
\item vers \textbf{4 à 6 mois}, le poids de naissance a \textbf{doublé} (environ 6 kg) ;
\item à \textbf{12 mois}, il a \textbf{triplé} : un enfant en bonne santé pèse donc environ \textbf{9 kg} à un an pour 3 kg à la naissance ;
\item à 24 mois, il a quadruplé (environ 12 kg).
\end{itemize}
Ces repères simples permettent déjà de détecter un retard de croissance évident.
\end{exemplebox}

\begin{popfigure}
\begin{center}
\begin{tikzpicture}
\begin{axis}[
width=0.95\linewidth, height=7.5cm,
xlabel={Âge (mois)}, ylabel={Poids (kg)},
xmin=0, xmax=36, ymin=2, ymax=16,
xtick={0,6,12,18,24,30,36},
ytick={2,4,6,8,10,12,14,16},
grid=both, grid style={popline!40},
legend style={at={(0.02,0.98)}, anchor=north west, font=\small, draw=popline}
]
% Courbe supérieure de référence
\addplot[popblue, thick, domain=0:36, samples=50] {3.2 + 3.4*sqrt(max(0,x/12))};
\addlegendentry{Courbe supérieure de référence}
% Courbe inférieure de référence
\addplot[poporange, thick, domain=0:36, samples=50] {2.4 + 2.4*sqrt(max(0,x/12))};
\addlegendentry{Courbe inférieure de référence}
% Enfant en bonne santé
\addplot[popgreen, very thick, mark=*, mark size=1.5pt] coordinates {
(0,3.0) (3,5.2) (6,6.9) (9,8.2) (12,9.2) (15,10.0) (18,10.8) (21,11.5) (24,12.1) (27,12.8) (30,13.4) (36,14.4)
};
\addlegendentry{Enfant bien nourri}
% Enfant en malnutrition
\addplot[poppink, very thick, mark=square*, mark size=1.5pt] coordinates {
(0,2.9) (3,4.6) (6,5.6) (9,6.2) (12,6.5) (15,6.5) (18,6.3) (21,6.4) (24,6.2)
};
\addlegendentry{Enfant en malnutrition}
\end{axis}
\end{tikzpicture}
\end{center}
\legende{Courbe de croissance poids/âge (0 à 36 mois). Les courbes de référence délimitent la zone normale. La courbe verte (enfant bien nourri) progresse régulièrement entre les deux références ; la courbe rose (enfant en malnutrition) marque un palier vers 12--15 mois, chute, puis passe sous la courbe inférieure : la malnutrition est installée.}
\end{popfigure}

\courssub{Lecture et interprétation de la courbe de croissance}

\begin{methode}[Lire une courbe de croissance en 4 étapes]
\begin{enumerate}
\item \textbf{Repérer l'âge} : sur l'axe horizontal, localiser l'âge de l'enfant au moment de la pesée (en mois) et monter verticalement.
\item \textbf{Repérer le poids} : sur l'axe vertical, localiser le poids mesuré (en kg) et se déplacer horizontalement.
\item \textbf{Situer le point} : le point de croisement des deux repères représente l'état pondéral de l'enfant à cette date ; le reporter et le relier au point précédent.
\item \textbf{Comparer aux courbes de référence} : déterminer dans quelle zone se trouve le point (au-dessus de la courbe supérieure, entre les deux courbes, ou sous la courbe inférieure) et observer la \textbf{pente} du tracé : monte-t-il régulièrement, fait-il un palier, descend-il ?
\end{enumerate}
\end{methode}

\textbf{Les signes d'alerte sur la courbe.} L'essentiel n'est pas seulement la position du point, mais la \emph{direction} de la courbe. Deux situations doivent alerter les parents et l'agent de santé :
\begin{itemize}
\item le \textbf{palier} : le poids \textbf{stagne} entre deux mesures consécutives (la courbe devient horizontale) ; à cet âge, un enfant bien nourri grossit chaque mois, donc une stagnation signifie que l'apport alimentaire ne couvre plus les besoins de croissance ;
\item la \textbf{chute} : le poids \textbf{diminue} entre deux mesures consécutives (la courbe descend) ; c'est le signe d'une malnutrition aiguë ou d'une maladie en cours (diarrhées répétées, paludisme, rougeole...).
\end{itemize}
Si la courbe franchit la courbe inférieure de référence, le diagnostic de \textbf{malnutrition} est posé et une prise en charge nutritionnelle s'impose.

\begin{attention}[Un palier est déjà un signal d'alerte]
Erreur fréquente, chez les parents comme chez les candidats au Probatoire : croire que seul le passage \emph{sous} la courbe inférieure est inquiétant. \textbf{Faux !} Un \textbf{palier de poids} (stagnation sur deux mesures consécutives) est déjà un signal d'alerte, \emph{même si l'enfant se trouve encore dans la zone normale}. C'est précisément l'intérêt de la courbe de croissance : elle permet d'agir \textbf{avant} que la malnutrition ne s'installe, en corrigeant l'alimentation ou en traitant une maladie intercurrente. Attendre que la courbe sorte de la zone normale, c'est attendre que l'enfant soit déjà malade.
\end{attention}

\begin{popfigure}
\begin{center}
\begin{tikzpicture}
\begin{axis}[
width=0.95\linewidth, height=7cm,
xlabel={Âge (mois)}, ylabel={Poids (kg)},
xmin=0, xmax=25, ymin=2, ymax=11,
xtick={0,3,6,9,12,15,18,21,24},
ytick={2,3,4,5,6,7,8,9,10,11},
grid=both, grid style={popline!40},
legend style={at={(0.02,0.98)}, anchor=north west, font=\small, draw=popline}
]
\addplot[poporange, thick, domain=0:25, samples=40] {2.4 + 2.4*sqrt(max(0,x/12))};
\addlegendentry{Courbe inférieure de référence}
\addplot[poppurple, very thick, mark=*, mark size=2pt] coordinates {
(0,3.1) (2,4.4) (4,5.6) (6,6.5) (8,7.3) (10,7.9) (12,8.3) (14,8.3) (16,8.3) (18,8.0) (20,7.6) (22,7.3)
};
\addlegendentry{Poids mensuel de l'enfant}
% annotations
\draw[popdark, thick, ->] (axis cs:15,8.55) -- (axis cs:15,9.3) node[above, font=\small, align=center]{Palier (12--16 mois)};
\draw[popdark, thick, ->] (axis cs:20,7.35) -- (axis cs:20,6.4) node[below, font=\small, align=center]{Chute (18--22 mois)};
\end{axis}
\end{tikzpicture}
\end{center}
\legende{Nuage de points des pesées mensuelles d'un enfant. La croissance est normale jusqu'à 12 mois, puis le poids stagne (palier de 12 à 16 mois : premier signal d'alerte) avant de diminuer (chute à partir de 18 mois). Vers 20--22 mois, la courbe passe sous la courbe inférieure de référence : la malnutrition est installée. Une correction alimentaire aurait dû intervenir dès le palier.}
\end{popfigure}

\courssub{Interprétation guidée d'un cas}

\begin{exoresolu}[Le cas de la petite Aïcha, 12 mois]
\textbf{Énoncé.} Aïcha, née avec un poids de 3,1 kg, est pesée chaque mois au centre de santé. Voici ses pesées récentes : à 9 mois, 7,0 kg ; à 10 mois, 7,0 kg ; à 11 mois, 6,8 kg ; à 12 mois, 6,7 kg. Sur le carnet de santé, la courbe inférieure de référence passe à environ 7,2 kg à 12 mois.
\begin{enumerate}
\item Quel poids aurait dû avoir Aïcha à 12 mois si sa croissance était normale ?
\item Décris l'allure de sa courbe de croissance entre 9 et 12 mois.
\item Conclus sur l'état nutritionnel d'Aïcha et propose une conduite à tenir.
\end{enumerate}
\tcblower
\textbf{Solution détaillée.}
\begin{enumerate}
\item À 12 mois, un enfant en bonne santé pèse environ le \textbf{triple de son poids de naissance} : $3{,}1 \times 3 \approx 9{,}3$ kg. Aïcha, avec 6,7 kg, est très en dessous de ce repère.
\item Entre 9 et 10 mois, le poids \textbf{stagne} (7,0 kg puis 7,0 kg) : c'est un \textbf{palier}, premier signal d'alerte. Entre 10 et 12 mois, le poids \textbf{diminue} (7,0 kg $\rightarrow$ 6,8 kg $\rightarrow$ 6,7 kg) : c'est une \textbf{chute} sur deux mesures consécutives, signe de malnutrition évolutive.
\item À 12 mois, le poids d'Aïcha (6,7 kg) est \textbf{sous la courbe inférieure de référence} (7,2 kg) : Aïcha est en \textbf{malnutrition}. Conduite à tenir : consulter sans tarder le personnel de santé pour rechercher une cause (infection, diarrhées, sevrage mal conduit, diversification alimentaire insuffisante), enrichir l'alimentation (bouillies énergétiques enrichies en huile, arachide, soja, œuf, poisson), poursuivre l'allaitement si possible et assurer un suivi rapproché des pesées.
\end{enumerate}
\end{exoresolu}

\courssub{L'IMC et l'utilisation conjointe des deux outils}

\textbf{Rappel sur l'IMC.} Chez l'adulte et l'adolescent, la courbe de croissance n'est plus utilisable ; on emploie l'\cle{indice de masse corporelle} (IMC), déjà rencontré à propos de l'obésité :
\[ \mathrm{IMC} = \frac{\text{masse (en kg)}}{\text{taille}^2\ (\text{en m}^2)} \]
L'interprétation chez l'adulte est la suivante :

\begin{center}
\begin{tabularx}{\linewidth}{|l|X|}\hline
\rowcolor{popblueL} \textbf{Valeur de l'IMC (kg/m$^2$)} & \textbf{Interprétation nutritionnelle} \\ \hline
IMC $< 18{,}5$ & Maigreur (dénutrition de l'adulte) \\ \hline
$18{,}5 \leq \mathrm{IMC} < 25$ & Corpulence normale \\ \hline
$25 \leq \mathrm{IMC} < 30$ & Surpoids \\ \hline
IMC $\geq 30$ & Obésité \\ \hline
\end{tabularx}
\end{center}

\textbf{Utilisation conjointe pour un diagnostic nutritionnel simple.} Les deux outils sont \emph{complémentaires} et couvrent toute la vie :
\begin{itemize}
\item de la \textbf{naissance à 5 ans}, la \textbf{courbe de croissance} du carnet de santé est l'outil de référence : elle suit le poids mois après mois et détecte paliers et chutes ;
\item chez l'\textbf{adulte}, l'\textbf{IMC} permet de classer l'état nutritionnel : maigreur, corpulence normale, surpoids ou obésité.
\item chez l'\textbf{adolescent}, l'IMC s'interprète à l'aide de \textbf{courbes de référence âge/sexe} (percentiles OMS), et non avec les seuils fixes de l'adulte.
\end{itemize}
Ainsi, un même raisonnement s'applique dans les deux cas : on \textbf{mesure} (poids, taille, âge), on \textbf{reporte ou calcule}, on \textbf{compare à une référence}, puis on \textbf{conclut} et on propose une mesure de prévention ou de correction adaptée.

\begin{exemplebox}[Deux diagnostics rapides]
\begin{itemize}
\item \textbf{Enfant} : Junior, 18 mois, pèse 8,5 kg alors que la courbe inférieure de référence passe à 9 kg à cet âge, et son poids n'a pas augmenté depuis 3 mois. Conclusion : palier prolongé et passage sous la courbe inférieure $\Rightarrow$ malnutrition ; consultation et correction alimentaire urgentes.
\item \textbf{Adulte} : M. Etoundi mesure 1,70 m et pèse 92 kg. $\mathrm{IMC} = \dfrac{92}{1{,}70^2} = \dfrac{92}{2{,}89} \approx 31{,}8$ kg/m$^2$. Conclusion : IMC $\geq 30$ $\Rightarrow$ obésité ; régime hypocalorique équilibré et activité physique régulière sont recommandés.
\end{itemize}
\end{exemplebox}

\begin{attention}[Ces outils orientent, ils ne remplacent pas le médecin]
La courbe de croissance et l'IMC sont des outils de \textbf{dépistage et d'orientation} : ils signalent un risque ou une anomalie, mais ils ne suffisent pas à poser un diagnostic médical complet. Un IMC élevé peut par exemple correspondre à une forte masse musculaire chez un sportif, et une chute de poids chez l'enfant peut révéler une maladie (tuberculose, paludisme, VIH/Sida) que seul un examen médical identifiera. Toute anomalie détectée doit conduire à \textbf{consulter un personnel de santé qualifié}, qui seul peut confirmer le diagnostic et prescrire la prise en charge adaptée.
\end{attention}

\begin{aretenir}[L'essentiel de la section]
\begin{itemize}
\item La \textbf{courbe de croissance} (carnet de santé) suit le \textbf{poids de l'enfant en fonction de l'âge} : âge en abscisses (mois), poids en ordonnées (kg).
\item Les \textbf{courbes de référence OMS} délimitent une zone normale, une zone de risque et une zone de malnutrition.
\item Signes d'alerte : un \textbf{palier} (stagnation) ou une \textbf{chute} (diminution) du poids sur deux mesures consécutives, \emph{même dans la zone normale}.
\item Un enfant dont la courbe passe \textbf{sous la courbe inférieure} est en \textbf{malnutrition}.
\item Repère : à 12 mois, le poids de naissance a \textbf{triplé} (environ 9 kg pour 3 kg).
\item Chez l'adulte, l'\textbf{IMC} $= \dfrac{\text{masse}}{\text{taille}^2}$ classe l'état nutritionnel (maigreur $< 18{,}5$ ; normal $18{,}5$--$25$ ; surpoids $25$--$30$ ; obésité $\geq 30$).
\item Chez l'adolescent, l'IMC s'interprète avec des courbes de référence âge/sexe (percentiles).
\item Ces outils \textbf{orientent} mais ne remplacent pas l'avis médical.
\end{itemize}
\end{aretenir}

\begin{exercice}[À toi de jouer]
Voici les pesées d'un nourrisson : naissance 3,2 kg ; 3 mois : 5,5 kg ; 6 mois : 7,0 kg ; 9 mois : 7,6 kg ; 12 mois : 7,6 kg ; 15 mois : 7,4 kg. La courbe inférieure de référence passe à 7,2 kg à 12 mois et 7,6 kg à 15 mois.
\begin{enumerate}
\item Trace (sur ton cahier) la courbe de croissance de cet enfant.
\item Identifie le ou les signes d'alerte et précise à quel âge ils apparaissent.
\item L'enfant est-il en malnutrition à 15 mois ? Justifie.
\item Calcule l'IMC de sa mère (1,62 m ; 78 kg) et conclus sur son état nutritionnel.
\end{enumerate}
\end{exercice}

\begin{corrige}[Exercice : courbe de croissance et IMC]
\begin{enumerate}
\item On reporte les points (0 ; 3,2), (3 ; 5,5), (6 ; 7,0), (9 ; 7,6), (12 ; 7,6), (15 ; 7,4) et on les relie.
\item Entre 9 et 12 mois, le poids stagne (7,6 kg puis 7,6 kg) : \textbf{palier} à partir de 12 mois. Entre 12 et 15 mois, le poids diminue (7,6 kg $\rightarrow$ 7,4 kg) : \textbf{chute}.
\item À 15 mois, le poids (7,4 kg) est \textbf{sous la courbe inférieure} (7,6 kg) : l'enfant est en \textbf{malnutrition}. L'alerte aurait dû intervenir dès le palier à 12 mois.
\item $\mathrm{IMC} = \dfrac{78}{1{,}62^2} = \dfrac{78}{2{,}62} \approx 29{,}7$ kg/m$^2$. La mère est en \textbf{surpoids} (proche de l'obésité) : rééquilibrage alimentaire et activité physique conseillés, avec avis médical.
\end{enumerate}
\end{corrige}

\coursec{Principes d'une alimentation préventive}

Dans la section précédente, nous avons appris à \cle{interpréter} une courbe de croissance et à calculer l'\cle{IMC} pour repérer une malnutrition ou une surcharge pondérale. Mais diagnostiquer ne suffit pas : il faut \textbf{prévenir}. Or la grande leçon de toutes les études nutritionnelles menées au Cameroun comme ailleurs est simple : la plupart des maladies nutritionnelles — kwashiorkor, marasme, avitaminoses, anémie, obésité — pourraient être évitées par une alimentation correctement conduite, sans médicament coûteux. Cette dernière section dégage les \textbf{principes d'une alimentation préventive}, applicables au quotidien, dans nos familles et nos cantines scolaires.

\courssub{Principe 1 : la diversité alimentaire — les trois groupes d'aliments à chaque repas}

\textbf{Situation concrète.} Dans un quartier de Douala, deux familles voisines ont des revenus comparables. Dans la première, le repas quotidien est presque toujours le même : un grand plat de riz blanc avec un peu d'huile. Dans la seconde, on sert du couscous de maïs avec une sauce de feuilles de folong, du poisson fumé et une banane mûre en dessert. Au bout de quelques années, les enfants de la première famille présentent fréquemment fatigue, infections à répétition et retard de croissance, alors que ceux de la seconde grandissent normalement. \emph{Même budget, résultats différents : la différence tient à la diversité de l'assiette.}

\begin{definition}[Les trois groupes d'aliments]
On classe les aliments en trois grands groupes selon leur rôle principal dans l'organisme :
\begin{itemize}
\item les \cle{aliments énergétiques} : riches en glucides et/ou lipides, ils fournissent l'énergie nécessaire au métabolisme de base et à l'activité physique (céréales : maïs, riz, mil ; tubercules : manioc, igname, taro ; banane plantain ; huiles) ;
\item les \cle{aliments de construction} : riches en protéines de bonne qualité nutritionnelle, ils servent à fabriquer et réparer les tissus (viande, poisson, œufs, lait, légumineuses : niébé, arachide, soja) ;
\item les \cle{aliments protecteurs} : riches en vitamines, sels minéraux et fibres, ils protègent contre les carences et les infections et régulent le transit (légumes feuilles : folong, ndolé, gombo, éru ; fruits ; huile de palme rouge non raffinée, source de provitamine A).
\end{itemize}
\end{definition}

\begin{propriete}[Principe fondamental de l'alimentation préventive (admis)]
Une alimentation \textbf{variée}, associant à chaque repas principal un aliment énergétique, un aliment de construction et un aliment protecteur, couvre l'ensemble des besoins en macronutriments et micronutriments de l'organisme et prévient la \textbf{majorité des maladies nutritionnelles} (malnutrition protéino-énergétique, avitaminoses, anémie ferriprive). Cette propriété, établie par de nombreuses enquêtes épidémiologiques et les apports nutritionnels conseillés (ANC), est admise.
\end{propriete}

\begin{attention}[Erreur fréquente : « bien manger » = « manger beaucoup »]
Beaucoup croient que manger une grande quantité d'un seul aliment (riz, fufu de manioc) suffit à être bien nourri. C'est faux : ces aliments sont presque exclusivement énergétiques. Un enfant rassasié de fufu peut être carencé en protéines, en fer ou en vitamine A : on parle de \textbf{faim cachée} (carence en micronutriments malgré un apport calorique suffisant). La quantité ne remplace jamais la diversité.
\end{attention}

\begin{popfigure}
\begin{center}
\begin{tikzpicture}
% Assiette idéale : trois secteurs égaux
\draw[fill=popcream] (0,0) circle (3.2);
\draw[fill=popblueL] (0,0) -- (90:3) arc (90:210:3) -- cycle;
\draw[fill=poporangeL] (0,0) -- (210:3) arc (210:330:3) -- cycle;
\draw[fill=popgreenL] (0,0) -- (330:3) arc (330:450:3) -- cycle;
\draw[popdark,thick] (0,0) circle (3);
\draw[popdark,thick] (0,0) circle (3.2);
% Étiquettes internes
\node[popdark,font=\bfseries\small,align=center] at (150:1.6) {Aliments\\énergétiques};
\node[popdark,font=\bfseries\small,align=center] at (270:1.6) {Aliments de\\construction};
\node[popdark,font=\bfseries\small,align=center] at (30:1.6) {Aliments\\protecteurs};
% Légendes externes avec lignes fines
\node[anchor=west,align=left,font=\small] (A) at (4.2,2.2) {Maïs, riz, mil, manioc,\\igname, plantain, huile};
\draw[popmuted,thin] (150:2.6) -- (4.1,2.2);
\node[anchor=west,align=left,font=\small] (B) at (4.2,0) {Poisson, viande, œufs,\\niébé, arachide, soja};
\draw[popmuted,thin] (270:2.6) -- (4.1,0);
\node[anchor=west,align=left,font=\small] (C) at (4.2,-2.2) {Folong, ndolé, éru,\\fruits, huile de palme rouge};
\draw[popmuted,thin] (30:2.6) -- (4.1,-2.2);
\end{tikzpicture}
\end{center}
\legende{L'assiette idéale : chaque repas principal associe les trois groupes d'aliments — énergétiques, de construction et protecteurs — avec des exemples d'aliments couramment disponibles au Cameroun.}
\end{popfigure}

\begin{methode}[Construire un repas équilibré avec les aliments locaux]
Pour composer un repas équilibré à partir de ce qui est disponible au marché :
\begin{enumerate}
\item \textbf{Choisir une base énergétique} : riz, couscous de maïs, fufu, bâton de manioc, plantain, igname (une portion modérée, adaptée à l'appétit et à l'activité).
\item \textbf{Ajouter un aliment de construction} : poisson fumé ou frais, viande, œufs, ou une légumineuse (niébé, arachide, soja) — au moins une source de protéines (animale ou végétale) par repas principal.
\item \textbf{Ajouter un protecteur} : une sauce de feuilles (folong, ndolé, éru, gombo), des légumes crus ou cuits, puis un fruit de saison (mangue, papaye, banane, orange, goyave) en dessert.
\item \textbf{Vérifier la présence des trois groupes} dans l'assiette avant de servir ; si un groupe manque, compléter (un œuf, une poignée d'arachides, une orange suffisent souvent).
\item \textbf{Varier d'un jour à l'autre} : alterner les sources de chaque groupe (céréales, tubercules ; viandes, poissons, légumineuses ; différents légumes et fruits) pour couvrir tous les besoins en micronutriments.
\end{enumerate}
\end{methode}

\begin{savaistu}[Le ndolé, un plat presque complet]
Le ndolé, plat emblématique du Cameroun, est un modèle involontaire d'équilibre nutritionnel : les \textbf{feuilles de ndolé} sont des aliments protecteurs (vitamines A, C, fer, folates), la pâte d'\textbf{arachides} et les crevettes ou le bœuf apportent les protéines de construction, et l'accompagnement (miondo, bâton de manioc, riz, plantain) fournit l'énergie. Un seul plat traditionnel réunit donc les trois groupes : la preuve que nos cuisines locales savent, depuis longtemps, prévenir les carences.
\end{savaistu}

\courssub{Principe 2 : adapter les quantités aux besoins}

La diversité ne suffit pas : il faut aussi que les \textbf{quantités} correspondent aux besoins, qui varient selon l'individu.

\begin{definition}[Besoins énergétiques]
Les \cle{besoins énergétiques} sont la quantité d'énergie (exprimée en kilocalories, kcal) qu'un individu doit recevoir chaque jour pour couvrir son métabolisme de base, sa croissance éventuelle et son activité physique. Ils dépendent principalement de l'\textbf{âge}, du \textbf{sexe}, de l'\textbf{activité physique} et de situations physiologiques particulières (\textbf{grossesse}, allaitement, convalescence).
\end{definition}

\begin{exemplebox}[Un exemple chiffré : l'adolescent actif]
Selon les Apports Nutritionnels Conseillés (ANC OMS/FAO), un adolescent de 16 à 18 ans, actif (sport, marche quotidienne), a besoin d'environ \textbf{\num{2500} à \num{3000} kcal par jour} (garçon) ou \num{2200} à \num{2500} kcal (fille), réparties sur trois repas :
\begin{itemize}
\item petit-déjeuner : environ 20 à 25 \% de l'apport (environ 600 kcal) ;
\item déjeuner : environ 35 à 40 \% (environ \num{1000} kcal) ;
\item dîner : environ 30 à 35 \% (environ 900 kcal), plus une collation éventuelle.
\end{itemize}
À titre de comparaison, une femme enceinte doit augmenter ses apports (environ +350 kcal/jour au dernier trimestre) et surtout ses apports en fer, folates et protéines ; un jeune enfant a des besoins absolus plus faibles mais des besoins \emph{relatifs} (par kilogramme de masse corporelle) plus élevés, car il grandit rapidement.
\end{exemplebox}

\begin{attention}[Deux erreurs symétriques]
Manger \textbf{trop peu} par rapport à ses besoins conduit à la malnutrition protéino-énergétique et aux carences ; manger \textbf{trop} conduit au surpoids puis à l'obésité. L'objectif n'est ni de se restreindre ni de se gaver, mais d'\textbf{adapter} la ration à ses besoins réels — un élève qui fait deux heures de football n'a pas les mêmes besoins qu'une journée sédentaire.
\end{attention}

\courssub{Principe 3 : limiter les excès}

L'alimentation préventive ne consiste pas seulement à ajouter les bons aliments, mais aussi à \textbf{limiter} ceux qui, en excès, provoquent les maladies de surcharge étudiées précédemment.

\begin{aretenir}[Les quatre excès à limiter]
\begin{itemize}
\item \textbf{Sucres rapides} (boissons sucrées industrielles, bonbons, excès de sucre dans le thé ou les bouillies) : favorisent l'obésité, le diabète de type 2 et les caries dentaires.
\item \textbf{Graisses} (fritures répétées, excès d'huile dans les sauces, viandes très grasses, peau de poulet) : favorisent l'obésité et les maladies cardiovasculaires (athérosclérose, hypertension).
\item \textbf{Sel} (chlorure de sodium) : l'excès favorise l'hypertension artérielle ; attention aux cubes d'assaisonnement industriels et aux poissons très salés, déjà riches en sodium.
\item \textbf{Alcool} : il apporte des calories « vides » (7 kcal/g sans nutriments essentiels), détruit les cellules du foie (cirrhose) et du système nerveux, et aggrave les carences en perturbant l'absorption et le métabolisme des nutriments (notamment vitamine B1, folates). Il est à proscrire chez l'enfant, l'adolescent et la femme enceinte.
\end{itemize}
\end{aretenir}

\courssub{Principe 4 : valoriser les ressources alimentaires locales}

Une idée reçue veut que « bien manger » coûte cher et exige des produits importés. C'est faux : le Cameroun dispose de ressources locales remarquables, accessibles et nutritives.

\begin{center}
\begin{tabularx}{\linewidth}{|l|X|X|}
\hline
\rowcolor{popblueL} \textbf{Aliment local} & \textbf{Groupe} & \textbf{Atout nutritionnel majeur} \\
\hline
Niébé, soja & Construction & Protéines végétales bon marché, fer, fibres \\
\hline
Arachide & Construction / Énergétique & Protéines, lipides (acides gras insaturés) ; base de nombreuses sauces \\
\hline
Poisson fumé (ex. bonga, capitaine) & Construction & Protéines de haute qualité, calcium (arêtes consommées), iode, accessible partout \\
\hline
Folong, ndolé, éru, gombo & Protecteurs & Vitamines A et C, fer, folates, fibres, antioxydants \\
\hline
Fruits de saison (mangue, papaye, goyave, orange, ananas) & Protecteurs & Vitamine C (améliore l'absorption du fer non hémique), bêta-carotène, fibres \\
\hline
Huile de palme rouge (non raffinée) & Énergétique / Protecteur & Très riche en provitamine A (bêta-carotène) : prévient la xérophtalmie et la cécité nocturne \\
\hline
\end{tabularx}
\end{center}

\begin{exemplebox}[Un geste simple et puissant]
Remplacer l'huile raffinée par de l'\textbf{huile de palme rouge non raffinée} dans la sauce de feuilles, et ajouter une \textbf{orange} ou quelques quartiers de goyave en fin de repas, améliore à la fois l'apport en vitamine A (provitamine A de l'huile) et l'absorption du fer non hémique des feuilles grâce à la vitamine C du fruit : deux carences majeures (avitaminose A et anémie ferriprive) prévenues pour quelques francs.
\end{exemplebox}

\courssub{Principe 5 : hygiène alimentaire et hydratation}

Les infections et les carences s'\textbf{aggravent mutuellement} : un enfant carencé s'infecte plus facilement (immunité affaiblie), et chaque infection (diarrhée, paludisme, parasitose intestinale) augmente les pertes (malabsorption, pertes sanguines, catabolisme) et les besoins nutritionnels. Casser ce cercle vicieux fait partie intégrante de la prévention nutritionnelle.

\begin{aretenir}[Règles d'hygiène alimentaire et d'hydratation]
\begin{itemize}
\item Boire une \textbf{eau potable} (traitée, bouillie, filtrée ou de source protégée) en quantité suffisante : environ 1,5 à 2 litres par jour pour un adolescent, davantage par forte chaleur ou lors d'efforts physiques intenses.
\item \textbf{Laver les mains} à l'eau et au savon avant de préparer ou de manger les aliments, et après être allé aux toilettes (interruption de la chaîne fécale-orale).
\item \textbf{Laver} soigneusement fruits et légumes consommés crus (eau + vinaigre ou javel alimentaire diluée) ; \textbf{bien cuire} viandes et poissons (cuisson à cœur) ; protéger les aliments des mouches, de la poussière et des rongeurs.
\item Conserver les restes au frais (réfrigérateur ou glacière) et les réchauffer correctement (ébullition) ; ne pas consommer d'aliments avariés (aspect, odeur, gonflement de la boîte).
\end{itemize}
\end{aretenir}

\courssub{Mesures collectives de prévention}

La lutte contre les maladies nutritionnelles ne se joue pas seulement dans l'assiette individuelle : l'État, les services de santé et les communautés mènent des actions collectives essentielles, relayées par l'éducation à la santé.

\begin{itemize}
\item \textbf{Allaitement maternel exclusif} pendant les six premiers mois de vie : le lait maternel couvre tous les besoins nutritionnels et hydriques du nourrisson, lui transmet des anticorps (immunité passive), est gratuit, stérile et toujours à la bonne température. Complémentation à partir de 6 mois avec poursuite de l'allaitement jusqu'à 2 ans ou plus.
\item \textbf{Fortification des aliments de base} : ajout obligatoire de micronutriments à des aliments de consommation courante — l'exemple phare est le \textbf{sel iodé} (iodation à 20-40 ppm), qui prévient le goitre et les troubles du développement cognitif liés à la carence en iode ; certaines farines de blé et huiles sont enrichies en fer, acide folique, zinc, vitamine A.
\item \textbf{Vaccination} (PEV - Programme Élargi de Vaccination) : en prévenant les infections graves (rougeole, coqueluche, diphtérie, tuberculose...), elle évite les poussées de malnutrition aiguë que ces maladies déclenchent chez l'enfant.
\item \textbf{Déparasitage} régulier des enfants en âge scolaire (vers intestinaux : ascaris, ankylostomes, trichures ; schistosomes) : ces parasites volent les nutriments, provoquent des saignements intestinaux (anémie) et altèrent la croissance.
\item \textbf{Éducation nutritionnelle en milieu scolaire} : cantines scolaires équilibrées (respect des trois groupes), jardins scolaires pédagogiques, leçons d'éducation à la santé — c'est précisément le rôle de ce module.
\end{itemize}

\begin{savaistu}[Le sel iodé, une victoire silencieuse de santé publique]
La carence en iode provoquait autrefois de nombreux cas de goitre endémique et de retard mental (crétinisme) dans certaines régions montagneuses du Cameroun (Adamaoua, Nord-Ouest, Ouest), où les sols lessivés sont pauvres en iode. L'iodation obligatoire du sel de cuisine, mesure simple, universelle et très peu coûteuse (quelques centimes par kg de sel), a fait reculer spectaculairement ces troubles depuis les années 1990 : c'est l'un des plus beaux succès de la nutrition préventive en Afrique.
\end{savaistu}

\courssub{Conduite à tenir devant un cas suspect}

Savoir reconnaître les signes (section précédente) et connaître les principes de prévention ne doit pas conduire à soigner soi-même : la place de l'élève, de l'enseignant et de la famille est d'\textbf{alerter et d'orienter} vers le système de santé.

\begin{methode}[Conduite à tenir devant un cas suspect de maladie nutritionnelle]
\begin{enumerate}
\item \textbf{Repérer} les signes d'alerte : amaigrissement récent ou prise de poids excessive rapide, œdèmes (pieds, visage), pâleur des conjonctives et paumes, cheveux décolorés/dépigmentés (signe du kwashiorkor), fatigue chronique, retard de croissance staturo-pondéral (courbe de croissance qui s'aplatit ou descend).
\item \textbf{Ne pas se contenter d'automédication} : acheter des vitamines, du « sanguin » ou des « toniques » en pharmacie ou au marché sans diagnostic médical peut masquer la maladie, être inadapté (mauvais dosage, mauvaise forme) ou dangereux (survitaminose A, surcharge en fer).
\item \textbf{Orienter} la personne vers un centre de santé intégré (CSI), un hôpital de district ou un hôpital régional : seul un personnel de santé qualifié (médecin, infirmier, diététicien) peut établir le diagnostic étiologique (examen clinique, pesée, taille, tour de bras, dosage de l'hémoglobine, recherche de parasites, sérologies...) et prescrire le traitement adapté (supplémentation ciblée, alimentation thérapeutique prête à l'emploi - ATPE pour la malnutrition aiguë sévère, traitement d'une cause associée : paludisme, tuberculose, VIH, parasitose).
\item \textbf{Accompagner} la correction médicale par l'application rigoureuse des principes de cette leçon : repas diversifiés (trois groupes), hygiène alimentaire, eau potable, suivi régulier de la courbe de croissance ou de l'IMC.
\end{enumerate}
\end{methode}

\begin{exoresolu}[Construire un menu équilibré pour un élève de Première]
Énoncé. Un élève de Première de 17 ans, actif (marche, sport occasionnel), dispose des aliments suivants au marché local : riz, huile de palme rouge non raffinée, poisson fumé (bonga), niébé (haricots blancs), feuilles de folong, tomates, bananes plantain, oranges, arachides grillées, pain, œufs, avocat. Proposer un menu d'une journée respectant les trois groupes d'aliments à chaque repas principal, puis justifier chaque choix.
\tcblower
Solution détaillée.
\begin{center}
\begin{tabularx}{\linewidth}{|l|X|X|X|}
\hline
\rowcolor{popgreenL} \textbf{Repas} & \textbf{Énergétique} & \textbf{Construction} & \textbf{Protecteur} \\
\hline
Petit-déjeuner & Pain (glucides complexes) & Œuf bouilli (protéines animales, fer, vitamine B12) & Une orange (vitamine C, fibres) \\
\hline
Déjeuner & Riz cuit à l'huile de palme rouge (glucides + provitamine A) & Sauce de niébé au poisson fumé (protéines végétales + animales, fer, zinc) & Feuilles de folong cuites (vitamines A, C, fer, folates) ; banane plantain mûre en dessert (énergie, potassium, bêta-carotène) \\
\hline
Dîner & Banane plantain bouillie (glucides, fibres) & Sauce arachide (arachides grillées écrasées) + reste de poisson fumé (protéines, lipides, calcium) & Salade de tomates et avocat (vitamine C, potassium, acides gras insaturés, fibres) \\
\hline
\end{tabularx}
\end{center}
\textbf{Justification :} Chaque repas principal associe explicitement les trois groupes. La complémentarité protéique est assurée (céréales + légumineuses au déjeuner ; légumineuses + poisson au dîner). L'apport en fer (niébé, poisson, folong) est optimisé par la présence systématique de vitamine C (orange, folong, tomate) qui en favorise l'absorption intestinale. L'huile de palme rouge assure l'apport en provitamine A. L'ensemble atteint environ \num{2600} à \num{2800} kcal, adapté à un adolescent actif, et n'utilise que des aliments locaux accessibles et bon marché.
\end{exoresolu}

\begin{exercice}[À toi de jouer]
Dans ta famille ou ton quartier, relève la composition détaillée des repas d'une journée complète (petit-déjeuner, déjeuner, dîner). Pour chaque repas, indique quels groupes d'aliments (énergétique, construction, protecteur) sont présents et lesquels manquent, puis propose une correction simple, peu coûteuse et culturellement acceptable utilisant uniquement des aliments locaux disponibles.
\end{exercice}

\begin{corrige}[Exercice 1]
Exemple de réponse attendue (scénario fréquent en milieu urbain) :
\begin{itemize}
\item \textbf{Petit-déjeuner :} Bouillie de maïs très sucrée, sans lait. \rightarrow \textbf{Manquent :} Construction + Protecteur. \textbf{Correction :} Ajouter du lait caillé (lait caillé + bouillie = « bouillie enrichie ») ou un œuf bouilli + un fruit de saison (banane, orange).
\item \textbf{Déjeuner :} Bâton de manioc (fufu) + sauce de feuilles (ndolé/folong) sans viande ni poisson ni légumineuse. \rightarrow \textbf{Manque :} Construction. \textbf{Correction :} Ajouter du poisson fumé (bonga), des crevettes séchées, du niébé ou de la pâte d'arachide dans la sauce.
\item \textbf{Dîner :} Riz blanc + huile + piment. \rightarrow \textbf{Manquent :} Construction + Protecteur. \textbf{Correction :} Préparer une sauce tomate/arachide avec du poisson ou des œufs, servir avec des légumes verts (feuilles) et un fruit en dessert.
\end{itemize}
La démarche correcte consiste toujours à : (1) dresser l'inventaire des groupes présents/absents, (2) combler les manques par des aliments locaux accessibles (légumineuses, poisson fumé, fruits de saison, huile rouge), (3) rappeler que les quantités doivent rester adaptées aux besoins de chacun (âge, sexe, activité).
\end{corrige}

\begin{aretenir}[L'essentiel de la section — Alimentation préventive]
\begin{itemize}
\item \textbf{Diversité :} Chaque repas principal doit associer les \textbf{trois groupes} (énergétique, construction, protecteurs) — la diversité prime sur la quantité ; c'est la clé pour éviter la faim cachée.
\item \textbf{Quantités adaptées :} Les apports caloriques s'ajustent aux besoins réels (âge, sexe, activité, grossesse/allaitement) : environ \num{2500}--\num{3000} kcal/jour pour un adolescent actif (garçon).
\item \textbf{Modération :} Limiter les quatre excès délétères : sucres rapides, graisses saturées/fritures, sel (cubes, poisson salé), alcool (interdit < 18 ans, femme enceinte).
\item \textbf{Valorisation du local :} Bien manger ne coûte pas cher : niébé, soja, arachide, poisson fumé, folong, ndolé, éru, fruits de saison, huile de palme rouge non raffinée sont des trésors nutritionnels accessibles.
\item \textbf{Hygiène \& Hydratation :} Eau potable (1,5--2 L/j), lavage des mains, lavage/cuisson des aliments, conservation sûre brisent le cercle vicieux infections $\leftrightarrow$ carences.
\item \textbf{Action collective :} Allaitement maternel exclusif 6 mois, sel iodé, vaccination (PEV), déparasitage scolaire, cantines équilibrées, éducation nutritionnelle.
\item \textbf{Conduite à tenir :} Devant tout signe de maladie nutritionnelle $\rightarrow$ \textbf{Orienter vers un centre de santé} (diagnostic + traitement), jamais d'automédication ; accompagner par une alimentation diversifiée et un suivi anthropométrique.
\end{itemize}
\end{aretenir}

\newpage

\coursec{Activité d'intégration}

\courssub{Objectifs de l'activité}

À l'issue de cette activité d'intégration, tu seras capable de :

\begin{itemize}
\item identifier les signes d'une maladie nutritionnelle à partir d'un cas concret (malnutrition protéino-énergétique, anémie, obésité) ;
\item calculer et interpréter l'\cle{indice de masse corporelle} (IMC) d'un patient ;
\item lire et interpréter une \cle{courbe de croissance} d'un jeune enfant ;
\item proposer des mesures de \cle{prévention} et de \cle{correction} adaptées à chaque cas, en utilisant des aliments locaux disponibles au Cameroun ;
\item travailler en équipe et présenter oralement des conclusions argumentées.
\end{itemize}

\courssub{Situation}

Tu es stagiaire au centre de santé intégré du quartier Nkol-Eton, à Yaoundé. Ce samedi matin, l'infirmier major t'accueille avec trois dossiers de patients et te demande de l'aider à préparer la consultation d'éducation nutritionnelle de l'après-midi. Par groupes de quatre, tu vas analyser chaque dossier comme un véritable conseiller nutritionnel.

\textbf{Dossier 1 — Petit Arnaud, 3 ans.} Sa mère l'amène car « il gonfle et ses cheveux changent de couleur ». À l'examen : poids 9 kg (au lieu d'environ 14 kg attendus à 3 ans), œdèmes des pieds et du visage, cheveux décolorés tirant sur le roux, apathie. Le carnet de santé montre que sa courbe de poids, qui suivait normalement le couloir de croissance jusqu'à 18 mois, chute nettement depuis le sevrage et l'arrivée d'un petit frère. L'alimentation actuelle est constituée presque exclusivement de bouillie de maïs.

\textbf{Dossier 2 — Junior, 16 ans.} Lycéen, il mesure \SI{1,62}{m} et pèse \SI{88}{kg}. Il consomme chaque jour deux à trois bouteilles de boissons sucrées et grignote beignets et chips devant la télévision. Il ne pratique aucune activité physique. Il se plaint d'essoufflement à l'effort.

\textbf{Dossier 3 — Sandrine, 15 ans.} Elle est pâle, fatiguée, essoufflée à la montée des escaliers, et ses résultats scolaires baissent. Elle déclare ne manger « presque jamais de viande ni de légumes », essentiellement du riz et du pain. Le dosage sanguin prescrit montre une hémoglobine à \SI{9}{g/dL} (norme : 12 à \SI{16}{g/dL} chez l'adolescente).

\courssub{Consignes}

Pour chaque dossier, et dans l'ordre :

\begin{enumerate}
\item Identifie la maladie nutritionnelle probable et justifie ton diagnostic en citant précisément les signes cliniques du dossier qui t'orientent.
\item Pour Junior, calcule son IMC, classe-le selon les normes et conclus sur son état nutritionnel.
\item Pour Arnaud, décris ce que montre sa courbe de croissance et explique pourquoi la période du sevrage est critique.
\item Pour Sandrine, nomme le type de carence en cause et cite le nutriment manquant ainsi que deux aliments locaux qui en sont riches.
\item Rédige pour chaque patient une fiche-conseil comportant : un menu type d'une journée (aliments locaux uniquement), deux mesures d'hygiène de vie, et un conseil de prévention adressé à la famille.
\item Prépare une restitution orale de 3 minutes par groupe ; les autres groupes évaluent avec la grille commune (exactitude du diagnostic, justesse des calculs, qualité des conseils, clarté de l'exposé).
\end{enumerate}

\courssub{Ressources à mobiliser}

\begin{aretenir}[Outils de la leçon à réutiliser]
\begin{itemize}
\item \cle{IMC} : $\mathrm{IMC} = \dfrac{\text{masse (kg)}}{\text{taille}^2\ (\mathrm{m}^2)}$. Interprétation : $< 18,5$ maigreur ; $18,5$ à $25$ corpulence normale ; $25$ à $30$ surpoids ; $> 30$ obésité.
\item Signes du \cle{kwashiorkor} (MPE) : œdèmes, cheveux décolorés, retard de croissance, apathie, chute de la courbe de poids après le sevrage.
\item \cle{Anémie} ferriprive : pâleur, fatigue, hémoglobine abaissée ; liée au manque de fer (viande, légumes verts, légumineuses).
\item \cle{Obésité} : excès de graisses et de sucres, sédentarité ; risques de diabète et de maladies cardiovasculaires.
\item Alimentation préventive : repas variés associant aliments énergétiques, de construction (protéines) et protecteurs (vitamines et minéraux).
\end{itemize}
\end{aretenir}

\courssub{Démarche et corrigé commenté}

\begin{exoresolu}[Corrigé commenté des trois dossiers]
\textbf{Dossier 1 — Arnaud.} Les signes (œdèmes, cheveux décolorés, poids très inférieur à la norme, apathie) sont typiques du \cle{kwashiorkor}, forme de malnutrition protéino-énergétique par carence en protéines. La courbe de croissance montre une cassure à 18 mois, au moment du sevrage : la bouillie de maïs seule couvre l'énergie mais pas les besoins en protéines. \emph{Fiche-conseil :} enrichir la bouillie (farine de soja, pâte d'arachide), introduire poisson fumé, œufs, haricots niébé ; petits repas fréquents ; suivi mensuel de la courbe de poids au centre de santé.

\textbf{Dossier 2 — Junior.} Calcul de l'IMC :
\[ \mathrm{IMC} = \frac{88}{1,62^2} = \frac{88}{2,6244} \approx 33,5 \]
Un IMC supérieur à 30 indique une \cle{obésité}. Causes : excès de boissons sucrées et de grignotage, sédentarité. Risques : diabète de type 2, hypertension, maladies cardiovasculaires. \emph{Fiche-conseil :} remplacer les sodas par de l'eau, réduire fritures et sucreries, 30 minutes de marche ou de sport par jour, repas équilibrés (légumes, poisson, tubercules en quantité modérée).

\textbf{Dossier 3 — Sandrine.} Pâleur, fatigue et hémoglobine à \SI{9}{g/dL} (norme 12 à 16) : il s'agit d'une \cle{anémie ferriprive}, carence en fer liée à l'absence de viande et de légumes. Le fer est indispensable à la fabrication de l'hémoglobine. \emph{Fiche-conseil :} consommer régulièrement viande, poisson, folong, feuilles de manioc, niébé ; ajouter un fruit riche en vitamine C (orange, goyave) au repas pour améliorer l'absorption du fer ; consultation médicale pour un éventuel supplément en fer.
\end{exoresolu}

\courssub{Bilan de l'activité}

Cette activité t'a placé dans la peau d'un véritable conseiller nutritionnel. Elle a permis de mettre en évidence que :

\begin{itemize}
\item chaque maladie nutritionnelle possède des \cle{signes caractéristiques} qu'une observation attentive suffit souvent à repérer ;
\item des outils simples et chiffrés — l'\cle{IMC} et la \cle{courbe de croissance} — permettent d'objectiver un diagnostic et de suivre l'état nutritionnel ;
\item les déséquilibres nutritionnels ont deux visages opposés, la \cle{carence} et la \cle{surcharge}, mais une cause commune : une alimentation mal équilibrée ;
\item la prévention repose sur des repas \emph{variés} construits à partir des trois groupes d'aliments, facilement réalisables avec les produits locaux (niébé, arachide, poisson, légumes-feuilles, fruits) ;
\item l'éducation nutritionnelle en milieu communautaire, comme au centre de santé de quartier, est un acte de santé publique essentiel au Cameroun.
\end{itemize}

Tu es désormais capable d'identifier, d'expliquer et de prévenir les principales maladies nutritionnelles autour de toi : c'est une compétence utile pour ta famille, ton quartier et ton avenir.

\newpage

\coursec{Méthodes et exercices résolus}

\coursec{Lutte contre les maladies nutritionnelles}

\courssub{Généralités sur les maladies nutritionnelles}

\begin{definition}[Malnutrition et déséquilibres nutritionnels]
On appelle \cle{malnutrition} tout état pathologique résultant d'un apport inadapté en nutriments (énergie, protéines, vitamines, sels minéraux) par rapport aux besoins de l'organisme. Elle se présente sous deux formes majeures :
\begin{itemize}
\item la \cle{sous-nutrition} (ou dénutrition) : apports insuffisants $\rightarrow$ maladies de carence (malnutrition protéino-énergétique, avitaminoses, anémies) ;
\item la \cle{surnutrition} : apports excessifs (surtout énergétiques) $\rightarrow$ maladies de surcharge (surpoids, obésité).
\end{itemize}
\end{definition}

\begin{propriete}[Classification des maladies nutritionnelles]
\begin{center}
\begin{tabularx}{\linewidth}{|l|X|X|}
\hline
\rowcolor{popblueL} \textbf{Type} & \textbf{Maladies principales} & \textbf{Cause fondamentale} \\
\hline
Maladies de carence & \begin{itemize}\item Malnutrition protéino-énergétique (kwashiorkor, marasme)\item Avitaminoses (A, C, D, B\textsubscript{1}...)\item Anémies (fer, folates, B\textsubscript{12})\end{itemize} & Déficit en nutriments spécifiques ou globaux (apport insuffisant, pertes accrues, besoins augmentés, malabsorption) \\
\hline
Maladies de surcharge & \begin{itemize}\item Surpoids\item Obésité (modérée, sévère, massive)\end{itemize} & Excès d'apports énergétiques (lipides, glucides simples) par rapport aux dépenses, sédentarité \\
\hline
\end{tabularx}
\end{center}
\textbf{Conséquences communes} : affaiblissement du système immunitaire, retard de croissance (enfant), baisse de la capacité de travail (adulte), surmortalité infantile et maternelle, charge économique pour les familles et le système de santé.
\end{propriete}

\begin{aretenir}[Vocabulaire du Probatoire]
\begin{itemize}
\item \textbf{MPE} = Malnutrition Protéino-Énergétique (terme parapluie).
\item \textbf{Kwashiorkor} = carence protéique prédominante (avec œdèmes).
\item \textbf{Marasme} = carence énergétique globale (amaigrissement extrême, sans œdème).
\item \textbf{Anémie ferriprive} = anémie par carence en fer (sidéropénie).
\item \textbf{IMC} = Indice de Masse Corporelle ($m/t^2$), valide \textbf{uniquement chez l'adulte}.
\item \textbf{Courbe de croissance} = outil de référence pour le suivi nutritionnel de l'enfant (0--5 ans).
\end{itemize}
\end{aretenir}

\courssub{La malnutrition protéino-énergétique (MPE)}

\begin{definition}[Kwashiorkor]
Forme de MPE due à une \cle{carence en protéines} relativement isolée (apport énergétique souvent maintenu par les féculents). Survient typiquement après un \textbf{sevrage précoce} vers une alimentation à base de bouillie de maïs, de manioc ou de banane, pauvre en protéines de bonne qualité.
\end{definition}

\begin{definition}[Marasme]
Forme de MPE due à une \cle{carence énergétique globale} (protéines + lipides + glucides). L'organisme consomme ses propres réserves (graisse, puis muscle) pour survivre. Fréquent chez le nourrisson non allaité ou mal nourri, ou lors d'infections répétées (diarrhées, rougeole, paludisme).
\end{definition}

\begin{propriete}[Signes cliniques distinctifs : Kwashiorkor vs Marasme]
\begin{center}
\begin{tabularx}{\linewidth}{|l|X|X|}
\hline
\rowcolor{popblueL} \textbf{Signe} & \textbf{Kwashiorkor (carence protéique)} & \textbf{Marasme (carence globale)} \\
\hline
\cle{Œdèmes} & \textbf{Présents} (pieds, visage, abdomen) : hypoalbuminémie $\downarrow$ pression oncotique & \textbf{Absents} \\
\hline
\cle{Amaigrissement} & Modéré, masqué par les œdèmes & \textbf{Extrême} : peau sur les os, \emph{aspect vieillard} \\
\hline
\cle{Cheveux} & Décolorés (roux), cassants, \emph{flag sign} & Fins, clairsemés mais pas décolorés \\
\hline
\cle{Peau} & \cle{Dermatose} : plaques squameuses, hyperpigmentées (\emph{peau en peinture écaillée}) & Sèche, ridée, atrophique \\
\hline
\cle{Abdomen} & Ballonné (hépatomégalie graisseuse + ascite) & Rentré (perte de graisse mésentérique) \\
\hline
\cle{Comportement} & Apathique, misérable, anorexique & Irritable, affamé, alerte (au début) \\
\hline
\cle{Âge fréquent} & 1--3 ans (post-sevrage) & 6--18 mois (sevrage précoce, biberon mal dosé) \\
\hline
\end{tabularx}
\end{center}
\textbf{Note} : existent des formes mixtes \cle{marasme-kwashiorkor} (amaigrissement sévère + œdèmes).
\end{propriete}

\begin{experience}[Physiopathologie des œdèmes du kwashiorkor (démonstration guidée)]
\textbf{Question} : Pourquoi une carence en protéines provoque-t-elle des œdèmes ?
\begin{enumerate}
\item \textbf{Rappel} : l'albumine (protéine plasmatique majeure, synthétisée par le foie) maintient la \cle{pression oncotique} (pression colloïdosmotique) du sang $\approx \SI{25}{mmHg}$, qui retient l'eau dans les vaisseaux.
\item \textbf{Carence protéique} $\rightarrow$ synthèse hépatique d'albumine $\downarrow$ $\rightarrow$ hypoalbuminémie $\rightarrow$ pression oncotique $\downarrow$.
\item \textbf{Déséquilibre de Starling} : la pression hydrostatique (qui pousse l'eau hors des capillaires) devient supérieure à la pression oncotique $\rightarrow$ fuite d'eau vers l'espace interstitiel $\rightarrow$ \textbf{œdèmes}.
\item \textbf{Conséquence hépatique} : le foie, sollicité pour synthétiser des lipoprotéines de transport (VLDL) afin d'exporter les triglycérides, s'engorge en graisse $\rightarrow$ \textbf{stéatose hépatique} (hépatomégalie, abdomen ballonné).
\end{enumerate}
\end{experience}

\courssub{Identifier les signes d'une maladie nutritionnelle à partir d'un cas}

\begin{methode}[Reconnaître une maladie nutritionnelle à partir d'une observation]
Face à un cas clinique décrit dans un énoncé (enfant, adulte), procéder en 5 étapes :
\begin{enumerate}
\item \textbf{Repérer les données chiffrées} : âge, masse, taille, périmètre brachial, taux d'hémoglobine. Les comparer aux valeurs de référence (courbe de croissance, IMC, taux normal d'hémoglobine : \SI{12}{--}\SI{16}{g/dL} chez la femme, \SI{13}{--}\SI{17}{g/dL} chez l'homme).
\item \textbf{Relever les signes cliniques} décrits : œdèmes (gonflements), amaigrissement, cheveux décolorés et cassants, peau squameuse, fatigue, pâleur, retard de croissance, surpoids.
\item \textbf{Associer chaque signe à une cause possible} : œdèmes + cheveux décolorés $\Rightarrow$ kwashiorkor (carence en protéines) ; amaigrissement extrême sans œdème $\Rightarrow$ marasme (carence énergétique globale) ; pâleur + fatigue $\Rightarrow$ anémie (carence en fer) ; cécité nocturne $\Rightarrow$ carence en vitamine A ; saignement des gencives $\Rightarrow$ scorbut (carence en vitamine C).
\item \textbf{Nommer précisément la maladie} et la carence ou la surcharge en cause (jamais de réponse vague du type « l'enfant est malade »).
\item \textbf{Conclure par une phrase complète} reliant les signes observés au diagnostic : « L'enfant présente ... car ... ; il s'agit donc de ... ».
\end{enumerate}
\textbf{Réflexe Probatoire} : toujours citer dans la réponse au moins deux signes de l'énoncé pour justifier le diagnostic.
\end{methode}

\begin{exoresolu}[Diagnostic d'un cas de malnutrition à l'hôpital de district]
Dans un centre de santé de la région de l'Extrême-Nord, une mère amène son enfant âgé de 3 ans. Le médecin note : masse de \SI{9}{kg} (la masse attendue à cet âge est d'environ \SI{14}{kg}), présence d'œdèmes aux pieds et au visage, cheveux roux, fragiles et se détachant facilement, peau présentant des plaques squameuses, abdomen ballonné. L'enfant a été sevré précocement et nourri essentiellement de bouillie de maïs.
\begin{enumerate}
\item Identifier la maladie dont souffre cet enfant en justifiant par deux arguments.
\item Préciser la carence en cause.
\end{enumerate}
\tcblower
\textbf{1. Identification de la maladie.}

Analysons les données de l'énoncé :
\begin{itemize}
\item la masse de l'enfant (\SI{9}{kg}) est très inférieure à la masse attendue (\SI{14}{kg}) : il existe un \cle{retard de croissance} (insuffisance pondérale) ;
\item l'enfant présente des \cle{œdèmes} (pieds, visage), un abdomen ballonné, des cheveux décolorés (roux) et cassants, une peau squameuse (dermatose).
\end{itemize}
Or la présence d'œdèmes associée à des cheveux décolorés et à une peau anormale, malgré un amaigrissement modéré, est la signature du \cle{kwashiorkor}. Dans le marasme, au contraire, l'amaigrissement est extrême mais il n'y a \emph{pas} d'œdèmes.

\textbf{Conclusion} : l'enfant souffre de kwashiorkor, car il présente un déficit pondéral associé à des œdèmes, des cheveux décolorés et des lésions cutanées.

\textbf{2. Carence en cause.}

Le kwashiorkor est une forme de \cle{malnutrition protéino-énergétique} (MPE) due principalement à une \cle{carence en protéines} (protéino-malnutrition), l'apport énergétique étant relativement conservé. Cela correspond à la situation de l'énoncé : l'enfant est nourri de bouillie de maïs, aliment énergétique mais très pauvre en protéines de bonne qualité. Le déficit en protéines plasmatiques (albumine) diminue la pression oncotique du sang, d'où la fuite d'eau vers les tissus et les œdèmes.

\textbf{Conclusion} : la carence en cause est une carence protéique, consécutive à un sevrage précoce et à une alimentation déséquilibrée dominée par les féculents.
\end{exoresolu}

\courssub{Les avitaminoses et l'anémie}

\begin{definition}[Avitaminose]
Maladie due à une carence prolongée en une \cle{vitamine} (nutriment organique indispensable en très faible quantité, non synthétisé par l'organisme). Les principales avitaminoses au Cameroun :
\begin{itemize}
\item \textbf{Vitamine A} (rétinol) : \cle{xérophtalmie} (sécheresse conjonctivale, taches de Bitot, cécité nocturne $\rightarrow$ cécité totale), baisse de l'immunité. Sources : huile de palme rouge, foie, jaune d'œuf, carotte, mangue, patate douce à chair orange.
\item \textbf{Vitamine C} (acide ascorbique) : \cle{scorbut} (saignement des gencives, hématomes, retard de cicatrisation, douleurs ostéo-articulaires). Sources : goyave, agrumes, papaye, chou, poivron, folon, ndolé (cuisson courte).
\item \textbf{Vitamine D} (calciférol) : \cle{rachitisme} (enfant : déformations osseuses, jambes arquées, rosaire costal, retard de marche ; adulte : ostéomalacie). Synthèse cutanée sous UVB + apport alimentaire (poissons gras, foie, jaune d'œuf).
\item \textbf{Vitamine B\textsubscript{1}} (thiamine) : \cle{béribéri} (atteinte neurologique : paresthésies, paralysies ; cardiaque : insuffisance cardiaque, œdèmes). Sources : son de céréales, légumineuses, porc, levure de bière.
\end{itemize}
\end{definition}

\begin{definition}[Anémie]
Diminution de la concentration en \cle{hémoglobine} dans le sang en dessous des seuils normaux (\SI{<12}{g/dL} femme, \SI{<13}{g/dL} homme, \SI{<11}{g/dL} femme enceinte/enfant). L'anémie la plus fréquente est l'\cle{anémie ferriprive} (carence en fer, sidéropénie), car le fer est au centre du groupe hème de l'hémoglobine. Autres causes : carence en folates (vit B\textsubscript{9}) ou B\textsubscript{12} (anémie macrocytaire), hémolyses (drépanocytose, paludisme), hémorragies.
\end{definition}

\begin{propriete}[Signes cliniques évocateurs d'anémie]
Pâleur des téguments et des conjonctives, fatigue disproportionnée à l'effort, essoufflement, tachycardie, vertiges, céphalées, baisse des performances scolaires/professionnelles. \textbf{Confirmation biologique} : Numération-Formule Sanguine (NFS) : Hémoglobine basse, VGM (Volume Globulaire Moyen) bas (microcytaire) $\Rightarrow$ fer ; VGM haut (macrocytaire) $\Rightarrow$ folates/B\textsubscript{12}.
\end{propriete}

\courssub{Identifier une avitaminose ou une anémie}

\begin{methode}[Associer un signe clinique à une carence en vitamine ou en fer]
\begin{enumerate}
\item \textbf{Classer le signe observé} selon l'organe atteint : œil, gencives et peau, os, sang.
\item \textbf{Appliquer le tableau de correspondance} à connaître par cœur :
\end{enumerate}
\begin{center}
\begin{tabularx}{\linewidth}{|l|X|X|}
\hline
\rowcolor{popblueL} \textbf{Signe observé} & \textbf{Carence} & \textbf{Maladie} \\
\hline
Cécité nocturne, sécheresse oculaire, taches de Bitot & Vitamine A & Xérophtalmie \\
\hline
Saignement des gencives, retard de cicatrisation, hématomes & Vitamine C & Scorbut \\
\hline
Déformations osseuses (jambes arquées, rosaire costal) chez l'enfant & Vitamine D (et calcium) & Rachitisme \\
\hline
Pâleur, fatigue, essoufflement, taux d'hémoglobine bas (NFS) & Fer (sels minéraux) & Anémie ferriprive \\
\hline
\end{tabularx}
\end{center}
\begin{enumerate}
\setcounter{enumi}{2}
\item \textbf{Vérifier la cohérence avec le mode de vie} décrit (alimentation pauvre en fruits et légumes, en produits animaux, manque d'exposition au soleil, paludisme chronique).
\item \textbf{Rédiger la réponse} en nommant la vitamine ou le minéral, la maladie et le signe qui la trahit.
\end{enumerate}
\textbf{Réflexe} : devant un taux d'hémoglobine, comparer toujours à la norme avant de conclure à une anémie. Ne pas confondre \emph{vitamine} et \emph{sel minéral} (le fer n'est pas une vitamine).
\end{methode}

\begin{exoresolu}[Fatigue et pâleur chez une adolescente]
Une élève de Première A, âgée de 16 ans, se plaint de fatigue permanente, de vertiges et d'essoufflement à l'effort. À l'infirmerie du lycée, on observe une pâleur du visage et des conjonctives. Une analyse de sang prescrite au centre de santé donne un taux d'hémoglobine de \SI{9,5}{g/dL} (norme : \SI{12}{--}\SI{16}{g/dL} chez la femme). Son alimentation est pauvre en viande et en légumes verts.
\begin{enumerate}
\item Calculer le déficit en hémoglobine par rapport à la borne inférieure de la norme.
\item Identifier la maladie et sa cause la plus probable.
\end{enumerate}
\tcblower
\textbf{1. Calcul du déficit en hémoglobine.}

La borne inférieure de la norme est \SI{12}{g/dL}. Le déficit vaut :
\[ \Delta = 12 - 9,5 = \SI{2,5}{g/dL}. \]
Le taux d'hémoglobine de l'élève est donc inférieur de \SI{2,5}{g/dL} à la valeur minimale normale, soit un déficit relatif de :
\[ \frac{2,5}{12} \times 100 \approx 20,8\,\%. \]

\textbf{2. Identification de la maladie.}

Les signes observés — fatigue, vertiges, essoufflement, pâleur des conjonctives — associés à un taux d'hémoglobine abaissé (\SI{9,5}{g/dL} $<$ \SI{12}{g/dL}) caractérisent une \cle{anémie}.

L'hémoglobine est une protéine des globules rouges dont la synthèse nécessite du \cle{fer}. L'énoncé précise que l'alimentation de l'élève est pauvre en viande et en légumes verts, c'est-à-dire pauvre en fer. De plus, une adolescente a des besoins accrus en fer (croissance, pertes menstruelles). La cause la plus probable est donc une \cle{carence martiale} (carence en fer) : il s'agit d'une \textbf{anémie ferriprive}.

\textbf{Conclusion} : l'élève présente une anémie ferriprive (déficit de \SI{2,5}{g/dL} d'hémoglobine), liée à un apport insuffisant en fer ; la correction passe par une alimentation riche en fer (viande, foie, légumes verts comme le folon ou le ndolé) associée à de la vitamine C (goyave, agrumes) qui favorise l'absorption du fer non hémique.
\end{exoresolu}

\courssub{Les maladies de surcharge : l'obésité}

\begin{definition}[Obésité]
Maladie chronique caractérisée par une \cle{accumulation excessive de tissu adipeux} (graisse) présentant un risque pour la santé. Elle résulte d'un \cle{bilan énergétique positif prolongé} : apports énergétiques (alimentation) $>$ dépenses énergétiques (métabolisme de base + activité physique + thermogenèse).
\end{definition}

\begin{propriete}[Diagnostic de l'obésité chez l'adulte : l'IMC]
L'\cle{Indice de Masse Corporelle (IMC)} est l'outil de référence pour le diagnostic de surcharge pondérale chez l'adulte ($\ge$ 18 ans).
\[ \text{IMC} = \frac{m}{t^2} \quad (m \text{ en kg}, \; t \text{ en m}) \quad \text{Unité : } \si{kg/m^2} \]
\begin{center}
\begin{tabularx}{\linewidth}{|l|X|}
\hline
\rowcolor{popblueL} \textbf{IMC (\si{kg/m^2})} & \textbf{Interprétation (OMS)} \\
\hline
$< 18,5$ & Maigreur (dénutrition) \\
\hline
$18,5 \le \text{IMC} < 25$ & Corpulence normale \\
\hline
$25 \le \text{IMC} < 30$ & Surpoids \\
\hline
$30 \le \text{IMC} < 35$ & Obésité modérée (classe I) \\
\hline
$35 \le \text{IMC} < 40$ & Obésité sévère (classe II) \\
\hline
$\ge 40$ & Obésité massive / morbide (classe III) \\
\hline
\end{tabularx}
\end{center}
\textbf{Autre indicateur} : le \cle{tour de taille} (obésité androïde/abdominale : $> \SI{94}{cm}$ homme, $> \SI{80}{cm}$ femme) corrélé au risque métabolique (syndrome métabolique).
\end{propriete}

\begin{propriete}[Complications majeures de l'obésité]
\begin{itemize}
\item \textbf{Métaboliques} : résistance à l'insuline $\rightarrow$ \cle{diabète de type 2}, dyslipidémie (hypertriglycéridémie, HDL bas).
\item \textbf{Cardiovasculaires} : \cle{hypertension artérielle}, athérosclérose, coronaropathie, infarctus du myocarde, accident vasculaire cérébral (AVC).
\item \textbf{Articulaires} : arthrose prématurée (genoux, hanches, rachis) par surcharge mécanique.
\item \textbf{Respiratoires} : syndrome d'apnées du sommeil, hypoventilation.
\item \textbf{Cancéreuses} : sur-risque cancers hormono-dépendants (sein, endomètre, côlon, prostate).
\item \textbf{Psychosociales} : stigmatisation, dépression, baisse de l'estime de soi.
\end{itemize}
\end{propriete}

\begin{attention}[Conditions d'application de l'IMC]
L'IMC ne s'applique qu'aux \textbf{adultes} (18 ans et plus). Chez l'enfant et l'adolescent, on utilise les \textbf{courbes de croissance} (références OMS ou nationales). Il ne distingue pas la masse grasse de la masse musculaire : un sportif très musclé peut avoir un IMC élevé sans excès de graisse (faux positif). Chez la femme enceinte, l'IMC pré-grossesse est utilisé.
\end{attention}

\courssub{Interpréter des données simples : l'IMC}

\begin{methode}[Calculer et interpréter l'indice de masse corporelle]
\begin{enumerate}
\item \textbf{Écrire la formule} : $\text{IMC} = \dfrac{m}{t^{2}}$, avec $m$ la masse en kilogrammes (kg) et $t$ la taille en mètres (m). L'IMC s'exprime en \si{kg/m^2}.
\item \textbf{Convertir la taille en mètres} avant tout calcul (ex. : \SI{165}{cm} $= \SI{1,65}{m}$) et élever au carré.
\item \textbf{Effectuer le calcul} avec les unités, en arrondissant à un chiffre après la virgule.
\item \textbf{Interpréter le résultat} à l'aide de l'échelle de référence (adulte) ci-dessus.
\item \textbf{Conclure} en nommant l'état nutritionnel et en rappelant les risques associés si l'IMC est anormal (diabète de type 2, hypertension artérielle, maladies cardiovasculaires pour l'obésité).
\end{enumerate}
\end{methode}

\begin{exoresolu}[Bilan nutritionnel d'un commerçant de Yaoundé]
M. N., commerçant de 42 ans, mesure \SI{1,70}{m} et pèse \SI{96}{kg}. Sédentaire, il consomme régulièrement des plats très gras et des boissons sucrées.
\begin{enumerate}
\item Calculer son IMC (arrondir au dixième).
\item Interpréter le résultat et citer deux risques pour sa santé.
\end{enumerate}
\tcblower
\textbf{1. Calcul de l'IMC.}

On applique la formule $\text{IMC} = \dfrac{m}{t^{2}}$ avec $m = \SI{96}{kg}$ et $t = \SI{1,70}{m}$.

Étape intermédiaire — calcul du carré de la taille :
\[ t^{2} = 1,70^{2} = 1,70 \times 1,70 = \SI{2,89}{m^{2}}. \]

Calcul de l'IMC :
\[ \text{IMC} = \frac{96}{2,89} \approx 33,2 \ \si{kg/m^2}. \]

\textbf{2. Interprétation.}

On a $30 \leq 33,2 < 35$ : d'après l'échelle de référence, M. N. présente une \cle{obésité modérée} (classe I).

L'obésité est une maladie de \cle{surcharge pondérale} résultant d'un bilan énergétique positif prolongé : les apports (plats gras, boissons sucrées) dépassent durablement les dépenses (sédentarité). L'excès d'énergie est stocké sous forme de triglycérides dans le tissu adipeux.

\textbf{Risques pour la santé} (deux suffisent) :
\begin{itemize}
\item le \cle{diabète de type 2} (résistance des cellules à l'insuline) ;
\item l'\cle{hypertension artérielle} et les maladies cardiovasculaires (infarctus, accident vasculaire cérébral), favorisées par le dépôt de cholestérol sur les parois des artères (athérosclérose).
\end{itemize}

\textbf{Conclusion} : avec un IMC de \SI{33,2}{kg/m^2}, M. N. est en obésité modérée ; il s'expose au diabète de type 2 et aux maladies cardiovasculaires. Une correction de l'alimentation (réduction graisses/sucres, augmentation fibres) et une activité physique régulière sont nécessaires.
\end{exoresolu}

\courssub{Interpréter une courbe de croissance}

\begin{methode}[Lire la courbe de croissance d'un enfant]
La courbe de croissance (carnet de santé) porte la \textbf{masse} (kg) en ordonnée et l'\textbf{âge} (mois ou années) en abscisse. Elle comporte des courbes de référence (z-scores -3, -2, -1, 0, +1, +2, +3) délimitant une zone de croissance normale (entre -2 et +2).
\begin{enumerate}
\item \textbf{Repérer le point} correspondant à l'enfant : âge en abscisse, masse en ordonnée.
\item \textbf{Situer le point par rapport aux courbes de référence} :
\begin{itemize}
\item dans la zone normale (entre -2 et +2 SD) $\Rightarrow$ croissance satisfaisante ;
\item sous la courbe -2 SD $\Rightarrow$ \cle{insuffisance pondérale} (retard de croissance, MPE possible) ;
\item au-dessus de la courbe +2 SD $\Rightarrow$ surcharge pondérale.
\end{itemize}
\item \textbf{Examiner l'évolution de la courbe de l'enfant} au fil des visites (trajectoire) :
\begin{itemize}
\item courbe ascendante et régulière, parallèle aux courbes de référence $\Rightarrow$ bonne croissance ;
\item \textbf{aplatissement} (stagnation de la masse, courbe horizontale) $\Rightarrow$ signal d'alerte majeur ;
\item \textbf{chute} de la courbe (croisement des lignes de référence vers le bas) $\Rightarrow$ danger : malnutrition aiguë ou maladie évolutive.
\end{itemize}
\item \textbf{Conclure} en reliant l'observation à une conduite à tenir (consultation, enquête alimentaire, supplémentation, traitement des infections).
\end{enumerate}
\textbf{Réflexe Probatoire} : une stagnation de la masse pendant deux mois consécutifs chez un nourrisson doit toujours alerter, même si le point reste dans la zone normale. Le \textbf{perimètre brachial} (PB) $< \SI{115}{mm}$ (enfant 6--59 mois) signe une malnutrition aiguë sévère (MAS).
\end{methode}

\begin{exoresolu}[Suivi de la croissance d'un nourrisson]
Le carnet de santé d'un nourrisson suivi au centre de santé de Bafoussam donne les masses suivantes :
\begin{center}
\begin{tabular}{|l|c|c|c|c|c|}
\hline
\rowcolor{popblueL} \textbf{Âge (mois)} & 2 & 4 & 6 & 8 & 10 \\
\hline
\textbf{Masse (kg)} & 5,0 & 6,4 & 7,6 & 7,7 & 7,6 \\
\hline
\end{tabular}
\end{center}
À 10 mois, la limite inférieure de la zone normale (courbe -2 SD) est d'environ \SI{8}{kg}.
\begin{enumerate}
\item Décrire l'évolution de la masse de cet enfant.
\item Interpréter cette évolution et proposer une conduite à tenir.
\end{enumerate}
\tcblower
\textbf{1. Description de l'évolution.}

Décomposons la période en deux phases :
\begin{itemize}
\item \textbf{De 2 à 6 mois} : la masse passe de \SI{5,0}{kg} à \SI{7,6}{kg}, soit un gain de $7,6 - 5,0 = \SI{2,6}{kg}$ en 4 mois (environ \SI{0,65}{kg} par mois) : la courbe est \textbf{ascendante et régulière}, la croissance est satisfaisante.
\item \textbf{De 6 à 10 mois} : la masse passe de \SI{7,6}{kg} à \SI{7,7}{kg} puis redescend à \SI{7,6}{kg} : la courbe \textbf{s'aplatit puis chute}. Le gain sur 4 mois est nul (stagnation puis perte).
\end{itemize}

\textbf{2. Interprétation et conduite à tenir.}

À 10 mois, la masse de l'enfant (\SI{7,6}{kg}) est \textbf{sous la limite inférieure de la zone normale} (\SI{8}{kg}, courbe -2 SD) : l'enfant présente une \cle{insuffisance pondérale} (malnutrition aiguë modérée). L'aplatissement de la courbe à partir de 6 mois — âge habituel de la diversification alimentaire et parfois du sevrage — évoque le début d'une \cle{malnutrition protéino-énergétique}, souvent favorisée par un sevrage mal conduit (bouillie trop diluée, pas assez énergétique ni protéique) ou des infections répétées (diarrhées, paludisme, infections respiratoires).

\textbf{Conduite à tenir} :
\begin{itemize}
\item consultation médicale pour rechercher une cause (infection, parasitose, alimentation insuffisante) et traiter si besoin ;
\item enquête alimentaire auprès de la mère et conseil nutritionnel : poursuite de l'allaitement maternel jusqu'à 2 ans, enrichissement des bouillies infantiles (farine enrichie, arachide, soja, œuf, huile rouge) pour augmenter la densité énergétique et protéique ;
\item pesées régulières de contrôle (mensuelles, voire bimensuelles) pour vérifier la reprise de la croissance (rattrapage) ;
\item supplémentation en vitamine A (tous les 6 mois) et déparasitage selon protocole national.
\end{itemize}

\textbf{Conclusion} : la stagnation puis la chute de la courbe de poids à partir de 6 mois, avec une masse passant sous la limite normale (-2 SD), révèlent un retard de croissance évoquant une MPE débutante ; une prise en charge nutritionnelle et médicale urgente s'impose.
\end{exoresolu}

\courssub{Proposer des mesures de prévention et de correction adaptées}

\begin{methode}[Construire une réponse « prévention et correction »]
\begin{enumerate}
\item \textbf{Identifier d'abord le déséquilibre} (carence ou surcharge) et sa cause : la mesure proposée doit attaquer la cause, pas seulement le symptôme.
\item \textbf{Proposer une correction diététique ciblée} (mesure curative) :
\begin{itemize}
\item carence protéino-énergétique $\Rightarrow$ aliments énergétiques (tubercules, céréales, huile rouge) \emph{et} protéiques (haricots, niébé, arachide, soja, poisson, œufs, lait) ;
\item anémie ferriprive $\Rightarrow$ aliments riches en fer hémique (viande rouge, foie, sang, poisson) + fer non hémique (légumes verts : folon, ndolé, feuilles de manioc, niébé) \textbf{associés à de la vitamine C} (goyave, agrumes, papaye, poivron) qui améliore l'absorption du fer non hémique ; éviter le thé/café au repas ;
\item carence en vitamine A $\Rightarrow$ huile de palme rouge (crue ou peu chauffée), carotte, mangue, patate douce orange, foie ;
\item obésité $\Rightarrow$ réduction des graisses saturées et des sucres rapides, augmentation des fibres (légumes, fruits, céréales complètes), activité physique régulière (\SI{30}{min}/jour marche rapide).
\end{itemize}
\item \textbf{Ajouter les mesures d'accompagnement et de prévention} (mesures préventives) : éducation nutritionnelle des familles (culinaire, hygiène), suivi médical régulier (pesées, analyses), supplémentation médicamenteuse prescrite par un médecin (fer/acide folique femme enceinte, vitamine A enfant 6--59 mois, zinc diarrhée), lutte contre les infections associées (paludisme : MILDA, TPI ; parasitoses : déparasitage ; vaccination ; eau potable/assainissement).
\item \textbf{Rédiger des mesures concrètes et réalistes} dans le contexte local (aliments disponibles sur les marchés camerounais, pratiques culturales), et non des conseils vagues du type « bien manger ».
\end{enumerate}
\textbf{Réflexe Probatoire} : une bonne réponse comporte toujours au moins une mesure \emph{curative} (corriger l'état actuel) et une mesure \emph{préventive} (éviter la récidive / atteindre la population saine).
\end{methode}

\begin{exoresolu}[Plan de prévention dans un village]
Dans un village de la région du Nord, le personnel du centre de santé intégré constate que de nombreux enfants de 1 à 5 ans présentent des signes de kwashiorkor et que plusieurs femmes enceintes souffrent d'anémie. L'alimentation locale repose presque exclusivement sur le mil et le sorgho.

Proposer trois mesures de prévention et de correction adaptées à cette situation, en justifiant chacune.
\tcblower
Analysons la situation : le régime à base de mil et de sorgho est énergétique mais pauvre en protéines de qualité (limitées en lysine) et en fer biodisponible, ce qui explique la fréquence du kwashiorkor (carence protéique) chez les enfants et de l'anémie ferriprive chez les femmes enceintes (besoins accrus).

\textbf{Mesure 1 — Diversifier l'alimentation par les protéines végétales locales (Curatif + Préventif).}

Encourager la culture et la consommation du \cle{niébé}, de l'\cle{arachide} et du \cle{soja}, légumineuses riches en protéines, bon marché et adaptées au climat sahélien de la région du Nord. L'association céréales + légumineuses (mil + niébé, boule de mil + sauce arachide/soja) fournit des protéines complètes (profil d'acides aminés équilibré). \emph{Justification} : elle corrige directement la carence protéique à l'origine du kwashiorkor chez l'enfant et améliore l'apport protéique de la femme enceinte.

\textbf{Mesure 2 — Améliorer les apports en fer et en vitamine C, en particulier chez la femme enceinte (Curatif + Préventif).}

Promouvoir la consommation quotidienne de feuilles vertes locales riches en fer (folon, feuilles de niébé, moringa) et de viande/foie/poisson séché lorsque possible, \textbf{systématiquement accompagnées de fruits riches en vitamine C} (goyave, agrumes, ananas, baobab) qui augmentent l'absorption du fer non hémique végétal ; mise en place effective de la \cle{supplémentation en fer et acide folique} (comprimés \SI{60}{mg} Fe + \SI{400}{\mu g} acide folique) dès le premier trimestre lors des consultations prénatales (CPN). \emph{Justification} : elle corrige et prévient l'anémie ferriprive, dangereuse pour la mère (hémorragies de la délivrance, infection, fatigue) et le fœtus (retard de croissance intra-utérin, prématurité, faible poids de naissance, réserves en fer faibles).

\textbf{Mesure 3 — Éducation nutritionnelle et suivi de la croissance (Préventif).}

Organiser des séances d'éducation à la santé au centre de santé intégré et dans la communauté (causeries éducatives, démonstrations culinaires) sur : conduite du sevrage (âge, fréquence, consistance), enrichissement des bouillies infantiles (farine d'arachide/soja + huile rouge + sucre + lait/œuf), hygiène alimentaire et lavage des mains ; assurer la \cle{pesée mensuelle} systématique des enfants de 0 à 5 ans avec suivi de la courbe de croissance sur le carnet de santé et dépistage par le périmètre brachial (PB). \emph{Justification} : la détection précoce d'un aplatissement ou d'une chute de la courbe (ou PB $< \SI{115}{mm}$) permet d'intervenir \textbf{avant} l'apparition des signes graves de la MPE (kwashiorkor, marasme) et de réduire la mortalité.

\textbf{Conclusion} : la lutte contre les maladies nutritionnelles de ce village repose sur la diversification alimentaire locale (protéines végétales, fer, vitamines), la supplémentation ciblée des femmes enceintes (fer/acide folique) et des enfants (vitamine A, zinc), et l'éducation nutritionnelle assortie d'un suivi régulier de la croissance (pesée mensuelle, PB).
\end{exoresolu}

\courssub{Principes d'une alimentation préventive}

\begin{definition}[Alimentation équilibrée et préventive]
Régime alimentaire qui couvre \textbf{simultanément} les besoins en énergie, en éléments plastiques (protéines) et en micronutriments (vitamines, sels minéraux), sans excès ni carence, tout en limitant les facteurs de risque de maladies chroniques (sel, sucres ajoutés, graisses saturées, alcool). Elle repose sur la \cle{variété}, la \cle{modération} et l'\cle{adaptation} aux besoins individuels.
\end{definition}

\begin{propriete}[Les trois familles d'aliments à associer à chaque repas principal]
\begin{center}
\begin{tabularx}{\linewidth}{|l|X|X|}
\hline
\rowcolor{popblueL} \textbf{Famille} & \textbf{Rôle principal} & \textbf{Exemples locaux (Cameroun)} \\
\hline
\cle{Énergétiques} & Fournir l'énergie (ATP) : glucides complexes + lipides & Riz, maïs, mil, sorgho, manioc, igname, macabo, plantain, pain, pâtes ; huile rouge, arachide, noix de coco, beurre \\
\hline
\cle{De construction (Plastiques)} & Construire/réparer les tissus : protéines, acides aminés essentiels & Viande (bœuf, poulet, brousse), poisson (frais, fumé, séché), œufs, lait, fromage ; \textbf{haricots, niébé, soja, arachide, lentilles} (légumineuses) \\
\hline
\cle{Protecteurs (Régulateurs)} & Réguler le métabolisme, immunité : vitamines, sels minéraux, fibres, eau & \textbf{Fruits} (mangue, papaye, goyave, ananas, banane, agrumes, baobab) ; \textbf{Légumes} (ndolé, folon, okok, feuilles de manioc, tomate, carotte, chou, concombre) ; eau potable \\
\hline
\end{tabularx}
\end{center}
\end{propriete}

\begin{methode}[Vérifier qu'un menu respecte les principes d'une alimentation préventive]
\begin{enumerate}
\item \textbf{Vérifier la présence des trois grandes familles d'aliments} à chaque journée (idéalement à chaque repas principal).
\item \textbf{Vérifier la variété} : varier les sources au sein de chaque famille (céréales \emph{et} tubercules ; protéines animales \emph{et} végétales ; fruits \emph{et} légumes de couleurs variées).
\item \textbf{Vérifier la modération} : limiter sucres rapides (boissons sucrées, bonbons, pâtisseries), graisses saturées (fritures, charcuterie, peau de poulet), sel (cubes de bouillon, sel de table) ; boire \SI{1,5}{--}\SI{2}{L} d'eau potable par jour.
\item \textbf{Adapter les quantités aux besoins} de l'individu (âge, sexe, activité physique, grossesse, allaitement, pathologie).
\item \textbf{Conclure} : un menu est équilibré s'il couvre les besoins énergétiques et plastiques sans excès, et prévient à la fois les carences et les surcharges.
\end{enumerate}
\end{methode}

\begin{exoresolu}[Analyse d'un menu de cantine scolaire]
La cantine d'un lycée de Douala propose le menu suivant : riz blanc, sauce arachide avec viande de bœuf, banane mûre, eau potable.
\begin{enumerate}
\item Montrer que ce menu respecte les principes d'une alimentation préventive.
\item Proposer un complément pour l'améliorer encore.
\end{enumerate}
\tcblower
\textbf{1. Analyse du menu.}

Vérifions la présence des trois familles d'aliments :
\begin{itemize}
\item \textbf{Aliments énergétiques} : le riz blanc (glucides complexes, amidon) et l'huile de la sauce arachide (lipides) fournissent l'énergie nécessaire aux activités scolaires de la journée ;
\item \textbf{Aliments de construction} : la viande de bœuf (protéines animales, fer hémique, zinc, vit B\textsubscript{12}) et l'arachide (protéines végétales, acides gras insaturés) apportent des \cle{protéines} de bonne qualité ;
\item \textbf{Aliments protecteurs} : la banane mûre apporte des \cle{vitamines} (B\textsubscript{6}, C), des sels minéraux (potassium, magnésium) et des fibres ; l'eau potable assure l'\cle{hydratation} sans risque de contamination microbiologique.
\end{itemize}
Le menu comporte donc les trois familles d'aliments, en quantités raisonnables, sans excès de sucres rapides ni de graisses saturées : il est \textbf{varié et équilibré}, conforme aux principes d'une alimentation préventive.

\textbf{2. Complément proposé.}

Le menu est pauvre en \cle{légumes verts} (source majeure de vitamine A, B\textsubscript{9}, C, fer non hémique, fibres, antioxydants). On peut ajouter une portion de légumes verts cuits (ndolé, folon, feuilles de manioc, okok) ou une salade de crudités (tomate, concombre, carotte râpée, chou) avec un filet d'huile et du citron. Cela enrichirait le repas en micronutriments protecteurs, renforçant la prévention des avitaminoses (A, B\textsubscript{9}, C) et de l'anémie, fréquentes chez les adolescents en pleine croissance.

\textbf{Conclusion} : ce menu couvre les besoins énergétiques, plastiques et protecteurs de base ; l'ajout de légumes verts le rendrait encore plus complet sur le plan des vitamines, minéraux et fibres, et favoriserait la satiété.
\end{exoresolu}

\begin{attention}[Les 5 erreurs qui coûtent des points au Probatoire]
\begin{enumerate}
\item \textbf{Confondre kwashiorkor et marasme.} Le kwashiorkor (carence protéique) s'accompagne d'\textbf{œdèmes} et de cheveux décolorés ; le marasme (carence énergétique globale) est un amaigrissement extrême \textbf{sans œdème}. Cite le signe décisif de l'énoncé pour trancher.
\item \textbf{Oublier de mettre la taille en mètres dans le calcul de l'IMC.} Avec $t = 170$ au lieu de $1,70$, on obtient un IMC absurde de $0,003$. Écris toujours la conversion explicite ($170 \text{ cm} = 1,70 \text{ m}$), puis $t^{2}$, avant de diviser. Et n'oublie pas l'unité : \si{kg/m^2}.
\item \textbf{Appliquer l'échelle de l'IMC à un enfant.} Chez l'enfant et l'adolescent, l'IMC seul ne conclut pas : on utilise la \textbf{courbe de croissance} du carnet de santé (références OMS). Mentionner cette condition d'application est souvent valorisé.
\item \textbf{Donner un diagnostic sans justification.} « L'enfant souffre de kwashiorkor » sans citer au moins deux signes de l'énoncé (œdèmes, cheveux roux, déficit pondéral, dermatose) fait perdre la moitié des points : toute identification doit être \textbf{argumentée à partir des données fournies}.
\item \textbf{Proposer des mesures vagues ou inadaptées.} « Il faut bien manger » ou « manger équilibré » ne vaut rien. Propose des aliments \textbf{précis et locaux} (niébé, soja, huile de palme rouge, folon, poisson séché, goyave) ciblant la carence identifiée, et distingue mesure curative (pour le malade) et mesure préventive (pour la population). Attention aussi à ne pas écrire « carence en fer = manque de vitamines » : le fer est un \textbf{sel minéral}, pas une vitamine.
\end{enumerate}
\end{attention}

\newpage

\coursec{Exercices}

\courssub{Je vérifie mes connaissances}

\begin{exercice}[QCM : Maladies nutritionnelles]
Pour chaque question, une seule réponse est exacte. Recopie le numéro et la lettre correspondante.
\begin{enumerate}
    \item La malnutrition protéino-énergétique (MPE) qui se manifeste par un ventre ballonné, des œdèmes, une décoloration des cheveux (cheveux roux) et une dermatose en « peau de serpent » correspond à :
    \begin{itemize}
        \item[A.] Le marasme
        \item[B.] Le kwashiorkor
        \item[C.] L'anémie ferriprive
        \item[D.] Le rachitisme
    \end{itemize}
    \item La carence en vitamine A est responsable de :
    \begin{itemize}
        \item[A.] La pellagre
        \item[B.] Le scorbut
        \item[C.] La xérophtalmie (cécité nocturne, kératomalacie)
        \item[D.] Le goitre
    \end{itemize}
    \item L'indice de masse corporelle (IMC) se calcule par la formule :
    \begin{itemize}
        \item[A.] $\dfrac{\text{Taille (m)}}{\text{Masse (kg)}^2}$
        \item[B.] $\dfrac{\text{Masse (kg)}}{\text{Taille (m)}^2}$
        \item[C.] $\dfrac{\text{Masse (kg)}}{\text{Taille (cm)}^2}$
        \item[D.] $\text{Masse (kg)} \times \text{Taille (m)}$
    \end{itemize}
    \item Selon la classification de l'OMS, un adulte avec un IMC de 32,5 kg/m$^2$ est classé en :
    \begin{itemize}
        \item[A.] Surpoids
        \item[B.] Obésité classe I (modérée)
        \item[C.] Obésité classe II (sévère)
        \item[D.] Obésité classe III (massive ou morbide)
    \end{itemize}
    \item L'anémie ferriprive se caractérise biologiquement par :
    \begin{itemize}
        \item[A.] Une augmentation de l'hémoglobine et du volume globulaire moyen (VGM)
        \item[B.] Une diminution de l'hémoglobine et une augmentation du VGM
        \item[C.] Une diminution de l'hémoglobine et une diminution du VGM (microcytose)
        \item[D.] Une hémoglobine normale mais une baisse du taux de fer sérique seulement
    \end{itemize}
\end{enumerate}
\end{exercice}

\begin{exercice}[Vrai ou Faux ? Justifie chaque réponse]
Indique si l'affirmation est vraie (V) ou fausse (F) et justifie brièvement ton choix en t'appuyant sur les connaissances du cours.
\begin{enumerate}
    \item Le marasme touche préférentiellement l'enfant de 6 à 18 mois sevré brutalement et alimenté uniquement avec des bouillies de maïs ou de manioc très diluées.
    \item La carence en iode est la première cause évitable de retard mental dans le monde ; elle se manifeste cliniquement par le goitre et, chez le fœtus, par le crétinisme.
    \item L'obésité est uniquement due à une consommation excessive de lipides ; les glucides et les protéines n'y jouent aucun rôle.
    \item Sur une courbe de croissance poids-âge, un palier (stagnation) suivi d'une inflexion descendante (chute) est le signe typique d'une malnutrition aiguë installée.
    \item La vitamine D est synthétisée par la peau sous l'action des rayons ultraviolets (UVB) ; une exposition solaire quotidienne de 15 à 20 minutes au visage et aux bras suffit généralement à couvrir les besoins.
\end{enumerate}
\end{exercice}

\begin{exercice}[Définitions et notions clés]
Donne une définition précise et concise pour chacun des termes suivants :
\begin{enumerate}
    \item Malnutrition protéino-énergétique (MPE).
    \item Kwashiorkor.
    \item Marasme.
    \item Avitaminose.
    \item Anémie ferriprive.
    \item Obésité (définition OMS).
    \item IMC (Indice de Masse Corporelle).
    \item Alimentation préventive.
\end{enumerate}
\end{exercice}

\begin{exercice}[Calculs directs d'IMC et classification]
Calcule l'IMC des trois individus suivants et classe-les selon les catégories de l'OMS (Maigreur, Normal, Surpoids, Obésité classe I, II, III). Donne le résultat avec une décimale.
\begin{enumerate}
    \item Individu A : Masse = 68 kg ; Taille = 1,70 m.
    \item Individu B : Masse = 92 kg ; Taille = 1,65 m.
    \item Individu C : Masse = 48 kg ; Taille = 1,60 m.
\end{enumerate}
\end{exercice}

\courssub{Je m'entraîne}

\begin{exercice}[Identification de maladies nutritionnelles à partir de signes cliniques]
Pour chaque cas clinique ci-dessous, identifie la maladie nutritionnelle probable, cite la cause nutritionnelle principale (carence ou excès spécifique) et propose \textbf{deux} mesures de prévention adaptées au contexte camerounais.
\begin{enumerate}
    \item \textbf{Cas 1 :} Enfant de 14 mois, poids = 5,2 kg (très en dessous de la courbe), pas d'œdèmes, peau fripée (« peau de vieillard »), musculature fondue, regard vif, appétit conservé. Cheveux normaux.
    \item \textbf{Cas 2 :} Enfant de 20 mois, poids = 8,5 kg (en dessous de la courbe), présence d'œdèmes bilatéraux des membres inférieurs, ventre ballonné (hépatomégalie), cheveux clairsemés, décolorés (roux), faciles à arracher, dermatose en « peau de serpent » aux fesses et cuisses. Enfant apathique, anorexique.
    \item \textbf{Cas 3 :} Femme enceinte de 28 ans, teint pâle, conjonctives pâles, fatigue intense à l'effort, essoufflement, palpitations. Taux d'hémoglobine = 85 g/L (norme > 110 g/L). VGM bas.
    \item \textbf{Cas 4 :} Adolescent de 16 ans, taille 1,72 m, masse 95 kg. IMC calculé. Antécédents : grignotage fréquent de beignets (puff-puff), boissons sucrées, sedentarité (jeux vidéo). Plainte : douleurs aux genoux, essoufflement rapide.
\end{enumerate}
\end{exercice}

\begin{exercice}[Association : Vitamines, Minéraux, Signes de carence et Sources locales]
Complète le tableau ci-dessous en associant chaque nutriment à son principal signe de carence (ou maladie) et à \textbf{deux} sources alimentaires locales camerounaises riches en ce nutriment.
\begin{center}
\begin{tabularx}{\linewidth}{|l|X|X|}\hline
\textbf{Nutriment} & \textbf{Principal signe de carence / Maladie} & \textbf{Deux sources alimentaires locales (Cameroun)} \\ \hline
Vitamine A & & \\ \hline
Vitamine B1 (Thiamine) & & \\ \hline
Vitamine C & & \\ \hline
Vitamine D & & \\ \hline
Fer & & \\ \hline
Iode & & \\ \hline
\end{tabularx}
\end{center}
\end{exercice}

\begin{exercice}[Construction d'un menu équilibré : Le déjeuner de l'adolescent]
À partir de la liste suivante de 12 aliments couramment consommés au Cameroun, compose un déjeuner équilibré pour un adolescent de 16 ans (besoins énergétiques élevés, croissance). Le menu doit comprendre : une entrée, un plat principal (féculent + protides + légumes), un dessert et une boisson.
\medskip
\textbf{Liste des aliments :} Riz blanc, Ndolé (feuilles de vernonia), Poisson braisé (capitaine), Plantain mûr frit, Haricots blancs (haricots secs cuits), Tomates, Oignons, Huile rouge, Mangue, Orange, Eau, Vin de palme.
\medskip
Rédige ton menu et justifie \textbf{chaque choix} en indiquant le groupe d'aliments auquel il appartient (Céréales/Féculents, Protéines, Légumes/Fruits, Matières grasses, Boissons) et son rôle nutritionnel principal.
\end{exercice}

\begin{exercice}[Interprétation d'une courbe de croissance poids-âge]
On suit la croissance d'un enfant, Pierre, de 0 à 24 mois. Sa courbe de poids (en kg) par rapport à l'âge (en mois) présente l'allure suivante (données simulées) :
\begin{itemize}
    \item 0 mois : 3,2 kg (naissance)
    \item 3 mois : 5,8 kg (croissance régulière, suit le couloir de référence)
    \item 6 mois : 7,5 kg (croissance régulière)
    \item 9 mois : 8,2 kg (début de palier, divergence vers le bas par rapport à la courbe de référence)
    \item 12 mois : 8,5 kg (palier net, stagnation)
    \item 15 mois : 8,3 kg (légère baisse)
    \item 18 mois : 8,0 kg (chute franche, franchissement vers le bas de la courbe -2 SD / 3e percentile)
    \item 21 mois : 7,8 kg
    \item 24 mois : 7,5 kg
\end{itemize}
\begin{enumerate}
    \item Décris l'évolution de la courbe de poids de Pierre en distinguant trois périodes (0-6 mois, 6-12 mois, 12-24 mois).
    \item Quel est l'état nutritionnel probable de Pierre à 24 mois ? Justifie par l'analyse de la courbe.
    \item Quelle période de la vie correspond au début de la divergence (vers 6-9 mois) ? Quelle en est la cause nutritionnelle la plus fréquente au Cameroun ?
    \item Propose une conduite à tenir (prise en charge) en trois étapes pour Pierre.
\end{enumerate}
\end{exercice}

\courssub{Je me prépare au Probatoire}

\begin{exercice}[Cas comparés : Marasme vs Kwashiorkor]
Deux enfants, \textbf{Amadou} (18 mois) et \textbf{Bih} (20 mois), sont admis au Centre de Santé Intégré de leur arrondissement pour malnutrition sévère.
\medskip
\textbf{Observations cliniques :}
\begin{itemize}
    \item \textbf{Amadou :} Poids 5,8 kg. Pas d'œdème. Aspect « petit vieillard » : peau fripée, plis fessiers effacés, musculature très réduite (fonte musculaire). Cheveux normaux. Regard vif, pleure vigoureusement. Température 36,2°C. La mère allaite encore mais donne des bouillies de maïs très liquides 3 fois/jour depuis 6 mois.
    \item \textbf{Bih :} Poids 7,2 kg. Œdèmes bilatéraux des pieds, chevilles, tibias (signe du godet). Ventre globuleux (hépatomégalie). Cheveux clairs, roux, fins, « en drapeau » (se détachent facilement). Dermatose érythémateuse, squameuse aux zones de friction (cuisses, fesses) : « peau de serpent » ou « crazy pavement ». Enfant apathique, triste, anorexie. La mère a arrêté l'allaitement à 12 mois pour la nourrir uniquement de « bobolo » (bâton de manioc) et de sauce pimentée sans protéines animales ni légumes.
\end{itemize}
\begin{enumerate}
    \item Pour chaque enfant, identifie le type de MPE (Marasme ou Kwashiorkor).
    \item Complète le tableau comparatif suivant en t'appuyant sur les observations :
    \begin{center}
    \begin{tabularx}{\linewidth}{|l|X|X|}\hline
    \textbf{Critère} & \textbf{Amadou (Marasme ?)} & \textbf{Bih (Kwashiorkor ?)} \\ \hline
    Cause nutritionnelle dominante & & \\ \hline
    Présence d'œdèmes & & \\ \hline
    Aspect général (peau, muscles, cheveux) & & \\ \hline
    Comportement / Appétit & & \\ \hline
    Âge fréquent de survenue & & \\ \hline
    Alimentation incriminée (d'après le cas) & & \\ \hline
    \end{tabularx}
    \end{center}
    \item Dégage le \textbf{point commun} de prévention primaire pour ces deux formes de MPE.
    \item Cite deux aliments locaux riches en protéines de bonne qualité (valeur biologique élevée) que tu recommanderais d'introduire dans l'alimentation de sevrage pour prévenir ces maladies.
\end{enumerate}
\end{exercice}

\begin{exercice}[Calcul d'IMC, classification OMS et conseils personnalisés]
Trois adultes consultent lors d'une campagne de dépistage des maladies non transmissibles au Centre Hospitalier Régional de Bafoussam. Voici leurs données anthropométriques :
\begin{center}
\begin{tabular}{|c|c|c|c|}\hline
\textbf{Patient} & \textbf{Âge} & \textbf{Masse (kg)} & \textbf{Taille (m)} \\ \hline
M. Fouda & 42 ans & 88 & 1,72 \\ \hline
Mme Atangana & 35 ans & 58 & 1,62 \\ \hline
M. Njoya & 50 ans & 115 & 1,78 \\ \hline
\end{tabular}
\end{center}
\begin{enumerate}
    \item Calcule l'IMC de chaque patient (arrondi à 0,1 kg/m$^2$ près).
    \item Classe chaque patient selon la classification OMS de l'IMC (Maigreur, Normal, Surpoids, Obésité I, II, III).
    \item Pour chaque patient, propose \textbf{un conseil diététique précis} et \textbf{un conseil d'activité physique} adaptés à sa catégorie d'IMC et à son âge.
    \item M. Fouda a une tour de taille de 102 cm. Mme Atangana a un tour de taille de 78 cm. Rappelle la valeur seuil du tour de taille chez l'homme et chez la femme au-delà de laquelle le risque cardiovasculaire est majoré (obésité androïde). Conclure pour ces deux patients.
\end{enumerate}
\end{exercice}

\begin{exercice}[Analyse de document : Évolution pondérale et conduite à tenir]
Le document ci-dessous représente la courbe de croissance poids-âge d'un enfant, \textbf{Kevin}, suivi en consultation de nourrisson sain dans un Centre de Santé de Yaoundé. La courbe de référence (médiane OMS) et les z-scores -2 et -3 sont représentés.

\textit{Description de la courbe de Kevin (garçon) :}
\begin{itemize}
    \item De 0 à 6 mois : La courbe suit parfaitement le couloir de la médiane (z-score 0). Allaitement maternel exclusif déclaré.
    \item À 6 mois : Introduction des aliments de complément (bouillie de maïs + eau + sucre).
    \item De 6 à 12 mois : La courbe marque un \textbf{palier} (stagnation du poids vers 7,8 kg) puis s'infléchit vers le bas, franchissant la ligne -2 SD (z-score -2) vers 10 mois.
    \item À 12 mois : Poids 8,0 kg (z-score -2,5). L'enfant présente une diarrhée aiguë depuis 5 jours.
    \item De 12 à 18 mois : Malgré le traitement de la diarrhée, la courbe reste plate (palier) autour de 8,2 kg, restant sous -2 SD.
    \item À 18 mois : L'enfant est admis pour une pneumonie. Poids 8,1 kg.
    \item De 18 à 24 mois : Après prise en charge nutritionnelle (RUTF - Aliment Thérapeutique Prêt à l'Emploi / Plumpy'nut\textsuperscript{\textregistered} + rééducation nutritionnelle maternelle), la courbe repart en \textbf{rattrapage} (pente raide vers le haut) pour rejoindre le couloir -1 SD à 24 mois (Poids 10,5 kg).
\end{itemize}
\begin{enumerate}
    \item En te basant sur la description, dessine schématiquement (ou décris précisément l'allure) de la courbe de Kevin de 0 à 24 mois par rapport aux courbes de référence (médiane, -2SD, -3SD). Indique les moments clés (6 mois, 12 mois, 18 mois, 24 mois).
    \item Interprète l'évolution de la courbe entre 6 et 12 mois : quel phénomène physiopathologique explique le palier puis la chute ? Relie-le aux pratiques alimentaires décrites.
    \item Quel est le diagnostic nutritionnel de Kevin à 12 mois ? À 18 mois ? Justifie par les critères anthropométriques (z-score POIDS/ÂGE).
    \item Explique le mécanisme du « rattrapage » observé entre 18 et 24 mois. Quelles sont les deux composantes essentielles de la prise en charge qui ont permis ce rattrapage ?
    \item En tant qu'agent de santé, que conseilles-tu à la mère de Kevin \textbf{à 6 mois} (prévention) et \textbf{à 24 mois} (consolidation) pour éviter une rechute ? Donne 3 conseils précis pour chaque âge.
\end{enumerate}
\end{exercice}

\newpage

\coursec{Corrigés des exercices}

\coursec{Leçon NA05 : Lutte contre les maladies nutritionnelles}
\courssub{Corrigés des exercices d'application et d'approfondissement}

\begin{corrige}[Exercice 1 — QCM : Maladies nutritionnelles]
\begin{enumerate}
    \item \textbf{B.} Le \cle{kwashiorkor} (forme œdémateuse de la MPE) se caractérise par la présence d'œdèmes (signe du godet), un ventre ballonné (hépatomégalie par stéatose hépatique), une décoloration des cheveux (roux, « en drapeau »), une dermatose en « peau de serpent » (zones de friction) et une anorexie. Le marasme (A) se manifeste par une fonte musculaire sévère sans œdèmes, un aspect « petit vieillard » et un appétit souvent conservé. L'anémie ferriprive (C) et le rachitisme (D) sont des carences spécifiques (fer, vitamine D/calcium).
    \item \textbf{C.} La carence en \cle{vitamine A} entraîne la \cle{xérophtalmie} (sécheresse conjonctivale, taches de Bitot, cécité crépusculaire, kératomalacie). La pellagre (A) est due à une carence en vitamine B3 (niacine), le scorbut (B) à une carence en vitamine C, le goitre (D) à une carence en iode.
    \item \textbf{B.} La formule exacte de l'\cle{IMC} (Indice de Masse Corporelle) est : $\text{IMC} = \dfrac{\text{Masse (kg)}}{\text{Taille (m)}^2}$. L'unité est $\text{kg/m}^2$.
    \item \textbf{B.} Classification OMS de l'IMC chez l'adulte : Maigreur ($< 18,5$), Normal ($18,5 - 24,9$), Surpoids ($25,0 - 29,9$), \cle{Obésité classe I} ($30,0 - 34,9$), Obésité classe II ($35,0 - 39,9$), Obésité classe III ($\ge 40,0$). Un IMC de $32,5$ se situe dans la classe I (modérée).
    \item \textbf{C.} L'\cle{anémie ferriprive} est une anémie \emph{microcytaire} (VGM diminué $< 80 \text{ fL}$) et \emph{hypochrome} (TCHM diminué), due à un défaut de synthèse de l'hémoglobine par manque de fer. L'hémoglobine est basse. L'augmentation du VGM (B) signe une anémie macrocytaire (carence B9/B12).
\end{enumerate}
\end{corrige}

\begin{corrige}[Exercice 2 — Vrai ou Faux ? Justifie chaque réponse]
\begin{enumerate}
    \item \textbf{V (Vrai).} Le \cle{marasme} survient typiquement entre 6 et 18 mois, lors du sevrage. L'alimentation de relais (bouillies de maïs, manioc, riz très diluées) apporte un volume important mais une densité énergétique et protéique trop faible pour couvrir les besoins élevés de l'enfant en croissance rapide.
    \item \textbf{V (Vrai).} La carence en \cle{iode} est la première cause évitable de retard mental (crétinisme endémique) au monde. Le \cle{goitre} (hypertrophie thyroïdienne compensatrice) est le signe clinique visible. Chez la femme enceinte, la carence entraîne une hypothyroïdie fœtale irréversible (crétinisme : retard mental, surdité, paralysie spastique).
    \item \textbf{F (Faux).} L'\cle{obésité} résulte d'un \textbf{excès d'apport énergétique total} (lipides + glucides + protéines + alcool) par rapport aux dépenses, sur une longue période. Les glucides (surtout sucres ajoutés, boissons sucrées) et les protéines en excès contribuent au bilan énergétique positif et à la lipogenèse hépatique (transformation en acides gras). Dire qu'elle est « uniquement due aux lipides » est une erreur fréquente.
    \item \textbf{V (Vrai).} Sur une courbe \cle{poids-âge}, un \textbf{palier} (stagnation du poids) suivi d'une \textbf{inflexion descendante} (perte de poids) est le signe typique d'une \cle{malnutrition aiguë} (amaigrissement). Cela traduit un catabolisme protéique et lipidique pour fournir de l'énergie (marasme) ou un épisode infectieux aigu sur terrain fragile.
    \item \textbf{V (Vrai).} La \cle{vitamine D} (cholécalciférol) est synthétisée dans la peau (7-déhydrocholestérol) sous l'action des \cle{UVB} solaires. Une exposition de 15 à 20 minutes par jour, visage et bras découverts, entre 10h et 15h, suffit généralement à couvrir les besoins (sauf peaux très foncées, latitudes extrêmes, hiver, pollution, crème solaire).
\end{enumerate}
\end{corrige}

\begin{corrige}[Exercice 3 — Définitions et notions clés]
\begin{enumerate}
    \item \textbf{Malnutrition protéino-énergétique (MPE) :} Ensemble de syndromes cliniques dus à un apport insuffisant en \textbf{protéines} et/ou en \textbf{énergie} (calories) par rapport aux besoins de l'organisme, entraînant une altération de la composition corporelle, de la fonction et de l'issue clinique. Elle concerne surtout l'enfant de 0 à 5 ans.
    \item \textbf{Kwashiorkor :} Forme \emph{œdémateuse} de la MPE, due à une carence \textbf{protéique prédominante} (relative) avec apport énergétique souvent moins réduit. Signes majeurs : œdèmes bilatéraux, dermatose en « peau de serpent », cheveux roux/dépigmentés, hépatomégalie, anorexie, apathie.
    \item \textbf{Marasme :} Forme \emph{sèche} (non œdémateuse) de la MPE, due à une carence \textbf{globale énergétique et protéique}. Signes majeurs : amaigrissement extrême (poids/âge $< -3$ SD), fonte musculaire et graisseuse sévère (aspect « petit vieillard »), plis cutanés distendus, appétit souvent conservé, pas d'œdèmes, pas de dermatose spécifique.
    \item \textbf{Avitaminose :} État pathologique résultant d'une carence prolongée et sévère en \textbf{une ou plusieurs vitamines} essentielles, entraînant des signes cliniques spécifiques (ex : scorbut pour vit. C, béribéri pour vit. B1, xérophtalmie pour vit. A, rachitisme pour vit. D).
    \item \textbf{Anémie ferriprive :} Anémie \textbf{microcytaire} (VGM bas) et \textbf{hypochrome} (TCHM bas) due à l'épuisement des \textbf{stocks de fer} (ferritine basse), limitant la synthèse de l'hémoglobine. Cause principale : apports insuffisants (fer héminique), pertes chroniques (ménorragies, parasitose intestinale), besoins accrus (grossesse, croissance).
    \item \textbf{Obésité (OMS) :} « Accumulation anormale ou excessive de graisse corporelle qui peut nuire à la santé ». Chez l'adulte, définie par un \textbf{IMC $\ge 30 \text{ kg/m}^2$}. Elle est un facteur de risque majeur de maladies non transmissibles (diabète type 2, HTA, cardiopathies, cancers).
    \item \textbf{IMC (Indice de Masse Corporelle) :} Indice anthropométrique simple pour estimer la corpulence d'un adulte : $\text{IMC} = \frac{\text{Masse (kg)}}{\text{Taille}^2 \text{ (m}^2\text{)}}$. Il ne distingue pas masse grasse et masse maigre, ni la répartition du gras (android vs gynoïde).
    \item \textbf{Alimentation préventive :} Alimentation \textbf{variée, équilibrée, diversifiée} et \textbf{adaptée} aux besoins physiologiques (âge, sexe, activité, état physiologique), couvrant les besoins en macro- et micronutriments, limitant les excès (sel, sucres ajoutés, graisses saturées, alcool) et respectant l'hygiène alimentaire, pour prévenir les carences et les maladies chroniques.
\end{enumerate}
\end{corrige}

\begin{corrige}[Exercice 4 — Calculs directs d'IMC et classification]
\begin{enumerate}
    \item \textbf{Individu A :} $\text{IMC} = \dfrac{68}{1,70^2} = \dfrac{68}{2,89} = 23,529 \approx \textbf{23,5 kg/m}^2$. \quad Classification : \textbf{Normal} ($18,5 - 24,9$).
    \item \textbf{Individu B :} $\text{IMC} = \dfrac{92}{1,65^2} = \dfrac{92}{2,7225} = 33,792 \approx \textbf{33,8 kg/m}^2$. \quad Classification : \textbf{Obésité classe I (modérée)} ($30,0 - 34,9$).
    \item \textbf{Individu C :} $\text{IMC} = \dfrac{48}{1,60^2} = \dfrac{48}{2,56} = 18,75 \approx \textbf{18,8 kg/m}^2$. \quad Classification : \textbf{Normal} ($18,5 - 24,9$). \textit{Note : limite inférieure de la normale (18,5), surveillance recommandée.}
\end{enumerate}
\end{corrige}

\begin{corrige}[Exercice 5 — Identification de maladies nutritionnelles à partir de signes cliniques]
\begin{enumerate}
    \item \textbf{Cas 1 :} \textbf{Marasme} (MPE sèche). \textbf{Cause :} Carence globale énergétique et protéique (bouillies trop diluées, sevrage précoce/inadapté, infections à répétition). \textbf{Prévention (Cameroun) :} 1) Allaitement maternel exclusif jusqu'à 6 mois + poursuite jusqu'à 2 ans. 2) Diversification alimentaire à 6 mois avec bouillies \emph{épaisses et enrichies} (arachide/lait/huile/poisson/œuf + légumes) données 3 à 4 fois/jour + éducation de la mère sur l'hygiène alimentaire.
    \item \textbf{Cas 2 :} \textbf{Kwashiorkor} (MPE œdémateuse). \textbf{Cause :} Carence \textbf{protéique prédominante} (apport énergétique suffisant par le manioc/bobolo, mais quasi absence de protéines de bonne qualité). \textbf{Prévention (Cameroun) :} 1) Introduction systématique de \textbf{protéines animales} (poisson sec/poudre de poisson, œufs, lait) ou végétales associées (légumineuses + céréales) dans l'alimentation de complément dès 6 mois. 2) Éducation nutritionnelle sur la valeur des aliments « constructeurs » (protéines) et lutte contre les croyances alimentaires défavorables (interdits protéiques).
    \item \textbf{Cas 3 :} \textbf{Anémie ferriprive de la grossesse}. \textbf{Cause :} Carence en \textbf{fer} (besoins accrus : expansion volémie, fœtus, placenta + pertes menstruelles antérieures + apport alimentaire insuffisant en fer héminique). \textbf{Prévention (Cameroun) :} 1) \textbf{Supplémentation systématique} en fer (60 mg) + acide folique (400 µg) dès le 1er trimestre (politique nationale CPN). 2) Alimentation riche en fer héminique (viande rouge, foie, rognons, poisson, sang) + sources de vitamine C (agrumes, goyave, tomate) pour favoriser l'absorption du fer non héminique (légumineuses, feuilles vertes).
    \item \textbf{Cas 4 :} \textbf{Obésité} (IMC = $95 / 1,72^2 = 32,1 \text{ kg/m}^2$ $\rightarrow$ Obésité classe I). \textbf{Cause :} Excès énergétique chronique (grignotage \emph{puff-puff} riches en glucides/lipides, boissons sucrées) + \textbf{sédentarité} (jeux vidéo). \textbf{Prévention (Cameroun) :} 1) Rééquilibrage alimentaire : réduire drastiquement beignets, boissons sucrées, fritures ; augmenter légumes (ndolé, folong, tomates), fruits, protéines maigres (poisson grillé, poulet), féculents complets ; cuisson vapeur/braise sans excès d'huile. 2) Activité physique quotidienne $\ge 60$ min (marche, sport, danse) + limitation temps d'écran $< 2$ h/jour ; prise en charge multidisciplinaire si besoin.
\end{enumerate}
\end{corrige}

\begin{corrige}[Exercice 6 — Association : Vitamines, Minéraux, Signes de carence et Sources locales]
\begin{center}
\begin{tabularx}{\linewidth}{|l|X|X|}\hline
\rowcolor{popblueL}
\textbf{Nutriment} & \textbf{Principal signe de carence / Maladie} & \textbf{Deux sources alimentaires locales (Cameroun)} \\ \hline
Vitamine A & Xérophtalmie (cécité crépusculaire, taches de Bitot, kératomalacie), baisse immunité, retard croissance & \textbf{Huile rouge} (très riche), \textbf{Feuilles vertes} (ndolé, folong, épinards, amarante), mangue, papaye, carotte, foie, œufs \\ \hline
Vitamine B1 (Thiamine) & \textbf{Béribéri} : sec (polyneuropathie, paralysies) ou humide (insuffisance cardiaque, œdèmes), troubles neurologiques (encéphalopathie de Wernicke) & \textbf{Céréales complètes} (maïs grain, mil, sorgho, riz non décortiqué), \textbf{Légumineuses} (haricots, niébé, arachide, soja), porc, foie, levure de bière \\ \hline
Vitamine C & \textbf{Scorbut} : saignements gingivaux, ecchymoses, cicatrisation lente, fatigue, douleurs ostéo-articulaires, anémie & \textbf{Goyave} (très riche), \textbf{Agrumes} (orange, mandarine, citron, pamplemousse), ananas, papaye, poivron, tomate, feuilles vertes crues \\ \hline
Vitamine D & \textbf{Rachitisme} (enfant : déformations osseuses, retard marche/dentition) / \textbf{Ostéomalacie} (adulte : douleurs osseuses, fractures) & \textbf{Poissons gras} (capitaine, sardine, maquereau, hareng), \textbf{Huile de foie de morue}, jaune d'œuf, champignons séchés au soleil, \textbf{Lait/yaourts enrichis} (source majeure = soleil UVB) \\ \hline
Fer & \textbf{Anémie ferriprive} (pâleur, fatigue, essoufflement, VGM bas, ferritine basse), baisse performances cognitives & \textbf{Viande rouge} (bœuf, chevreuil), \textbf{Abats} (foie, rein, cœur), \textbf{Poisson} (sardine, capitaine), Légumineuses (haricots, soja, lentilles) + Vit C, Feuilles vertes (ndolé) \\ \hline
Iode & \textbf{Goitre} (hypertrophie thyroïde), \textbf{Crétinisme} (fœtus/enfant : retard mental irréversible), Hypothyroïdie, Avortements & \textbf{Sel iodé} (obligatoire au Cameroun, source n°1), \textbf{Poissons et fruits de mer} (crevettes, crabes), algues, lait, œufs \\ \hline
\end{tabularx}
\end{center}
\end{corrige}

\begin{corrige}[Exercice 7 — Construction d'un menu équilibré : Le déjeuner de l'adolescent]
\textbf{Menu proposé :}
\begin{itemize}
    \item \textbf{Entrée :} \textbf{Salade de tomates et oignons} (assaisonnée avec un filet d'huile rouge et sel iodé).
    \item \textbf{Plat principal :} \textbf{Riz blanc} + \textbf{Ndolé cuisiné au poisson braisé (capitaine) et huile rouge} (avec crevettes séchées si possible).
    \item \textbf{Dessert :} \textbf{Mangue} + \textbf{Orange} (salade de fruits).
    \item \textbf{Boisson :} \textbf{Eau potable} (1 grand verre).
\end{itemize}

\textbf{Justification nutritionnelle :}
\begin{center}
\begin{tabularx}{\linewidth}{|l|l|X|}\hline
\rowcolor{popblueL}
\textbf{Élément du menu} & \textbf{Groupe d'aliments} & \textbf{Rôle nutritionnel principal} \\ \hline
Salade tomates/oignons & Légumes (crus) & Apport en \textbf{vitamine C} (absorption fer), fibres, eau, antioxydants (lycopène), faible densité énergétique. \\ \hline
Riz blanc & Céréales / Féculents & Source principale d'\textbf{énergie} (glucides complexes/amidon), satiété, quelques protéines végétales, vitamines B. \\ \hline
Ndolé (feuilles) & Légumes (verts foncés) & Riche en \textbf{vitamine A} (bêta-carotène), \textbf{fer} non héminique, \textbf{calcium}, fibres, folates (vit B9). \\ \hline
Poisson braisé (capitaine) & Protéines (animales) & \textbf{Protéines de haute valeur biologique} (acides aminés essentiels), \textbf{fer héminique} (bien absorbé), \textbf{zinc}, \textbf{vitamine B12}, oméga-3. \\ \hline
Huile rouge (cuisson) & Matières grasses & Source concentrée d'\textbf{énergie}, \textbf{vitamine A} (bêta-carotène), \textbf{vitamine E}, acides gras essentiels. Améliore la palatabilité et l'absorption vit A. \\ \hline
Mangue + Orange & Fruits & \textbf{Vitamine C} majeure (orange), \textbf{vitamine A} (mangue), fibres, eau, potassium, antioxydants. Dessert sain en remplacement des sucreries. \\ \hline
Eau & Boissons & \textbf{Hydratation} indispensable, transport nutriments, élimination déchets, thermorégulation. Seule boisson indispensable. \\ \hline
\end{tabularx}
\end{center}
\textit{Note : Le plantain frit, les haricots blancs et le vin de palme ne sont pas retenus. Plantain frit = excès lipides (friture) ; haricots = redondance protéines végétales (poisson déjà présent) + risque ballonnements ; vin de palme = alcool, déconseillé à l'adolescent.}
\end{corrige}

\begin{corrige}[Exercice 8 — Interprétation d'une courbe de croissance poids-âge]
\begin{enumerate}
    \item \textbf{Évolution en trois périodes :}
    \begin{itemize}
        \item \textbf{0–6 mois :} \textbf{Croissance normale (satisfaisante).} La courbe suit le couloir de la médiane (z-score $\approx 0$). Poids : 3,2 $\rightarrow$ 7,5 kg. Gain $\approx$ 4,3 kg. Correspond à l'allaitement maternel exclusif (AME) bien conduit.
        \item \textbf{6–12 mois :} \textbf{Divergence / Palier / Début de malnutrition.} À 9 mois, la courbe marque un palier (8,2 kg) puis stagne à 12 mois (8,5 kg). Elle s'éloigne vers le bas de la courbe de référence. Gain sur 6 mois = seulement 1 kg (au lieu de $\approx$ 2 kg attendus). Signe d'une \textbf{malnutrition aiguë modérée} naissante.
        \item \textbf{12–24 mois :} \textbf{Aggravation / Malnutrition aiguë sévère installée.} Chute franche du poids : 8,5 $\rightarrow$ 7,5 kg (perte de 1 kg en 12 mois). Franchissement de la ligne -2 SD (z-score -2) vers 18 mois, poursuite sous -2 SD. À 24 mois, poids 7,5 kg (z-score $< -3$ SD probable). Tableau de \textbf{marasme sévère}.
    \end{itemize}
    \item \textbf{État nutritionnel à 24 mois :} \textbf{Malnutrition aiguë sévère (Marasme)}. Justification : Poids/âge $< -3$ SD (ou franchissement net de -2 SD avec perte de poids continue), absence d'œdèmes (données cliniques compatibles : amaigrissement extrême, poids 7,5 kg à 24 mois vs médiane $\approx$ 12 kg).
    \item \textbf{Début divergence (6–9 mois) :} Correspond à la \textbf{période de sevrage / introduction des aliments de complément (6–24 mois)}. \textbf{Cause nutritionnelle la plus fréquente au Cameroun :} Aliments de complément \textbf{trop dilués} (bouillies de maïs/manioc très liquides), \textbf{faible densité énergétique et protéique}, \textbf{faible fréquence des repas} ($< 3-4$/jour), \textbf{hygiène précaire} (eau, mains, ustensiles) $\rightarrow$ diarrhées à répétition $\rightarrow$ malabsorption + pertes nutritives.
    \item \textbf{Conduite à tenir (3 étapes) :}
    \begin{enumerate}
        \item \textbf{Prise en charge médicale urgente (Phase de stabilisation) :} Traiter les complications (infections : antibiotiques adaptés ; déshydratation : réhydratation orale/IV ; hypoglycémie : sucre/sérum glucosé ; hypothermie : réchauffement ; anémie sévère : transfusion si Hb $< 4$ g/dL ou $< 6$ g/dL avec détresse). Vitamine A (200 000 UI si $> 12$ mois), Zinc, Folates, Vermifuge (après stabilisation).
        \item \textbf{Réhabilitation nutritionnelle (Phase de transition puis réhabilitation) :} \textbf{Aliment Thérapeutique Prêt à l'Emploi (ATPE / RUTF type Plumpy'nut\textsuperscript{\textregistered})} à domicile (si appétit bon, pas de complication) : $\approx$ 200 kcal/kg/jour (ex: 2-3 sachets/jour pour 7,5 kg). Éducation de la mère : alimentation familiale enrichie (bouillie épaisse + huile + arachide/lait/poisson/œuf + légumes) 5-6 fois/jour, hygiène, poursuite allaitement.
        \item \textbf{Suivi rapproché et prévention rechute :} Pesée hebdomadaire puis bi-mensuelle au CSI jusqu'à récupération (Poids/âge $> -1$ SD ou $\ge 85\%$ médiane). Supplémentation Vit A (tous les 6 mois), Fer/Folates, Vaccination à jour (Rougeole prioritaire), Moustiquaire imprégnée, Éducation nutritionnelle continue (diversification, fréquence, hygiène, prise en charge précoce maladies).
    \end{enumerate}
\end{enumerate}
\end{corrige}

\begin{corrige}[Exercice 9 — Cas comparés : Marasme vs Kwashiorkor]
\begin{enumerate}
    \item \textbf{Amadou : Marasme.} \textbf{Bih : Kwashiorkor.}
    \item \textbf{Tableau comparatif :}
    \begin{center}
    \begin{tabularx}{\linewidth}{|l|X|X|}\hline
    \rowcolor{popblueL}
    \textbf{Critère} & \textbf{Amadou (Marasme)} & \textbf{Bih (Kwashiorkor)} \\ \hline
    Cause nutritionnelle dominante & Carence \textbf{globale énergétique et protéique} (apport total insuffisant). Bouillies très liquides = faible densité calorique. & Carence \textbf{protéique prédominante} (relative). Apport énergétique peut-être correct (manioc/bobolo) mais quasi absence de protéines de bonne qualité. \\ \hline
    Présence d'œdèmes & \textbf{Absents} (pas de signe du godet). & \textbf{Présents} : œdèmes bilatéraux des membres inférieurs (pieds, chevilles, tibias), signe du godet positif. Ventre globuleux (hépatomégalie + ascite possible). \\ \hline
    Aspect général (peau, muscles, cheveux) & Peau fripée (« peau de vieillard »), plis fessiers effacés, \textbf{fonte musculaire sévère}, cheveux \textbf{normaux} (pas de dépigmentation). & Peau : dermatose érythémato-squameuse « peau de serpent » / « crazy pavement » (zones friction). Cheveux : \textbf{clairs, roux, fins, « en drapeau »} (se détachent facilement). Muscles : moins fondus qu'au marasme (masquée par œdèmes). \\ \hline
    Comportement / Appétit & \textbf{Regard vif, pleure vigoureusement, appétit conservé} (réclame à manger). & \textbf{Apathique, triste, anorexie} (refuse de manger), irritable. \\ \hline
    Âge fréquent de survenue & \textbf{6–18 mois} (pic vers 1 an), période sevrage. & \textbf{18–36 mois} (pic vers 2 ans), post-sevrage. \\ \hline
    Alimentation incriminée (d'après le cas) & Allaitement maternel + \textbf{bouillies de maïs très liquides} (3 fois/jour) : faible énergie, faible protéines, grand volume. & Sevrage à 12 mois. \textbf{Bobolo (manioc) + sauce pimentée seule} : énergie (glucides) mais \textbf{zéro protéines animales, zéro légumes}. \\ \hline
    \end{tabularx}
    \end{center}
    \item \textbf{Point commun de prévention primaire :} \textbf{Allaitement maternel exclusif jusqu'à 6 mois} + \textbf{Diversification alimentaire adaptée dès 6 mois} (aliments de complément \textbf{épaisses, enrichies en protéines de bonne qualité, en énergie, en micronutriments}, données à \textbf{fréquence suffisante} 3-4 fois/jour + collations, dans des conditions d'\textbf{hygiène rigoureuse}) + \textbf{Poursuite de l'allaitement jusqu'à 2 ans et au-delà}.
    \item \textbf{Deux aliments locaux riches en protéines de haute valeur biologique (VB élevée) pour le sevrage :}
    \begin{enumerate}
        \item \textbf{Poisson (fraîchement pêché, séché, fumé ou en poudre) :} Capitaine, sardine, carpe, « kanga ». Protéines complètes, fer héminique, zinc, calcium (si mangé avec arêtes), oméga-3. Facile à incorporer dans les bouillies (poudre de poisson).
        \item \textbf{Œuf (de poule, de caille) :} Protéine de référence (score acides aminés = 100), fer, zinc, vitamine A, D, B12, choline. Un œuf/jour ou 3-4/semaine dans la bouillie ou en omelette/mouille. Alternative végétale de qualité : \textbf{Soja} (lait de soja, farine de soja toastée) bien préparé pour détruire facteurs antinutritionnels.
    \end{enumerate}
\end{enumerate}
\end{corrige}

\begin{corrige}[Exercice 10 — Calcul d'IMC, classification OMS et conseils personnalisés]
\begin{enumerate}
    \item \textbf{Calculs d'IMC (arrondi 0,1) :}
    \begin{align*}
        \text{M. Fouda} &: \frac{88}{1,72^2} = \frac{88}{2,9584} = 29,74 \approx \textbf{29,7 kg/m}^2 \\
        \text{Mme Atangana} &: \frac{58}{1,62^2} = \frac{58}{2,6244} = 22,10 \approx \textbf{22,1 kg/m}^2 \\
        \text{M. Njoya} &: \frac{115}{1,78^2} = \frac{115}{3,1684} = 36,30 \approx \textbf{36,3 kg/m}^2
    \end{align*}
    \item \textbf{Classification OMS :}
    \begin{itemize}
        \item M. Fouda (29,7) : \textbf{Surpoids} ($25,0 - 29,9$). \textit{Note : limite supérieure, risque bascule obésité.}
        \item Mme Atangana (22,1) : \textbf{Normal} ($18,5 - 24,9$).
        \item M. Njoya (36,3) : \textbf{Obésité classe II (sévère)} ($35,0 - 39,9$).
    \end{itemize}
    \item \textbf{Conseils personnalisés :}
    \begin{center}
    \begin{tabularx}{\linewidth}{|l|X|X|}\hline
    \rowcolor{popblueL}
    \textbf{Patient} & \textbf{Conseil diététique précis} & \textbf{Conseil activité physique} \\ \hline
    \textbf{M. Fouda} (Surpoids, 42 ans) & Réduire les portions (assiette : 1/2 légumes, 1/4 féculents complets, 1/4 protéines maigres). Supprimer boissons sucrées, alcool, fritures, charcuteries. Cuisson vapeur, braise, papillote. 3 repas + 1 collation fruit/yaourt. & \textbf{Marche rapide 30–45 min / jour, 5 j/sem} (ou équivalent : vélo, natation). Intégrer marche quotidienne (transport, pauses). Renforcement musculaire 2 j/sem (gainage, squats). \\ \hline
    \textbf{Mme Atangana} (Normal, 35 ans) & Maintenir alimentation variée, équilibrée, diversifiée (fruits/légumes à chaque repas, protéines variées, féculents complets, huiles végétales). Limiter sel, sucres ajoutés, gras saturés. Hydratation eau 1,5–2 L/j. & Maintenir une \textbf{activité régulière} : marche 30 min/j, sport loisir (danse, fitness, ballon) 2-3 fois/sem. Éviter sédentarité prolongée (lever toutes les heures). \\ \hline
    \textbf{M. Njoya} (Obésité II, 50 ans) & \textbf{Prise en charge diététique encadrée} (nutritionniste). Hypocalorique modérée (-500 à -600 kcal/j). Fractionnement (3 repas + 1-2 collations). Riches fibres (légumes, complets), protéines maigres (satiété), index glycémique bas. Stop sucres/gras cachés/alcool. & \textbf{Avis médical préalable} (bilan cardio-vasculaire, articulaire). Reprise \textbf{très progressive} : marche 15-20 min/j $\rightarrow$ 45 min. Activités portées (vélo, natation, aquagym) pour ménager articulations (genoux). Renforcement musculaire léger encadré. Objectif : 150–300 min/sem d'endurance modérée. \\ \hline
    \end{tabularx}
    \end{center}
    \item \textbf{Tour de taille (obésité androïde / risque cardiovasculaire) :}
    \begin{itemize}
        \item \textbf{Seuils OMS/IDF :} \textbf{Homme :} Risque accru $\ge 94$ cm, \textbf{Risque très élevé $\ge 102$ cm}. \textbf{Femme :} Risque accru $\ge 80$ cm, \textbf{Risque très élevé $\ge 88$ cm}.
        \item \textbf{M. Fouda :} Tour de taille = \textbf{102 cm} $\rightarrow$ \textbf{Seuil de risque très élevé atteint}. Obésité androïde (distribution abdominale) confirmée. Risque cardiovasculaire (HTA, diabète, dyslipidémie, syndrome métabolique) \textbf{majoré} malgré IMC « seulement » surpoids.
        \item \textbf{Mme Atangana :} Tour de taille = \textbf{78 cm} $\rightarrow$ \textbf{En dessous du seuil de risque accru (80 cm)}. Distribution gynoïde ou normale. Risque cardiovasculaire lié à l'adiposité abdominale \textbf{non majoré} par le tour de taille.
    \end{itemize}
    \textbf{Conclusion :} L'IMC seul sous-estime le risque chez M. Fouda (obésité androïde « cachée » par une masse musculaire possible ou une répartition centrale). Le tour de taille est un marqueur complémentaire indispensable du risque métabolique.
\end{enumerate}
\end{corrige}

\begin{corrige}[Exercice 11 — Analyse de document : Évolution pondérale et conduite à tenir (Kevin)]
\begin{enumerate}
    \item \textbf{Description schématique de la courbe de Kevin (0–24 mois) :}
    \begin{itemize}
        \item \textbf{0–6 mois :} Courbe \textbf{superposable à la médiane (z-score 0)}. Croissance harmonieuse.
        \item \textbf{6 mois :} Point d'inflexion. Introduction bouillie maïs/eau/sucre.
        \item \textbf{6–10 mois :} \textbf{Palier} (stagnation $\approx$ 7,8 kg). Courbe s'aplatit, s'éloigne vers le bas de la médiane.
        \item \textbf{10 mois :} \textbf{Franchissement de la ligne -2 SD (z-score -2)} vers le bas.
        \item \textbf{12 mois :} Point bas : \textbf{8,0 kg (z-score -2,5)}. Courbe sous -2 SD. Épisode diarrhée aiguë.
        \item \textbf{12–18 mois :} \textbf{Palier prolongé} sous -2 SD (poids stable $\approx$ 8,1–8,2 kg). Pas de rattrapage malgré traitement diarrhée. Épisode pneumonie à 18 mois (8,1 kg).
        \item \textbf{18–24 mois :} \textbf{Rattrapage (catch-up growth) :} Pente \textbf{raide, ascendante}, plus forte que la médiane. Courbe remonte pour \textbf{rejoindre le couloir -1 SD à 24 mois} (Poids 10,5 kg).
    \end{itemize}
    \item \textbf{Interprétation 6–12 mois (Palier + Chute) :} Phénomène : \textbf{Malnutrition aiguë modérée puis sévère (Marasme)}. Physiopathologie : \textbf{Catabolisme protéique et lipidique} pour couvrir les besoins énergétiques de base (métabolisme de base + activité + coût de la croissance) non couverts par l'apport alimentaire.
    \textbf{Lien avec pratiques alimentaires :} Introduction tardive/inadéquate d'aliments de complément à 6 mois : \textbf{bouillie de maïs + eau + sucre} = \textbf{faible densité énergétique} ($\ll 0,8$ kcal/g visé), \textbf{faible densité protéique}, \textbf{faible densité micronutriments}, \textbf{volume gastrique limitant} (estomac petit). Apport calorique/jour $<$ Besoins. Pas de protéines animales, pas de légumes, pas d'huile. Hygiène précaire (eau, mains) $\rightarrow$ diarrhée $\rightarrow$ malabsorption + pertes protéiques/électrolytes + anorexie. Cercle vicieux malnutrition-infection.
    \item \textbf{Diagnostic nutritionnel :}
    \begin{itemize}
        \item \textbf{À 12 mois :} \textbf{Malnutrition Aiguë Modérée (MAM)}. Critère OMS : Poids/Âge (P/A) z-score $= -2,5$ (entre -3 et -2 SD). \textit{Nuance :} Avec diarrhée aiguë et anorexie, prise en charge souvent alignée sur MAS (hospitalisation si complications).
        \item \textbf{À 18 mois :} \textbf{Malnutrition Aiguë Modérée (MAM) persistante / non résolue}. P/A reste $< -2$ SD (poids 8,1 kg, âge 18 mois, médiane $\approx$ 11 kg $\rightarrow$ z-score $\approx -2,5$ à -3). Aggravation par épisode infectieux (pneumonie) sur terrain fragilisé. Absence de rattrapage = échec de la prise en charge ambulatoire simple.
    \end{itemize}
    \item \textbf{Mécanisme du rattrapage (18–24 mois) :} \textbf{Anabolisme intense} permis par un apport nutritionnel \textbf{massif, dense et de haute qualité} (RUTF/Plumpy'nut\textsuperscript{\textregistered} : $\approx$ 5,3 kcal/g, 25-30\% protéines lait/arachide, vitamines/minéraux) $\rightarrow$ couverture large des besoins (énergie + protéines + micronutriments) + \textbf{résolution des infections} (pneumonie traitée) $\rightarrow$ cessation du catabolisme, restauration de la sensibilité à l'insuline/IGF-1, synthèse protéique et dépôt adipeux rapides. Pente de croissance $>$ pente de référence.
    \textbf{Deux composantes essentielles de la prise en charge :}
    \begin{enumerate}
        \item \textbf{Aliment Thérapeutique Prêt à l'Emploi (ATPE / RUTF)} à domicile (approche CTC/PCIME : Prise en Charge Intégrée de la Malnutrition Aiguë en Milieu Communautaire) : apport nutritionnel optimal, sans eau (pas contamination), appétent, géré par la mère.
        \item \textbf{Prise en charge médicale des complications} (antibiotiques pour pneumonie, vitamine A, zinc, vermifuge, palu) \textbf{+ Éducation nutritionnelle / Rééducation maternelle} (préparation bouillies enrichies, fréquence, hygiène, stimulation psycho-motrice) pour assurer la transition vers l'alimentation familiale et prévenir la rechute.
    \end{enumerate}
    \item \textbf{Conseils à la mère :}
    \begin{itemize}
        \item \textbf{À 6 mois (Prévention - moment critique de la divergence) :}
        \begin{enumerate}
            \item \textbf{Introduire les aliments de complément à 6 mois précis} (ni avant, ni après), en \textbf{maintenant l'allaitement maternel à la demande} (au moins 8-10 tétées/24h).
            \item \textbf{Donner des bouillies ÉPAISSES et ENRICHIES :} Farine de maïs/manioc/soja/arachide + \textbf{huile rouge} (1 c.à.c/bol) + \textbf{poudre de poisson/crevettes/œuf/lait} + \textbf{légumes verts finement mixés} (ndolé, folong). Densité énergétique $\ge 1$ kcal/g (idéal 1,5). Volume adapté (200-250 ml/bol).
            \item \textbf{Fréquence :} 3 repas + 1-2 collations nutritives/jour. \textbf{Hygiène stricte :} Lavage mains au savon (mère + enfant) avant repas, eau potable (bouillie/filtrée/chlorée), ustensiles propres, consommation immédiate (pas de conservation à température ambiante).
        \end{enumerate}
        \item \textbf{À 24 mois (Consolidation - prévention rechute) :}
        \begin{enumerate}
            \item \textbf{Alimentation familiale diversifiée et enrichie :} Kevin mange comme la famille (morceaux tendres) mais \textbf{assiette enrichie} : 1/2 légumes (ndolé, tomates, carottes), 1/4 féculents (riz, plantain, igname, patate), 1/4 protéines (poisson, viande, œufs, légumineuses) + huile rouge. 3 repas principaux + 2 collations (fruit, yaourt, bouillie enrichie, arachide grillée).
            \item \textbf{Suivi de croissance mensuel au CSI} jusqu'à 3 ans (puis trimestriel) : Pesée, traçage sur carte de croissance, détection précoce de tout palier ou chute (z-score P/A ou P/T). Rattrapage statural (taille/âge) à surveiller aussi.
            \item \textbf{Mesures de santé publique systématiques :} Vaccination complète (Rougeole 2e dose, Fièvre jaune, Méningite, etc.), \textbf{Vitamine A 200 000 UI tous les 6 mois} (jusqu'à 5 ans), \textbf{Vermifuge (Mébendazole/Albendazole) tous les 6 mois} (dès 12 mois), \textbf{Moustiquaire imprégnée à longue durée d'action (MILD)} chaque nuit, traitement précoce de toute fièvre/diarrhée (ORS + Zinc + alimentation continue).
        \end{enumerate}
    \end{itemize}
\end{enumerate}
\end{corrige}

\newpage

\coursec{Fiche bilan}

\coursec{NA05 : Lutte contre les maladies nutritionnelles}
\courssub{Synthèse illustrée et récapitulatifs}

\begin{popfigure}
\centering
\begin{tikzpicture}[
  mindmap,
  concept color=popblue,
  level 1/.append style={level distance=4.5cm, sibling angle=60},
  level 2/.append style={level distance=3cm, sibling angle=45},
  every node/.style={font=\small},
  text=white,
  grow cyclic
]
\node[concept, minimum size=2.2cm, align=center]{Maladies\\ Nutritionnelles}
  child[concept color=poporange] {
    node[concept, minimum size=1.8cm, align=center]{Maladies\\ de carence}
    child { node[concept color=poporange!70!white, minimum size=1.4cm, align=center]{MPE\\ (Kwashiorkor, Marasme)} }
    child { node[concept color=poporange!70!white, minimum size=1.4cm, align=center]{Avitaminoses\\ (A, D, B1, C, PP)} }
    child { node[concept color=poporange!70!white, minimum size=1.4cm, align=center]{Anémie\\ (Fer, B9, B12)} }
  }
  child[concept color=poppink] {
    node[concept, minimum size=1.8cm, align=center]{Maladies\\ de surcharge}
    child { node[concept color=poppink!70!white, minimum size=1.4cm, align=center]{Obésité\\ (IMC $\ge$ 30)} }
    child { node[concept color=poppink!70!white, minimum size=1.4cm, align=center]{Surpoids\\ (25 $\le$ IMC < 30)} }
  }
  child[concept color=poppurple] {
    node[concept, minimum size=1.8cm] {Causes}
    child { node[concept color=poppurple!70!white, minimum size=1.3cm, align=center]{Apports inadéquats\\ (quantité/qualité)} }
    child { node[concept color=poppurple!70!white, minimum size=1.3cm, align=center]{Perte/Utilisation\\ altérée (maladies)} }
    child { node[concept color=poppurple!70!white, minimum size=1.3cm, align=center]{Facteurs socio-éco\\ (pauvreté, éducation)} }
  }
  child[concept color=popgreen] {
    node[concept, minimum size=1.8cm] {Conséquences}
    child { node[concept color=popgreen!70!white, minimum size=1.3cm, align=center]{Retard croissance\\ Système immunitaire} }
    child { node[concept color=popgreen!70!white, minimum size=1.3cm, align=center]{Mortalité infantile\\ (MPE, Vit A)} }
    child { node[concept color=popgreen!70!white, minimum size=1.3cm, align=center]{Pathologies chroniques\\ (HTA, Diabète, CVD)} }
  }
  child[concept color=popblue] {
    node[concept, minimum size=1.8cm] {Diagnostic}
    child { node[concept color=popblue!70!white, minimum size=1.3cm, align=center]{Courbes croissance\\ (OMS/UNICEF)} }
    child { node[concept color=popblue!70!white, minimum size=1.3cm, align=center]{IMC = $\frac{P}{T^2}$\\ (kg/m$^2$)} }
    child { node[concept color=popblue!70!white, minimum size=1.3cm, align=center]{Signes cliniques\\ Bilans biologiques} }
  }
  child[concept color=popgold] {
    node[concept, minimum size=1.8cm, align=center]{Prévention\\ \& Lutte}
    child { node[concept color=popgold!70!white, minimum size=1.3cm, align=center]{Alimentation équilibrée\\ \& diversifiée} }
    child { node[concept color=popgold!70!white, minimum size=1.3cm, align=center]{Supplémentation ciblée\\ (Vit A, Fer, Iode)} }
    child { node[concept color=popgold!70!white, minimum size=1.3cm, align=center]{Éducation nutritionnelle\\ Hygiène, Activité physique} }
  };
\end{tikzpicture}
\legende{Carte mentale récapitulative de la leçon : classification, étiologie, diagnostic et stratégies de prévention des maladies nutritionnelles.}
\end{popfigure}

\begin{definition}[Définitions clés]
\begin{itemize}
  \item \textbf{Malnutrition} : État pathologique résultant d'un déséquilibre (carence ou excès) entre les apports nutritionnels et les besoins de l'organisme.
  \item \textbf{Malnutrition Protéino-Énergétique (MPE)} : Carence concomitante en protéines et en énergie. Formes cliniques : \textbf{Kwashiorkor} (œdèmes, dermatose en « peinture craquelée », hépatomégalie, cheveux décolorés « en drapeau »), \textbf{Marasme} (amaigrissement extrême type « vieillard miniature », perte masse musculaire/graisseuse, absence d'œdème), \textbf{Marasmo-kwashiorkor} (association des signes des deux formes).
  \item \textbf{Avitaminose} : Carence en une ou plusieurs vitamines. Exemples majeurs : Vitamine A (xérophtalmie, cécité crépusculaire), Vitamine D (rachitisme chez l'enfant, ostéomalacie chez l'adulte), Vitamine B1 (béribéri sec/humide), Vitamine C (scorbut), Vitamine PP ou B3 (pellagre : dermite, diarrhée, démence).
  \item \textbf{Anémie} : Diminution de la concentration en hémoglobine sanguine en dessous des seuils de référence (Enfant 6--59 mois $<$ 110 g/L, Femme enceinte $<$ 110 g/L, Femme non enceinte $<$ 120 g/L, Homme $<$ 130 g/L). Principales causes nutritionnelles : carence en \textbf{fer} (anémie microcytaire hypochrome), en vitamine \textbf{B9} ou \textbf{B12} (anémie macrocytaire mégaloblastique).
  \item \textbf{Obésité} : Excès de masse grasse nuisible à la santé. Diagnostic chez l'adulte par \textbf{IMC} $\ge$ 30 kg/m$^2$. Surpoids : 25 $\le$ IMC < 30. Chez l'enfant, utilisation des courbes de référence OMS (Poids/Taille $>+2$ Z : surpoids ; $>+3$ Z : obésité).
\end{itemize}
\end{definition}

\begin{propriete}[Indicateurs anthropométriques et biologiques : formules et seuils de référence]
\begin{center}
\begin{tabularx}{\linewidth}{|l|X|X|}\hline
\rowcolor{popblueL}
\textbf{Indicateur} & \textbf{Formule / Définition} & \textbf{Seuils de référence (Adulte)} \\ \hline
IMC (Indice de Masse Corporelle) & $\displaystyle \text{IMC} = \frac{\text{Poids (kg)}}{\text{Taille (m)}^2}$ & Maigreur $<$ 18,5 \textbar\ Normal 18,5--24,9 \textbar\ Surpoids 25--29,9 \textbar\ Obésité $\ge$ 30 \\ \hline
Courbe de croissance (Enfant 0--5 ans) & Poids/Âge, Taille/Âge, Poids/Taille (Z-scores ou percentiles OMS) & $<-2$ Z : Retard/Insuffisance \textbar\ $<-3$ Z : Sévère \textbar\ $>+2$ Z : Surpoids/Obésité \\ \hline
Hémoglobine (Anémie) & Dosage sanguin (g/L) & Enfant 6--59 mois $<$ 110 \textbar\ Femme enceinte $<$ 110 \textbar\ Femme non enceinte $<$ 120 \textbar\ Homme $<$ 130 \\ \hline
\end{tabularx}
\end{center}
\end{propriete}

\begin{methode}[Calculer et interpréter l'IMC (adulte)]
\begin{itemize}
  \item Mesurer le poids (kg) et la taille (m) dans de bonnes conditions (à jeun, vêtements légers, balance calibrée, toise verticale).
  \item Appliquer la formule : $\text{IMC} = P / T^2$.
  \item Classer selon les seuils OMS (voir tableau ci-dessus).
  \item \textbf{Attention :} L'IMC ne distingue pas masse grasse et masse musculaire ; il est invalide ou peu fiable pour la femme enceinte, le sportif de haut niveau, la personne âgée (sarcopénie) et l'enfant (utiliser les courbes de croissance).
\end{itemize}
\end{methode}

\begin{methode}[Lire une courbe de croissance (enfant 0--5 ans, standards OMS/UNICEF)]
\begin{itemize}
  \item Repérer l'âge (abscisse) et la mesure (poids ou taille, ordonnée) sur le carnet de santé ou le graphique.
  \item Localiser le point par rapport aux courbes de référence (lignes Z = -3, -2, -1, 0, +1, +2, +3).
  \item Interpréter les trois indicateurs principaux :
    \begin{itemize}
      \item \textbf{Poids/Âge $<-2$ Z} : \textbf{Insuffisance pondérale} (sous-poids), reflet d'une malnutrition globale (aigüe ou chronique).
      \item \textbf{Taille/Âge $<-2$ Z} : \textbf{Retard de croissance} (nanisme), témoin d'une malnutrition \emph{chronique} (installée dans la durée).
      \item \textbf{Poids/Taille $<-2$ Z} : \textbf{Émaciation} (maigreur aigüe), témoin d'une malnutrition \emph{aigüe} (récente, sévère). $<-3$ Z : émaciation sévère (risque vital).
      \item \textbf{Poids/Taille $>+2$ Z} : \textbf{Surpoids} ; $>+3$ Z : \textbf{Obésité}.
    \end{itemize}
  \item Surveiller la \textbf{vitesse de croissance} (pente de la courbe) : une inflexion descendante (cassure de la courbe) est un signal d'alerte précoce, même si le point reste dans la norme.
\end{itemize}
\end{methode}

\begin{methode}[Identifier une maladie nutritionnelle à partir d'un cas clinique]
\begin{itemize}
  \item Lister les \textbf{signes cliniques} observés (œdèmes, dermatose, amaigrissement, pâleur conjonctivale/palmaire, cécité crépusculaire, goitre, douleurs osseuses, hémorragies gingivales, etc.).
  \item Noter les \textbf{données anthropométriques} (IMC, position sur courbes de croissance, Z-scores) et \textbf{biologiques} (Hb, rétinolémie, ferritine, taux de vitamine D) si disponibles.
  \item Rechercher le \textbf{contexte étiologique} : régime alimentaire (fréquence, diversité, mode de cuisson), pathologies associées (diarrhées chroniques, parasitoses intestinales, VIH/Sida, tuberculose), statut vaccinal, niveau socio-économique, pratiques d'allaitement et de diversification.
  \item Conclure par le diagnostic le plus probable (ex : « Marasme sévère compliqué d'anémie ferriprive ») et proposer la prise en charge adaptée (thérapeutique urgente + éducation nutritionnelle + suivi).
\end{itemize}
\end{methode}

\begin{aretenir}[Checklist avant le Probatoire]
\begin{itemize}
  \item $\square$ Savoir définir la malnutrition, la MPE (différencier Kwashiorkor vs Marasme vs Marasmo-kwashiorkor), les avitaminoses majeures (A, D, B1, C, PP), l'anémie nutritionnelle (types et causes) et l'obésité.
  \item $\square$ Connaître les signes cliniques distinctifs du Kwashiorkor (œdèmes bilatéraux, dermatose en « peinture craquelée », cheveux « drapeau », hépatomégalie) et du Marasme (amaigrissement extrême « peau sur les os », pas d'œdème, aspect « vieillard miniature »).
  \item $\square$ Citer les causes principales des carences : apports insuffisants (quantité/qualité), pertes excessives (diarrhées, parasitoses, hémorragies), besoins accrus (croissance, grossesse, allaitement, infection), mauvaise utilisation (malabsorption, maladies métaboliques).
  \item $\square$ Savoir calculer l'IMC ($P/T^2$) et classer le statut pondéral adulte selon les seuils OMS (18,5 ; 25 ; 30).
  \item $\square$ Savoir interpréter les 3 indicateurs anthropométriques de l'enfant (P/Â, T/Â, P/T) en Z-scores : insuffisance pondérale, retard de croissance (chronique), émaciation (aigüe), surpoids/obésité.
  \item $\square$ Expliquer les conséquences de la MPE sur le système immunitaire (atrophie thymique, baisse IgA, risque infectieux majoré : rougeole, diarrhée, pneumonie) et le développement cognitif (retard irréversible si survenu < 2 ans).
  \item $\square$ Lister les conséquences métaboliques de l'obésité : résistance à l'insuline (diabète type 2), HTA, dyslipidémie (hypertriglycéridémie, HDL bas), risques cardiovasculaires (AVC, IDM), atteintes ostéo-articulaires, apnée du sommeil.
  \item $\square$ Connaître les stratégies de prévention primaire : alimentation diversifiée (céréales + légumineuses + fruits/légumes + produits animaux), allaitement maternel exclusif 6 mois, diversification alimentaire à 6 mois (bouillies enrichies), hygiène alimentaire et eau potable.
  \item $\square$ Citer les programmes de supplémentation ciblée au Cameroun : Vitamine A (enfants 6--59 mois \emph{et} femmes en post-partum immédiat), Fer/Acide folique (femmes enceintes \emph{pendant} la grossesse), Iode (sel iodé obligatoire), Déparasitage (Albendazole/Mebendazole enfants 12--59 mois et femmes enceintes au 2ème trimestre).
  \item $\square$ Savoir proposer une ration alimentaire équilibrée (répartition glucides 50--55\% / lipides 30--35\% / protéines 10--15\% de l'apport énergétique total), à densité énergétique et micronutritionnelle adaptée à l'âge et l'état physiologique.
  \item $\square$ Éviter les confusions fréquentes : Kwashiorkor $\neq$ Marasme (présence/absence d'œdèmes), Anémie ferriprive (microcytaire) $\neq$ Anémie mégaloblastique B9/B12 (macrocytaire), Surpoids $\neq$ Obésité (seuil IMC 30), Retard de croissance/Taille-Âge (chronique) $\neq$ Émaciation/Poids-Taille (aigüe).
\end{itemize}
\end{aretenir}

\findecours

\end{document}
